% ---------------------------------------------------------------------------
% Chapter 2: The Mother Gives Birth: Parton Showering
% The scripts, configurations and records discussed here live in example/PS/
% of https://github.com/ravindkv/JitJet . The book gives the physics and the
% instructions; the code itself is read on GitHub, not printed here.
% ---------------------------------------------------------------------------
%% chapter-local shorthands (\pt already carries a subscript, so it cannot take another)
\providecommand{\ptevol}{\ensuremath{p_{\mathrm{T,evol}}}\xspace}
\providecommand{\ptrad}{\ensuremath{p_{\mathrm{T,rad}}}\xspace}
\chapter{The Mother Gives Birth: Parton Showering}
\chaptermark{The Mother Gives Birth}
\label{ch:mother_gives_birth}

\begin{journeybox}
A coloured mother cannot travel alone. On her first steps she radiates gluons,
and those gluons split again: children and grandchildren appear, ordered in
transverse momentum. Some of them wander too far from the road and are lost
to the family forever. The two protons she left behind radiate too, and the
whole family is pushed sideways by children it never knew it had.
\end{journeybox}

\physicsline{collinear and soft singularities, the DGLAP splitting functions,
the Sudakov form factor and the veto algorithm, \PYTHIA's \pt-ordered dipole
shower for final- and initial-state radiation, backward evolution with PDF
ratios, interleaving, the dead cone of a heavy quark, out-of-cone radiation,
the CP5 shower settings, the parton-level record of our event}

\section{Welcome back: what we will see at this second stop}
\label{sec:mother_gives_birth:what_we_compute}

At the first stop we left the mother, a $b$ quark of $\pt=113.22\GeV$, at the
moment of her birth, together with her twin, the $\bar b$, and nothing else.
That is not a physical state. A quark carries colour, and a coloured charge
that has just been accelerated radiates gluons, exactly as an electron that
has just been accelerated radiates photons, only much more, because the
coupling is larger and because the gluons themselves radiate. By the time she
has travelled a femtometre the mother is surrounded by children. The parton
shower is the piece of the generator that produces them, and this stop is
about how.

There is again one number to take away, and it is the first entry in the
energy ledger that differs from the previous one: after the shower the
partons within $\dR<0.4$ of the mother carry $100\GeV$, not $113$. Where the
missing $13\GeV$ went, and why the ``family'' is defined with a cone at all,
is the physics of the chapter. As at the first stop we will meet the number
four times, from the top down:
\begin{enumerate}[nosep]
  \item \textbf{Watch it.} A film replays the shower of our event branching
    by branching, measures who radiates first, keeps the momentum accounts
    and shows the same hard event showered a dozen times
    (\cref{sec:mother_gives_birth:animation}).
  \item \textbf{Run it.} A single Python file, needing only NumPy and SciPy,
    showers the hard-process record of \cref{ch:ancestral_home} in a third
    of a second and writes the parton record we follow through the rest of
    the book (\cref{sec:mother_gives_birth:standalone_code}).
  \item \textbf{Derive it.} From the collinear limit of QCD we write down
    the splitting functions, the Sudakov form factor and the one-line
    algorithm that turns a form factor into a random emission, follow
    \PYTHIA's construction of the shower step by step, and read the record
    the script wrote
    (\crefrange{sec:mother_gives_birth:why_alone}{sec:mother_gives_birth:runaways}).
  \item \textbf{Check it against the official packages.} A validation
    script showers the same hard event two thousand times with the real
    \PYTHIA and compares the averages with the standalone shower's, and the
    CMS generator configuration of \cref{ch:ancestral_home}, which already
    ran the shower, gives us its log and its event to read
    (\cref{sec:mother_gives_birth:validation}).
\end{enumerate}
\MG does not appear on this list, and it is worth saying why at once: \MG
computes matrix elements and hands its events to \PYTHIA for the shower, so
at this stage there is no separate \MG number to compare with. The question of what
happens when the matrix element and the shower both want to produce the
same gluon is the subject of \cref{ch:whose_child}.

\Cref{tab:mother_gives_birth:numbers} gives away the ending, and it needs a
sentence of warning that \cref{tab:ancestral_home:three_numbers} did not.
A cross section is a number; a shower is a random process. The same hard
event showered twice gives two different families, and the spread between
them is not an error but physics: our event, showered once with the
standalone script, leaves the mother at $78\GeV$, while two thousand showers
of the same event leave her at $94\pm29\GeV$ on average, and the \CMSSW run
left her at $92\GeV$. The three numbers are one distribution. What can be
compared between two showers is the average over many events, and there the
standalone script and \PYTHIA agree to a few percent on everything hard and
to ten percent on the softest emissions.

\begin{table}[htbp]\centering\small
\caption[The shower of our event obtained three ways]{The shower of our
event obtained three ways. ``Family'' means the final-state partons within
$\dR<0.4$ of the $b$ quark after the shower. The first two columns are
averages over 2000 showers of the same hard event (the spread of a single
shower is quoted after $\pm$ where it matters; the error of a mean is 45
times smaller); the third column is the one shower with seed~1 that wrote
\file{example/PS/event_PS.lhe}; the last is the single event of the \CMSSW
log, in which multi-parton interactions and hadronisation were also on. All
runs use the CP5 shower settings of \cref{sec:mother_gives_birth:cp5}.}
\label{tab:mother_gives_birth:numbers}
\footnotesize
\begin{tabular}{lrrrr}\toprule
 & \PYTHIA~8.315 & standalone & standalone & \CMSSW \\
 & 2000 showers & 2000 showers & seed 1 & 1 event \\ \midrule
ISR branchings & $5.05\pm2.9$ & $4.88\pm2.9$ & $5$ & -- \\
FSR branchings & $19.5\pm11$ & $20.9\pm12$ & $33$ & -- \\
final partons (of which gluons) & $26.5$ ($21.3$) & $27.7$ ($22.0$) & $40$ ($27$) & $57$ drawn \\
partons with $\pt>20$ / $5$ / $1\GeV$ & $2.9$ / $7.1$ / $18.8$ & $2.8$ / $7.3$ / $19.5$ & $8$ / $18$ / $31$ & -- \\
$x$ after ISR, $+z$ / $-z$ parton & $0.051$ / $0.450$ & $0.042$ / $0.446$ & $0.017$ / $0.470$ & -- \\ \midrule
$b$: \pt after the shower [GeV] & $94.8\pm29$ & $93.9\pm29$ & $78.1$ & $91.6$ \\
$b$ family: \pt / $n$ / $m$ [GeV] & $108.4$ / $2.5$ / $8.1$ & $108.4$ / $2.7$ / $8.2$ & $100.5$ / $3$ / $13.5$ & -- \\
$\bar b$: \pt after the shower [GeV] & $91.6$ & $90.1$ & $73.2$ & $118.6$ \\
$\bar b$ family: \pt / $n$ / $m$ [GeV] & $104.8$ / $2.4$ / $8.0$ & $103.4$ / $2.7$ / $8.0$ & $92.1$ / $3$ / $7.0$ & -- \\
\bottomrule
\end{tabular}
\end{table}

\paragraph{Why a cone of $0.4$, already now.} We are five stops away from
the inn where jets are clustered (\cref{ch:the_inn}), and yet the table
speaks of partons ``within $\dR<0.4$''. I have chosen to measure the family
with the radius of the jets CMS will eventually build, because the question
the ledger asks at every stage is the same: how much of the mother's original
momentum will end up inside the jet? A parton radiated at a wide angle is
still her child, but it will never be counted as part of her jet, and the
shower is the first place where such children are lost. The cone is a
bookkeeping device here, not a jet algorithm; the algorithm comes later and
will not ask who the mother was.

\paragraph{A second picture to carry: the project that keeps being
subcontracted.} The family of this book tells you who is related to whom;
it does not tell you how the shower decides what happens next. For that I
find a second picture useful, and the numbers of our event fit it exactly.
Think of the hard process as a large project awarded to two main
contractors, the $b$ and the $\bar b$, each with a job already assigned.
What the shower reveals, round by round, is how that work gets
subcontracted.
\begin{enumerate}[nosep]
  \item \textbf{Contractors subcontract.} A quark emitting a gluon
    ($q\to qg$) is a contractor who keeps the job and brings in a
    subcontractor; a gluon splitting into two ($g\to gg$) is a
    subcontractor dividing its work between two of its own kind; a gluon
    turning into a quark pair ($g\to q\bar q$) hands the job to two
    complementary specialists. Every new subcontractor may subcontract
    again. That repeated subdivision \emph{is} the shower.
  \item \textbf{Everybody proposes, the largest proposal goes first.} At
    each round every eligible contractor, and the two suppliers upstream
    (the incoming partons), proposes its next restructuring with a
    ``size'' drawn at random according to the rules of QCD. In the first
    round of our event the $\bar b$ proposed $2.2\GeV$, the $b$ $5.2\GeV$,
    the $+z$ supplier $53.9\GeV$ and the $-z$ supplier $12.7\GeV$; the
    $53.9\GeV$ proposal is carried out and everybody proposes afresh below
    that size. A contractor may also propose \emph{nothing} above the
    minimum size: that is the no-emission probability, the Sudakov form
    factor. A losing proposal never becomes a contract: of the 467
    proposals accepted in our event, 38 were built.
  \item \textbf{Restructuring one contract changes a partner's
    allocation.} Contractors work in a fixed network of partners, the
    colour connections. When one divides its job, a connected partner must
    adjust its own allocation so that the overall accounts balance. That
    adjustment is \emph{recoil}, and it is why the mother's \pt moves from
    $113$ to $94$, $92$ and $112\GeV$ before she has subcontracted anything
    herself.
  \item \textbf{Initial-state radiation is the upstream history.} For a
    final-state contractor you ask ``whom does it subcontract to?''; for an
    incoming parton you ask ``which larger organisation supplied this
    contractor, and what else did that organisation hand out?''. The $\bar u$
    that entered our hard process turns out to have been supplied by a
    gluon that also sent a $u$ quark into the final state.
\end{enumerate}
The picture captures repeated branching, competition, no branching,
connected partners, recoil and backward ancestry. Its one limitation should
be said at once: the ``size'' is the shower's evolution scale, not a budget
and not a particle's energy. Four-momentum is conserved exactly by the
bookkeeping, and the division into individual emitters is an approximation
to quantum radiation, good in the collinear and soft limits where the
chapter derives it.

\section{Watch the shower happen: the animated journey}
\label{sec:mother_gives_birth:animation}

Out of the window first, as before. The script
\file{example/PS/Standalone/animate_standalone_parton_showering.py}
imports the standalone shower of \cref{sec:mother_gives_birth:standalone_code}
and renders a four-scene film of about forty seconds, so that every number
on the screen is computed by the same classes that wrote the record. It is
on the JitJet playlist,
\begin{center}
  \url{https://www.youtube.com/playlist?list=PLFtvDH9g5Km0}
\end{center}
and \cref{fig:mother_gives_birth:stills} shows one key frame per scene.
Each scene names the section that explains it; watch first, and let the
formulae catch up.
\begin{enumerate}
  \item \textbf{The shower unwinds} (\cref{sec:mother_gives_birth:step_by_step}).
    A scale runs down from $\pt^{\max}=113.3\GeV$ to the cut-off, and each
    time it passes the \pt of one of the 38 branchings of our event the
    branching happens: a ladder on the left fills rung by rung, orange
    squares for ISR and red or blue dots for the two families, while the
    $(\eta,\phi)$ picture on the right fills with partons, connected by the
    colour dipoles that do the radiating. The scene replays the seed-1
    shower branching by branching and checks itself against the record to
    half a MeV. Watch the first three rungs: all initial state, all above
    $40\GeV$, and every marker on the right jumps at each of them. Watch the
    ladder cross $\mb$ and $m_c$, where the number of flavours in the
    coupling changes, and watch a gluon near the $\bar b$ turn into a
    $\bbbar$ pair.
  \item \textbf{Who radiates first?} (\cref{sec:mother_gives_birth:sudakov,sec:mother_gives_birth:veto}).
    The four ends that can radiate from $\pt^{\max}$, the $b$, the $\bar b$
    and the two incoming partons, each propose their next emission 1500
    times. The survival curves on the left are the no-emission probabilities
    $\Delta(\pt^{\max},\pt)$ of each end, measured rather than computed:
    the $b$ and $\bar b$ curves lie on top of each other, the $+z$ beam at
    small $x$ falls fastest, and the dashed product is the probability that
    nobody has radiated yet. On the right, the histogram of the hardest of
    the four proposals, coloured by who made it: the $+z$ beam wins about
    $53\%$ of the time, the $-z$ beam $18\%$, the $b$ and the $\bar b$ some
    $15\%$ each,
    and in one trial in a thousand nobody radiates at all. The seed-1
    event's first branching at $53.9\GeV$ is marked.
  \item \textbf{Momentum bookkeeping} (\cref{sec:mother_gives_birth:runaways}).
    The shower runs again while the transverse plane on the left keeps the
    accounts: the $b$ and $\bar b$ arrows, their sum, the ISR partons and
    their sum, and the total, which stays at zero. On the right the two
    incoming partons walk up in $x$ on top of the PDFs at the current
    scale, from $2.2\times10^{-3}$ to $1.7\times10^{-2}$ and from $0.39$ to
    $0.47$, while the curves themselves change with the falling scale.
  \item \textbf{No two alike} (\cref{tab:mother_gives_birth:numbers}). The
    same hard event with seeds $1,2,3,\ldots$: twelve outcomes side by side,
    and the running means of the branching counts and of the family \pt of
    each $b$ quark settling, after 64 showers, onto $5.2$ and $22$
    branchings and $113$ and $106\GeV$, next to the \PYTHIA values of the
    README. Our event, the first frame, is one of the less lucky ones.
\end{enumerate}
To render it you need matplotlib in addition to NumPy and SciPy, and
\code{ffmpeg} for an \code{.mp4}:
\begin{minted}[linenos=false,frame=single,framesep=1.5mm,fontsize=\small]{bash}
cd example/PS/Standalone
python3 animate_standalone_parton_showering.py           # the film, about a minute
python3 animate_standalone_parton_showering.py --fast    # low-resolution preview
python3 animate_standalone_parton_showering.py --scene sudakov -o sudakov.gif
python3 animate_standalone_parton_showering.py --stills stills.pdf --no-video
\end{minted}
The physics options are those of the shower script we run next, so \code{--seed 7}
films another family, \code{--no-dampen} shows what
\cref{eq:mother_gives_birth:dampen} does to the survival curves, and
\code{--ptmax 50} starts the shower lower.

\begin{figure}[htbp]\centering
\includegraphics[page=1,width=0.49\textwidth]{PS/Standalone/animate_standalone_parton_showering_stills.pdf}\hfill
\includegraphics[page=2,width=0.49\textwidth]{PS/Standalone/animate_standalone_parton_showering_stills.pdf}\\[2mm]
\includegraphics[page=3,width=0.49\textwidth]{PS/Standalone/animate_standalone_parton_showering_stills.pdf}\hfill
\includegraphics[page=4,width=0.49\textwidth]{PS/Standalone/animate_standalone_parton_showering_stills.pdf}
\caption[Key frames of the animated standalone shower]{Key frames of the
four scenes of the animated standalone shower, drawn by
\file{animate_standalone_parton_showering.py} with \code{--stills}. Top left:
the ladder of branchings in decreasing \pt and the $(\eta,\phi)$ picture
filling up. Top right: the measured survival curves $\Delta(\pt^{\max},\pt)$
of the four radiators and the histogram of who radiates first. Bottom left:
the transverse-momentum balance and the walk of the incoming partons in $x$.
Bottom right: twelve showers of the same hard event and the running means
converging to the \PYTHIA reference. The film is on the JitJet YouTube
playlist.}
\label{fig:mother_gives_birth:stills}
\end{figure}

\section{Press the button: the standalone shower}
\label{sec:mother_gives_birth:standalone_code}

Now make a family yourself. The material of this chapter is in
\file{example/PS/} of the JitJet repository,
\begin{center}
  \url{https://github.com/ravindkv/JitJet}
\end{center}
with the shower and its validation under \file{example/PS/Standalone/}, the
event record and the plotting script one level up, and the \CMSSW log and
figures under \file{example/PS/CMSSW/}. As before the code is not printed;
each script gets a description, the command and the output to expect.

\subsection{The standalone shower script}

\file{example/PS/Standalone/standalone_parton_showering.py} contains the
construction that \cref{sec:mother_gives_birth:dipole_shower} will describe,
and nothing else: about a thousand lines of Python needing NumPy and SciPy, no LHAPDF,
no \PYTHIA. The PDF grid (member~0 of NNPDF3.1 on LHAPDF's own knots, all
flavours, $x\ge10^{-3}$, $Q$ up to $140\GeV$) is embedded between the markers
\code{BEGIN/END GENERATED TABLES}, produced by \file{dump_shower_tables.py}
in an environment that has LHAPDF. A map of the file: \code{AlphaStrong}
(second-order running with thresholds, \cref{eq:mother_gives_birth:lambdas});
\code{TablePDF} (bicubic spline in $(\ln x,\ln Q^2)$, frozen below the
grid); \code{Particle}, \code{Event}, \code{read\_event} and \code{listing}
(\PYTHIA's record format, in and out); \code{Shower} with
\code{pT2next\_fsr} and \code{pT2next\_isr} (the veto algorithm,
\cref{sec:mother_gives_birth:veto}), \code{fsr\_kinematics} and
\code{isr\_kinematics} (the recoils), \code{branch\_fsr}, \code{branch\_isr}
and \code{run} (steps~4 and~5); \code{verify} (energy--momentum, masses and
colour lines after every branching, for \code{--check}); and
\code{family}, \code{summarise} and \code{main}. From the repository root:
\begin{minted}[linenos=false,frame=single,framesep=1.5mm,fontsize=\small]{bash}
cd example/PS/Standalone
python3 standalone_parton_showering.py                       # seed 1 -> ../event_PS.lhe
python3 standalone_parton_showering.py --seed 7 --history h.json
python3 standalone_parton_showering.py --no-isr              # the two b quarks only
python3 standalone_parton_showering.py --check               # assert conservation laws
python3 standalone_parton_showering.py --repeat 2000 --quiet --history s.json
python3 standalone_parton_showering.py -v > standalone_parton_showering.log   # learning mode
python3 standalone_parton_showering.py -vv | less                             # ... every trial too
\end{minted}
The first command takes about a second, most of it parsing the tables, and
prints
\begin{minted}[linenos=false,frame=single,framesep=1.5mm,fontsize=\footnotesize]{text}
input ../../ME/event_ME.lhe: hard process g u -> bbar b, pTmax = 113.32 GeV, seed = 1
shower finished: 5 ISR + 33 FSR branchings, 40 final-state partons (27 gluons); ...
incoming after ISR: g x = 1.7140e-02 (+z), u x = 4.6999e-01 (-z)
parton  pT before  pT after     eta     phi  family n  family pT  family E  family m
bbar       113.22     73.18  -3.409   0.577         3      92.13   1414.97      7.02
b          113.22     78.12  -1.379  -2.728         3     100.45    224.59     13.52
branchings in pT order:
    1 ISR q<-g   side A(+z)  g  -> ubar + u   pT =  53.877  z = 0.779  x -> 2.7914e-03
    2 ISR q<-q   side B(-z)  u  -> u    + g   pT =  51.361  z = 0.850  x -> 4.5920e-01
    ...
\end{minted}
followed by the 38 branchings and the record, written to
\file{example/PS/event_PS.lhe}. With \code{--history} the branchings and the
summary go to a JSON file as well; with \code{--repeat N} the script showers
the event $N$ times with seeds $1\ldots N$ and writes the per-event
quantities of \cref{tab:mother_gives_birth:numbers} to
\file{s.stats.json}, about ten minutes for $2000$. The physics options are
\code{--alphas} and \code{--order} (the coupling), \code{--ptmin-fsr},
\code{--ptmin-isr}, \code{--pt0} and \code{--ptmax} (the scales of
\cref{sec:mother_gives_birth:scales}), \code{--no-isr}, \code{--no-fsr},
\code{--no-dampen} (\cref{eq:mother_gives_birth:dampen} off) and
\code{--radius} (the family cone). The output with \code{--seed 1} is
byte-for-byte the committed record; any other seed gives another family,
which is the point of Exercise~1.

The two last commands are the ones to run with this chapter open.
\code{-v} is a learning mode: before the shower it prints the running
coupling at the scales that matter, the input record with $x_A$, $x_B$,
$\shat$ and the starting scale, the parton densities and PDF ratios the
backward evolution needs, and the colour dipoles with their masses; then,
for every step, the competition (each dipole end and each beam with its
proposed scale, $z$ and channel, or ``no emission''), the winner, and the
winning branching worked out in full with every momentum before and after
and the conservation checks. \code{-vv} adds every trial of the veto
algorithm: the overestimate integrals, the random number and the scale it
produces, the chosen $z$ and each factor of the acceptance weight. The
random numbers are not touched, so the record is identical with and
without the flags. The $1700$ lines of the \code{-v} output for seed~1 are
committed as \file{example/PS/Standalone/standalone_parton_showering.log},
and every number quoted in the theory sections that follow is taken from
that file, so that you can put the formula and the log side by side.
The plots of the theory sections come from a sibling script,
\begin{minted}[linenos=false,frame=single,framesep=1.5mm,fontsize=\small]{bash}
python3 example/PS/Standalone/plot_theory_PS.py      # -> theory_parton_showering.pdf, one page per section
\end{minted}
which replays the seed-1 shower with the classes of the shower script,
reads the committed record and parses the \code{-v} log, so that every
curve and every marker in those figures can be traced to a line of the
log; \code{make figures} regenerates the PDF.

\subsection{Reading the record}
\label{sec:mother_gives_birth:record_format}

The record keeps \PYTHIA's conventions so that the same reader serves the
standalone and the \CMSSW output. Status codes: $-21$ and $23$ are the hard
process as before; $43$ a parton emitted by ISR, $-41$ the new incoming
parton, $44$ a final-state parton shifted by an ISR recoil; $51$ the two
daughters of an FSR branching, $52$ a final-state recoiler, $-53$ an
incoming recoiler. A negative status means the parton has been replaced by
its daughters and is no longer in the event; the final state is every row
with a positive status, and its momentum sum equals that of the two current
incoming partons. Every row names its mother, so the family of any final
parton can be found by walking up until a hard-process row ($-23$) or an
ISR emission ($43$) is reached; that is what the plotting script does.

\subsection{Drawing the event}

\file{example/PS/plot_journey_PS.py} reads the record and draws every
final-state parton in the $(\eta,\phi)$ plane in the layout of the \CMSSW
analyser of \cref{sec:mother_gives_birth:cmssw}
(\cref{fig:mother_gives_birth:journey}), classifying each parton by walking
its mother chain back to the hard $b$, the hard $\bar b$ or an ISR emission,
and prints a table of all partons with their origin, the momentum sums of
\cref{sec:mother_gives_birth:runaways} and the $x$ walk of the incoming
partons. It needs matplotlib only:
\begin{minted}[linenos=false,frame=single,framesep=1.5mm,fontsize=\small]{bash}
python3 example/PS/plot_journey_PS.py                  # -> example/PS/journey_PS.pdf
python3 example/PS/plot_journey_PS.py -i my.lhe --label-min 10 --png
\end{minted}

\section{Why a quark cannot travel alone: collinear and soft singularities}
\label{sec:mother_gives_birth:why_alone}

You have watched the family appear, and you have produced one yourself. Now
we earn it. The next five sections derive the branching probability, turn
it into the algorithm the script executes, follow \PYTHIA's construction of
the shower step by step, and end by reading the record you have just
written.

Before we derive anything, a puzzle to carry through the next pages. Every
gluon that is emitted costs one power of \alphas, which over the scales of
our shower lies between $0.11$ and $0.2$. Three emissions should then be
suppressed by $\alphas^3\approx10^{-3}$ relative to the Born process, and
the $38$ branchings of our event by $\alphas^{38}\approx10^{-33}$. Yet
every event looks like ours: \PYTHIA averaged $5$ initial- and $20$
final-state branchings over the $2000$ showers of
\cref{sec:mother_gives_birth:pythia}. Why does a chain that is suppressed
at every link go on for dozens of links? Because \alphas never comes alone.
Each power of the coupling is multiplied by an integral over the angle of
the emission and by an integral over its energy, and both integrals are
logarithms that in our event reach five. A small number times a large
logarithm squared is of order one, and the expansion parameter of a shower
is not \alphas but $\alphas\ln^2$. The propagator below shows where the two
logarithms come from.

\begin{figure}[htbp]\centering
\resizebox{\textwidth}{!}{% ---------------------------------------------------------------------------
% The emission of a gluon by the outgoing b: the propagator that makes a
% quark radiate, its two singular limits, and the double-logarithmic phase
% space in the (ln 1/theta, ln E_g) plane. Numbers from the -v log of
% standalone_parton_showering.py (section 2 and branching 5) and from
% chapter 2's text (beta_b = 0.99982, 7.5 octaves, ln = 5.2).
% ---------------------------------------------------------------------------
\begin{tikzpicture}[x=1cm,y=1cm,
  every node/.style={font=\footnotesize},
  fermion/.style={thick,postaction={decorate},decoration={markings,mark=at position 0.55 with {\arrow[scale=1.1]{Stealth}}}},
  gluon/.style={thick,decorate,decoration={coil,aspect=0,segment length=4.5pt,amplitude=2pt}},
  vertex/.style={fill=black,circle,inner sep=1.6pt},
  lab/.style={font=\scriptsize,align=center},
  fac/.style={font=\scriptsize,align=center,text=orange!80!black},
  head/.style={font=\small\bfseries},
  axis/.style={-{Stealth},thick}]
% ---------------- (a) the propagator -----------------------------------------------------------
\node[head,anchor=west] at (-0.3,3.4) {(a) the real correction $q\bar q\to\bbbar g$: the propagator};
\node[circle,draw,thick,fill=orange!25,minimum size=11mm] (H) at (0.6,0) {$\hat\sigma$};
\draw[fermion] (H.east) -- (3.2,0) node[midway,below,lab] {$b^*$, off shell:\\$(p_b+k)^2-\mb^2=2p_b\!\cdot\!k$};
\draw[fermion] (3.2,0) -- (6.0,0.9) node[right,lab,text=red!80!black] {$b$, $E_b=252\GeV$, $\beta_b=0.99982$};
\draw[gluon] (3.2,0) -- (6.0,-1.1) node[right,lab] {$g$, energy $E_g$, angle $\theta$ to the $b$};
\node[vertex] at (3.2,0) {};
\draw[-{Stealth},thin] (4.4,0.39) arc[start angle=18,end angle=-21,radius=1.25];
\node[lab] at (4.75,-0.05) {$\theta$};
\draw[fermion] (H.north east) -- (2.2,1.8) node[above right,lab,text=blue!80!black] {$\bar b$};
\node[fac,text width=6.8cm,anchor=north west] at (-0.3,-1.6) {$\dfrac{1}{2p_b\!\cdot\!k}=\dfrac{1}{2E_bE_g\,(1-\beta_b\cos\theta)}$: large when $\theta\to0$ (collinear) and when $E_g\to0$ (soft). For a massless quark $1-\cos\theta\to0$; for our mother $1-\beta_b\cos\theta\ge1.8\times10^{-4}$, the dead cone.};
% ---------------- (b) the double-logarithmic phase space -----------------------------------------
\begin{scope}[xshift=11.4cm]
\node[head,anchor=west] at (-0.3,3.4) {(b) the phase space of one emission: an area};
\fill[blue!8] (0,0) -- (5.6,0) -- (5.6,2.4) -- cycle;
\draw[axis] (-0.3,0) -- (6.4,0) node[below left,lab] {$\ln(1/\theta)$: more collinear $\to$};
\draw[axis] (0,-0.3) -- (0,2.7) node[right,lab,anchor=south west] at (0,2.5) {$\ln(E_b/E_g)$: softer $\uparrow$};
\draw[thick] (0,0) -- (5.6,2.4) node[pos=0.6,above,sloped,lab] {$\pt=E_g\theta=\pt^{\min}$: the cut-off};
\draw[thick] (5.6,-0.1) -- (5.6,2.5) node[right,lab] {$\theta=\theta_0=\mb/E_b$\\the dead cone};
\node[lab,anchor=west] at (0.15,0.25) {$\pt\approx\pt^{\max}$};
\draw[thick,gray!70] (0,0.0) -- (0,0.0);
\node[star,star points=5,star point ratio=2.2,fill=red!80!black,draw=white,inner sep=0pt,minimum size=3mm] at (2.7,0.55) {};
\node[lab,text=red!80!black,anchor=west] at (2.85,0.55) {br.\ 5: $\ptevol=32\GeV$, $z=0.60$};
\node[star,star points=5,star point ratio=2.2,fill=red!80!black,draw=white,inner sep=0pt,minimum size=3mm] at (4.4,1.45) {};
\node[lab,text=red!80!black,anchor=west] at (4.55,1.45) {br.\ 13: $4.4\GeV$, $z=0.86$};
\node[lab,text=black!65,text width=6.6cm,anchor=north west] at (-0.3,-0.55) {Every point of the triangle costs $\alphas/\pi$ and the measure $\mathrm{d}\ln\theta^2\,\mathrm{d}\ln(1/(1-z))$ is flat in it, so the emission probability is $\alphas/\pi$ times an \emph{area}: the two logarithms of \cref{eq:mother_gives_birth:double_log}. For our shower both sides reach $\ln(113/0.611)=5.2$, and $(\alphas/\pi)\times5.2^2\approx1$.};
\end{scope}
\end{tikzpicture}
}
\caption[The propagator that makes a quark radiate, and the double-logarithmic phase space]{(a) The real correction to the hard process: the outgoing $b$ is off shell before it emits the gluon, and the propagator $1/(2p_b\cdot k)$ of \cref{eq:mother_gives_birth:propagator} is large for a collinear and for a soft gluon; for our massive mother the collinear growth stops at $1-\beta_b$, the dead cone of \cref{sec:mother_gives_birth:dead_cone}. (b) The phase space of one emission in the plane of $\ln(1/\theta)$ and $\ln(E_b/E_g)$: the emission probability is $\alphas/\pi$ times the area of the triangle bounded by the cut-off and the dead cone, the double logarithm of \cref{eq:mother_gives_birth:double_log}. The mother's two hardest emissions of the \code{-v} log are marked.}
\label{fig:mother_gives_birth:propagator_sketch}
\end{figure}

Look out of the window at the mother leaving the hard process. In the
partonic frame she carries $E=\sqrt{\shat}/2=198\GeV$ and is, at Born level,
an on-shell quark of mass $4.8\GeV$. Now ask what the next order of
perturbation theory does to her. The real correction to $q\bar q\to\bbbar$
is $q\bar q\to\bbbar g$, and when the gluon is emitted by the outgoing $b$
the amplitude contains the propagator of the $b$ just before the emission,
\begin{equation}
  \frac{1}{(p_b+k)^2-\mb^2}=\frac{1}{2\,p_b\cdot k}
  =\frac{1}{2E_bE_g\,(1-\beta_b\cos\theta)},
  \label{eq:mother_gives_birth:propagator}
\end{equation}
where $\theta$ is the angle between the $b$ and the gluon and
$\beta_b$ the $b$ velocity. For a massless quark $\beta_b=1$, and the
denominator vanishes twice: when $\theta\to0$ (a \emph{collinear} emission)
and when $E_g\to0$ (a \emph{soft} emission). Our mother is not quite
massless: in the laboratory she leaves the hard process with
$E_b=252.06\GeV$ and $m_b=4.8\GeV$, so $\beta_b=0.99982$ and the factor
$1-\beta_b\cos\theta$ never falls below $1.8\times10^{-4}$. Keep that
number; it is the dead cone of \cref{sec:mother_gives_birth:dead_cone}, and
for everything but the softest emissions it makes no difference. Squaring the amplitude and
adding phase space, the emission probability in these limits factorises
from the Born cross section,
\begin{equation}
  \mathrm{d}\sigma_{\bbbar g}\;\simeq\;\mathrm{d}\sigma_{\bbbar}\;
  \frac{\alphas}{2\pi}\,\frac{\mathrm{d}\theta^2}{\theta^2}\;
  P_{qq}(z)\,\mathrm{d}z,
  \qquad
  P_{qq}(z)=C_F\,\frac{1+z^2}{1-z},\quad C_F=\tfrac43 ,
  \label{eq:mother_gives_birth:collinear_limit}
\end{equation}
with $z$ the fraction of the parent's energy kept by the quark and $1-z$
handed to the gluon. Put two numbers into the kernel before reading on:
$P_{qq}(0.5)=\tfrac43\cdot\tfrac{1.25}{0.5}=3.33$ and
$P_{qq}(0.9)=\tfrac43\cdot\tfrac{1.81}{0.1}=24.1$. When the quark keeps
$90\%$ of the energy and the gluon takes a soft $10\%$, the kernel is seven
times larger than for an equal split; that is the preference for soft
gluons, and the kernel is not a probability but a rate, which is why it may
exceed one. Read the formula as we read the factorisation formula of the
previous stop: the Born cross section does not know about the emission, the
emission does not know about the Born process, and what connects them is a
universal function of $z$ and a logarithmic measure in the angle. The
$\mathrm{d}\theta^2/\theta^2$ is the collinear singularity and the
$1/(1-z)$ the soft one, and where both apply the probability is
\emph{double} logarithmic,
\begin{equation}
  \int\frac{\mathrm{d}\theta^2}{\theta^2}\int\frac{\mathrm{d}z}{1-z}
  \;\sim\;\ln^2 ,
  \label{eq:mother_gives_birth:double_log}
\end{equation}
which multiplied by $\alphas/\pi\approx0.04$ is not small once the
logarithms reach five or ten. They do: the shower of our event fills the
range from $113.3\GeV$ down to the cut-off at $0.611\GeV$, seven and a half
octaves, and $\ln(113.3/0.611)=5.2$, so that
$(\alphas/\pi)\ln^2\approx0.04\times27\approx1$. That resolves the puzzle: each
extra emission costs one power of \alphas and buys a logarithm squared of
the same size, so the chain is not suppressed at all, and it is the whole
reason for the shower.
Fixed-order perturbation theory in $\alphas$ is an expansion in a small
number multiplied by large logarithms, and the product is of order one; the
shower resums the logarithms to all orders by iterating
\cref{eq:mother_gives_birth:collinear_limit} emission after emission, at the
price of keeping only the collinear and soft parts of each matrix element.

\Cref{fig:mother_gives_birth:propagator_plot} shows the two statements in
numbers. The propagator factor of a massless quark grows as $2/\theta^2$
without limit; for our mother it saturates at $1/(1-\beta_b)=5515$ below
$\theta_0=0.019$. And the $38$ branchings of our event, counted per octave
of the evolution variable, rise from one or two per octave at the top to
sixteen in the last octave above the cut-off, where the coupling has grown
from $0.11$ to $0.5$: the first ten branchings span three octaves, the
last eighteen less than one.

\begin{figure}[htbp]\centering
\includegraphics[page=1,width=\textwidth]{PS/Standalone/theory_parton_showering.pdf}
\caption[The collinear singularity and the branchings per octave of our event]{Why a quark cannot travel alone, from \file{plot_theory_PS.py}. (a) The factor $1/(1-\beta\cos\theta)$ of \cref{eq:mother_gives_birth:propagator} against the emission angle, for a massless quark and for our mother with $E_b=252\GeV$; the mass turns the singularity into a plateau below $\theta_0=\mb/E_b$. (b) The $38$ branchings of the seed-1 shower per octave of $\ptevol$ below $\pt^{\max}$, coloured by kind, with the running coupling on the right axis: the rate per octave grows as the scale falls, the $\mathrm{d}\theta^2/\theta^2$ measure of \cref{eq:mother_gives_birth:collinear_limit} helped by the growing coupling.}
\label{fig:mother_gives_birth:propagator_plot}
\end{figure}

Two remarks before we move on. First, the singularities are not artefacts:
they tell us that a bare quark is not an observable, and that any
calculation that stops at a fixed order needs to be combined with a
definition of what is measured, a jet, that is insensitive to arbitrarily
soft and collinear splittings. That is a theme for \cref{ch:the_inn}.
Second, the same structure appears for the \emph{incoming} partons. The
$\bar u$ and the $u$ that entered the hard process were resolved out of the
protons at the scale $\muF$, and the DGLAP evolution that produced the PDFs
of \cref{sec:ancestral_home:pdfs} is the very same
\cref{eq:mother_gives_birth:collinear_limit}, integrated over the angle. The
shower run backwards from the hard process undoes that evolution one step
at a time, and a parton emitted along the way is a real gluon or quark in
the final state. That is initial-state radiation, and we will see that in
our event it is louder than the mother's own.

\section{The splitting functions and the Sudakov form factor}
\label{sec:mother_gives_birth:splitting_sudakov}

\subsection{The four splitting functions}
\label{sec:mother_gives_birth:splitting_functions}

Set the stage before the formulas. After the three initial-state branchings
of our event there are five final-state partons and seven dipole ends;
after thirty there are forty partons and sixty-seven ends, and at every
round each of them may branch in one of several ways. A quark can only emit
a gluon. A gluon may emit a gluon or split into a quark--antiquark pair of
any flavour light enough at that scale. An incoming parton may be found to
come from a quark or from a gluon. So which kind of branching do you expect
to win most rounds, and which one almost never? Make your guess from three
things: the colour charge of the emitter, the number of emitters of each
kind at that moment, and whether the kernel favours a soft daughter. Then
compare with the list of our event: $38$ branchings, of which $5$ are
initial-state, $9$ are $q\to qg$, $19$ are $g\to gg$ and only $5$ are
$g\to q\bar q$. \Cref{fig:mother_gives_birth:ladder} below the table shows
them in the order in which they happened; the functions that set these
proportions are the following.

\begin{figure}[htbp]\centering
\resizebox{\textwidth}{!}{% ---------------------------------------------------------------------------
% The four splitting functions as vertices, with their kernels, colour
% factors and the number of times each occurred among the 38 branchings of
% our event (-v log of standalone_parton_showering.py: 9 q->qg, 19 g->gg,
% 5 g->qqbar, 5 ISR of which 2 q<-q, 1 q<-g, 2 g<-g).
% ---------------------------------------------------------------------------
\begin{tikzpicture}[x=1cm,y=1cm,
  every node/.style={font=\footnotesize},
  fermion/.style={thick,postaction={decorate},decoration={markings,mark=at position 0.6 with {\arrow[scale=1.0]{Stealth}}}},
  antifermion/.style={thick,postaction={decorate},decoration={markings,mark=at position 0.4 with {\arrowreversed[scale=1.0]{Stealth}}}},
  gluon/.style={thick,decorate,decoration={coil,aspect=0,segment length=4pt,amplitude=1.8pt}},
  vertex/.style={fill=black,circle,inner sep=1.5pt},
  lab/.style={font=\scriptsize,align=center},
  kern/.style={font=\scriptsize,align=left},
  head/.style={font=\small\bfseries},
  count/.style={draw,rounded corners=2pt,fill=orange!20,font=\scriptsize,inner sep=2pt}]
\node[head,anchor=west] at (-0.4,2.6) {the four branchings of the shower, their kernels, and how often each happened in the 38 branchings of our event};
% (1) q -> q g
\begin{scope}
\draw[fermion] (0,0) -- (1.4,0); \draw[fermion] (1.4,0) -- (2.6,0.8) node[right,lab] {$q$, keeps $z$};
\draw[gluon] (1.4,0) -- (2.6,-0.8) node[right,lab] {$g$, takes $1-z$};
\node[vertex] at (1.4,0) {};
\node[lab,above] at (0.6,0.05) {$q$};
\node[kern,anchor=north west] at (-0.2,-1.3) {$P_{qq}(z)=C_F\dfrac{1+z^2}{1-z}$\\[3pt] soft pole at $z\to1$;\\$C_F=\tfrac43$};
\node[count] at (1.3,1.6) {$9\times$ ($b$, $\bar b$, $s$, $\bar s$ emitting)};
\end{scope}
% (2) g -> g g
\begin{scope}[xshift=5.6cm]
\draw[gluon] (0,0) -- (1.4,0); \draw[gluon] (1.4,0) -- (2.6,0.8) node[right,lab] {$g$, $z$};
\draw[gluon] (1.4,0) -- (2.6,-0.8) node[right,lab] {$g$, $1-z$};
\node[vertex] at (1.4,0) {};
\node[lab,above] at (0.6,0.1) {$g$};
\node[kern,anchor=north west] at (-0.2,-1.3) {$P_{gg}(z)=2C_A\dfrac{(1-z(1-z))^2}{z(1-z)}$\\[3pt] poles at both ends;\\$C_A=3$, $C_A/C_F=\tfrac94$};
\node[count] at (1.3,1.6) {$19\times$, the most frequent};
\end{scope}
% (3) g -> q qbar
\begin{scope}[xshift=11.2cm]
\draw[gluon] (0,0) -- (1.4,0); \draw[fermion] (1.4,0) -- (2.6,0.8) node[right,lab] {$q$, $z$};
\draw[antifermion] (1.4,0) -- (2.6,-0.8) node[right,lab] {$\bar q$, $1-z$};
\node[vertex] at (1.4,0) {};
\node[lab,above] at (0.6,0.1) {$g$};
\node[kern,anchor=north west] at (-0.2,-1.3) {$P_{qg}(z)=T_R\,[z^2+(1-z)^2]$\\[3pt] no pole: a soft quark\\is not enhanced; $T_R=\tfrac12$};
\node[count] at (1.3,1.6) {$5\times$ ($s\bar s$, $\bbbar$, \dots)};
\end{scope}
% (4) initial state: q <- q, q <- g, g <- g, g <- q
\begin{scope}[xshift=16.8cm]
\draw[fermion] (0,0.8) -- (1.4,0) node[pos=0.1,above,lab] {$b$ from the proton, $x/z$};
\draw[fermion] (1.4,0) -- (2.8,0) node[right,lab] {$a$ into $\hat\sigma$, $x$};
\draw[gluon] (1.4,0) -- (2.4,-1.0) node[right,lab] {$c$ emitted};
\node[vertex] at (1.4,0) {};
\node[kern,anchor=north west] at (-0.2,-1.3) {backward: $P_{ba}(z)\,\dfrac{x'f_b(x',\pt^2)}{xf_a(x,\pt^2)}$\\[3pt] the same kernels read\\upstream, times a PDF ratio};
\node[count] at (1.3,1.6) {$5\times$ ($2\;q\leftarrow q$, $1\;q\leftarrow g$, $2\;g\leftarrow g$)};
\end{scope}
\node[lab,text=black!65,text width=19cm,anchor=north west] at (-0.4,-3.1) {The daughter that continues keeps the fraction $z$ of the parent's momentum; the kernels are rates per unit $\ln t$ and may exceed one. $P_{gq}(z)=P_{qq}(1-z)$ is the first vertex read with the gluon as the continuing parton, needed for $g\leftarrow q$ in the initial state. The gluon's double pole and its colour factor are why gluons take over the radiating as the shower proceeds: $27$ of the $40$ final partons of our event are gluons.};
\end{tikzpicture}
}
\caption[The four branchings of the shower as vertices, with their kernels and their counts in our event]{The four branchings of the shower as vertices: the parent on the left, the daughter that keeps the fraction $z$ above, the emitted parton below, with the kernel of \cref{eq:mother_gives_birth:splitting}, its colour factor and the number of times it occurred among the $38$ branchings of the seed-1 shower (\code{-v} log). The fourth vertex is the initial state: the same kernels read upstream, weighted by the ratio of parton densities of \cref{eq:mother_gives_birth:isr}.}
\label{fig:mother_gives_birth:splittings_sketch}
\end{figure}

Every branching of the shower is one of four, and the probability of each is
a leading-order Altarelli--Parisi splitting function of the momentum fraction
$z$~\cite{ParticleDataGroup:2024cfk}:
\begin{align}
  P_{qq}(z)&=C_F\,\frac{1+z^2}{1-z}, & &q\to qg, \notag\\
  P_{gg}(z)&=2C_A\left[\frac{z}{1-z}+\frac{1-z}{z}+z(1-z)\right]
            =2C_A\,\frac{\bigl(1-z(1-z)\bigr)^2}{z(1-z)}, & &g\to gg, \notag\\
  P_{qg}(z)&=T_R\left[z^2+(1-z)^2\right], & &g\to q\bar q, \label{eq:mother_gives_birth:splitting}\\
  P_{gq}(z)&=C_F\,\frac{1+(1-z)^2}{z}, & &q\to gq\ (\text{the gluon keeps }z), \notag
\end{align}
with $C_F=\tfrac43$, $C_A=3$ and $T_R=\tfrac12$. The first three are the
final-state branchings; the last is the first one read the other way round
and is needed for the initial state, where the parton that continues into the
hard process may be the gluon. Let me point at each as we pass. $P_{qq}$ and
$P_{gg}$ diverge as $1-z\to0$: that is the soft-gluon singularity of
\cref{eq:mother_gives_birth:propagator}, and $P_{gg}$ diverges at both ends
because either of the two gluons may be the soft one. $P_{qg}$ is finite:
a soft quark is not enhanced, which is why gluons split into quarks so much
less often than into gluons, and why in \cref{tab:mother_gives_birth:numbers}
three quarters of the final-state partons are gluons. The colour factors say
that a gluon radiates $C_A/C_F=9/4$ times more than a quark; a gluon jet is
broader and busier than a quark jet for exactly this reason, and every jet
tagger in the later chapters exploits it.

\Cref{tab:mother_gives_birth:kernels} evaluates the four kernels at the
$z$ values the shower of our event actually drew, so that the formulas
have faces. Two readings. Along a row, the gluon's $P_{gg}$ is always the
largest, by the factor $C_A/C_F$ and by its second pole. Down the
$q\to qg$ column, the kernel grows from $4.5$ for the mother's first
emission at $z=0.60$ to $145$ for the soft proposal at $z=0.98$ she made in
the first round; the latter was accepted and then lost, which is the
shower's way of saying that soft gluons are frequent but not important.

\begin{table}[htbp]\centering\small
\caption[The splitting functions at the $z$ of our event]{The four
splitting functions of \cref{eq:mother_gives_birth:splitting} at $z$
values drawn in the shower of our event; the bold entry is the kernel the
branching actually used. The kernels are rates, not probabilities, and
exceed one freely.}
\label{tab:mother_gives_birth:kernels}
\begin{tabular}{llrrrr}\toprule
branching & $z$ & $P_{qq}$ & $P_{gg}$ & $P_{qg}$ & $P_{gq}$ \\ \midrule
mother's first emission, $b\to bg$ (branching 5) & $0.601$ & $\mathbf{4.55}$ & $14.5$ & $0.26$ & $2.57$ \\
mother's first proposal, $b\to bg$ (round 1, lost) & $0.982$ & $\mathbf{145}$ & $326$ & $0.48$ & $1.36$ \\
$g\to gg$ (branching 6) & $0.275$ & $1.98$ & $\mathbf{19.3}$ & $0.30$ & $7.39$ \\
$g\to \bbbar$ (branching 10) & $0.350$ & $2.30$ & $15.7$ & $\mathbf{0.27}$ & $5.42$ \\
ISR $\bar u\leftarrow g$ (branching 1) & $0.779$ & $9.69$ & $23.9$ & $\mathbf{0.33}$ & $1.80$ \\
ISR $u\leftarrow u$ (branching 2) & $0.850$ & $\mathbf{15.3}$ & $35.7$ & $0.37$ & $1.61$ \\
ISR $g\leftarrow g$ (branching 3) & $0.206$ & $1.75$ & $\mathbf{25.7}$ & $0.34$ & $10.6$ \\
\bottomrule
\end{tabular}
\end{table}

One identity is worth checking yourself, because the standalone script relies
on it: the second form of $P_{gg}$ follows from the first by putting the three
terms over the common denominator $z(1-z)$ and recognising
$z^2+(1-z)^2+z^2(1-z)^2=(1-z+z^2)^2$. Exercise~2 asks for it.
\Cref{fig:mother_gives_birth:kernels_plot} draws the four kernels over the
whole range of $z$ with the $33$ final-state branchings of our event placed
on them, and the distribution of the $z$ values by kind: the quark
emissions cluster at $z>0.6$ where $P_{qq}$ is large, the gluon splittings
pile up at both ends where $P_{gg}$ has its two poles, and the five
$g\to q\bar q$ splittings sit where they fall, because $P_{qg}$ has no
preference.

\begin{figure}[htbp]\centering
\includegraphics[page=2,width=\textwidth]{PS/Standalone/theory_parton_showering.pdf}
\caption[The splitting functions and the $z$ values drawn in our event]{The splitting functions at work, from \file{plot_theory_PS.py}. (a) The four kernels of \cref{eq:mother_gives_birth:splitting} and the overestimate $2C_F/(1-z)$ of the veto algorithm; each dot is one of the $33$ final-state branchings of the seed-1 shower at its drawn $z$, the mother's two hardest emissions labelled. (b) The distribution of $z$ for the $38$ branchings by kind: soft daughters dominate the gluon splittings, the quark emissions of our event drew $z$ between $0.6$ and $1$.}
\label{fig:mother_gives_birth:kernels_plot}
\end{figure}

\begin{figure}[htbp]\centering
\includegraphics[page=1,width=\textwidth]{PS/Standalone/animate_standalone_parton_showering_stills.pdf}
\caption[The ladder of branchings of our event]{The first scene of the
film (\cref{sec:mother_gives_birth:animation}), stopped at branching 25 of
38. Left: the branchings of our event in decreasing \ptevol, orange squares
for the initial state, red and blue dots for the families of the $b$ and of
the $\bar b$. The rows of \cref{tab:mother_gives_birth:kernels} are rungs of
this ladder: the three orange squares at the top are the initial-state
branchings 1--3, the blue and red dots at rungs 4 and 5 are the first
emissions of the $\bar b$ and of the mother, rung 6 is the $g\to gg$ at
$25.6\GeV$ and rung 10 the $g\to\bbbar$ at $6.0\GeV$. Right: the
$(\eta,\phi)$ plane at that moment, with 27 final-state partons, 20 of them
gluons (circles) and 7 quarks (diamonds), joined by the colour dipoles that
do the radiating.}
\label{fig:mother_gives_birth:ladder}
\end{figure}

Now read the table and the figure together. The ladder is steep at the top
and flat at the bottom: the first ten branchings run from $53.9$ down to
$6.0\GeV$, three octaves, while the last eighteen all lie between
$1.02\GeV$ and the cut-off, less than one octave. That is the
$\mathrm{d}\theta^2/\theta^2$ measure of
\cref{eq:mother_gives_birth:collinear_limit} at work, helped by a coupling
that grows from $0.14$ at $30\GeV$ to $0.53$ at $1\GeV$: the number of
branchings per octave rises as the scale falls. Who branches is decided by
the kernels. Every branching adds a gluon to the final state, and a gluon
radiates with $C_A$ rather than $C_F$ from each of its two colour lines, so
down the ladder the gluons take over the radiating: at rung 25 twenty of
the 27 partons are gluons, at the end 27 of 40, and 19 of the 38
branchings are $g\to gg$. The finite $P_{qg}$ keeps the $g\to q\bar q$
splittings to five, and because that kernel has no soft enhancement its
$z$ values, $0.35$ for the $g\to\bbbar$ of rung 10, sit in the middle of
the range rather than near the ends. The two quark emissions at rungs 4
and 5 both drew $z\approx0.6$, where $P_{qq}$ is a modest $4.5$; the soft
end of the kernel, where it reaches $145$ at $z=0.98$, produced the
mother's round-1 proposal at $5.2\GeV$, which lost. Soft emissions are the
most frequent, but in a shower ordered in \pt they come late.

\subsection{The Sudakov form factor}
\label{sec:mother_gives_birth:sudakov}

Start at the moment before the shower begins. Nothing has been emitted, so
for each of the four candidates, the $b$, the $\bar b$ and the two incoming
partons, the probability of having radiated is zero and the probability of
not having radiated is exactly one. Now let the scale run down from
$113\GeV$ and ask how that no-emission probability falls for each of them.
Guess before you look at \cref{fig:mother_gives_birth:survival}: which
curve falls fastest, and which two should lie on top of each other? The
$b$ and the $\bar b$ carry the same colour charge and sit in dipoles of the
same mass, $377.9\GeV$, so their curves must coincide. The two beams
differ: the $u$ at $x=0.39$ is a valence quark whose density falls when it
is asked for a larger $x$, while the $\bar u$ at $x=2.2\times10^{-3}$ can
be resolved out of a gluon density more than fifteen times larger. The end
whose curve falls fastest radiates first most often, and the figure,
measured by letting each end propose $1500$ times from $\pt^{\max}$,
confirms the guess: the $+z$ beam falls fastest and wins $53\%$ of the
first branchings, the $b$ and $\bar b$ curves lie on top of each other and
still stand at $0.8$ when the scale has fallen to $10\GeV$, and the dashed
product of all four is the probability that nobody has radiated yet. The
quantity on the vertical axis is the Sudakov form factor, and here is its
definition.

\begin{figure}[htbp]\centering
\includegraphics[page=2,width=\textwidth]{PS/Standalone/animate_standalone_parton_showering_stills.pdf}
\caption[Measured no-emission probabilities of the four radiators]{The
second scene of the film (\cref{sec:mother_gives_birth:animation}). Left:
the probability that an end has not yet radiated when the scale has fallen
from $\pt^{\max}=113.3\GeV$ to \pt, measured from $1500$ proposals per end
with the standalone shower; the dashed curve is the product of the four.
Right: who made the hardest of the four proposals, that is who radiates
first, as a histogram of that first scale. The first branching of the
seed-1 event, at $53.9\GeV$, is marked.}
\label{fig:mother_gives_birth:survival}
\end{figure}

Now the central object of the chapter. Take a parton produced at some hard
scale and ask: what is the probability that it does \emph{not} branch while
its evolution variable $t$ falls from $t_{\max}$ down to some $t$? Call the
branching probability per unit $\ln t$ at scale $t$
\begin{equation}
  \frac{\mathrm{d}\mathcal P}{\mathrm{d}\ln t}
  =\frac{\alphas(t)}{2\pi}\int_{z_{\min}}^{z_{\max}}\!\mathrm{d}z\,P(z)
  \equiv\frac{\alphas(t)}{2\pi}\,\hat P(t) ,
  \label{eq:mother_gives_birth:dP}
\end{equation}
where the $z$ limits depend on $t$ through the kinematics and keep the
integral finite. Divide the interval into many small steps; the probability
of surviving each step is one minus the probability of branching in it;
multiply them, take the limit, and you have
\begin{equation}
  \Delta(t_{\max},t)=\exp\left[-\int_{t}^{t_{\max}}\frac{\mathrm{d}t'}{t'}\,
  \frac{\alphas(t')}{2\pi}\,\hat P(t')\right],
  \label{eq:mother_gives_birth:sudakov}
\end{equation}
the Sudakov form factor. It is the same mathematics as radioactive decay
with a time-dependent decay rate, and it carries the same lesson: the
probability of the \emph{first} branching happening at $t$ is the rate at
$t$ times the probability of having survived down to $t$,
\begin{equation}
  \frac{\mathrm{d}\mathcal P_{\mathrm{first}}}{\mathrm{d}\ln t}
  =\frac{\alphas(t)}{2\pi}\,\hat P(t)\;\Delta(t_{\max},t)
  =-\frac{\mathrm{d}\Delta(t_{\max},t)}{\mathrm{d}\ln t}.
  \label{eq:mother_gives_birth:first}
\end{equation}
Integrate the right-hand side from $t_{\max}$ down to the cut-off $t_{\min}$
and you get $1-\Delta(t_{\max},t_{\min})$: the branching probabilities of
all scales, plus the probability of never branching, add up to one. This is
\emph{unitarity}, and it is the property that makes the shower a
probabilistic process rather than a perturbative expansion. The singular
kernels of \cref{eq:mother_gives_birth:splitting} are not a problem any more:
a divergence in $\hat P$ makes the form factor vanish, which just says that a
parton certainly branches before reaching the singular region.

\begin{figure}[htbp]\centering
\resizebox{\textwidth}{!}{% ---------------------------------------------------------------------------
% The Sudakov form factor as a product of survival probabilities, step by
% step down the scale, and its two faces (branch at t / never branch) with
% the numbers of our event (-v log: 3513 trials, 433 proposals, 33 winners
% for FSR; "no emission" in round 2 for the mother's end; 69 silent
% competitors in round 39).
% ---------------------------------------------------------------------------
\begin{tikzpicture}[x=1cm,y=1cm,
  every node/.style={font=\footnotesize},
  lab/.style={font=\scriptsize,align=center},
  head/.style={font=\small\bfseries},
  axis/.style={-{Stealth},thick},
  stepbox/.style={draw,thick,fill=blue!8,minimum height=7mm,minimum width=20mm,font=\scriptsize,align=center},
  branch/.style={draw,thick,fill=red!15,minimum height=7mm,minimum width=16mm,font=\scriptsize,align=center},
  flow/.style={-{Stealth},thick}]
% ---------------- (a) the product of survivals --------------------------------------------------
\node[head,anchor=west] at (-0.3,2.4) {(a) a no-emission probability is a product of survivals};
\draw[axis] (-0.3,0) -- (11.2,0); \node[lab,below] at (10.3,-0.1) {$\ln t$ runs \emph{down} $\to$};
\foreach \i/\x/\t in {1/0.3/{$t_{\max}$},2/2.6/{$t_1$},3/4.9/{$t_2$},4/7.2/{$t_3$},5/9.5/{$\cdots$}} {
  \draw (\x,-0.1) -- (\x,0.1); \node[lab,below] at (\x,-0.1) {\t};}
\foreach \x/\lab in {1.45/{$1-\frac{\alphas\hat P}{2\pi}\delta\ln t\big|_{t_1}$},3.75/{$1-\frac{\alphas\hat P}{2\pi}\delta\ln t\big|_{t_2}$},6.05/{$1-\frac{\alphas\hat P}{2\pi}\delta\ln t\big|_{t_3}$},8.35/{$\cdots$}} {
  \node[stepbox] at (\x,0.9) {\lab};}
\node[lab] at (9.9,0.9) {$\times\;\cdots$};
\node[lab,text=black!65,text width=10.4cm,anchor=north west] at (-0.3,-0.75) {Survive the first step, \emph{and} the second, \emph{and} the third: multiply. In the limit $\delta\ln t\to0$ the product becomes $\exp[-\int\mathrm{d}\ln t\,\frac{\alphas}{2\pi}\hat P]=\Delta(t_{\max},t)$, \cref{eq:mother_gives_birth:sudakov}, the same mathematics as a decay with a time-dependent rate. The probability of the \emph{first} branching at $t$ is the rate there times the survival down to there, \cref{eq:mother_gives_birth:first}.};
% ---------------- (b) two faces -------------------------------------------------------------
\begin{scope}[xshift=13.4cm]
\node[head,anchor=west] at (-0.3,2.4) {(b) the two faces in the log of our event};
\node[stepbox,minimum width=30mm] (S) at (1.6,1.2) {a dipole end at scale $t_{\max}$};
\node[branch,minimum width=30mm] (B) at (-0.6,-0.6) {``proposes \dots\ pT = \dots''\\branches at $t$};
\node[stepbox,minimum width=30mm,fill=gray!12] (N) at (3.8,-0.6) {``no emission above\\the cut-off''};
\draw[flow] (S) -- (B) node[midway,left,lab] {$1-\Delta(t_{\max},t_{\min})$};
\draw[flow] (S) -- (N) node[midway,right,lab] {$\Delta(t_{\max},t_{\min})$};
\node[lab,text=black!65,text width=7.4cm,anchor=north west] at (-2.1,-1.35) {The two add up to one: unitarity. In round~2 the mother's own end took the right branch from $53.9\GeV$ straight through the cut-off; in round~39 all $69$ competitors did, and the shower stopped. Three levels of bookkeeping over the event: $3513$ trial scales, $433$ accepted proposals, $33$ winners (final state).};
\end{scope}
\end{tikzpicture}
}
\caption[The Sudakov form factor as a product of survival probabilities, and its two faces in the log]{(a) The no-emission probability as a product of survivals: divide the evolution in $\ln t$ into small steps, survive each with probability one minus the branching probability of the step, and multiply; the limit is the exponential of \cref{eq:mother_gives_birth:sudakov}. (b) The two faces of the form factor in the \code{-v} log of our event: every line that reads ``proposes'' is the branching face, every ``no emission above the cut-off'' is the surviving face, and the two add up to one. The counts are those of the seed-1 shower.}
\label{fig:mother_gives_birth:sudakov_sketch}
\end{figure}

The log shows both faces of the form factor. The ``branches at $t$'' face
is every line that reads ``proposes''; the ``never branches'' face is every
line that reads ``no emission above the cut-off'', and the second kind
appears already in round~2 of our event, where the mother's own dipole end
evolved from $53.9\GeV$ straight through the cut-off without an accepted
emission, and dominates the end: in round~39 all $69$ competitors ($67$
colour ends of the $40$ final partons plus the two beams) propose nothing,
and the shower stops. Over the whole event the veto algorithm generated
$3513$ final-state and $212$ initial-state trials, of which $433$ and $34$
survived their own acceptance tests and $33$ and $5$ won their round. Three
levels, then: a trial is a random number, an accepted proposal is a
possible emission, and only the winner of a round becomes a parton.

\paragraph{The double logarithm, once by hand.} Let me do for the mother
what we did with the Jacobian at the previous stop, because every shower
hides it. Keep only the soft-collinear part of $P_{qq}$, $2C_F/(1-z)$,
use as evolution variable the transverse momentum $\pt$ of the emission with
respect to the emitter, and take as $z$ limit the collinear kinematics
$1-z\ge\pt/Q$ for a parton of energy of order $Q$. Then
$\hat P=2C_F\ln(Q/\pt)$ and, for a fixed coupling,
\begin{equation}
\begin{split}
  -\ln\Delta(\pt^{\max},\pt)
  &=\frac{\alphas C_F}{\pi}\int_{\pt}^{\pt^{\max}}\frac{2\,\mathrm{d}q}{q}\ln\frac{Q}{q}
  =\frac{\alphas C_F}{\pi}\left[\ln^2\frac{Q}{\pt}-\ln^2\frac{Q}{\pt^{\max}}\right].
\end{split}
  \label{eq:mother_gives_birth:double_log_sudakov}
\end{equation}
Differentiate $\ln^2(Q/q)$ with respect to $q$ to check the integral; you
get $-2\ln(Q/q)/q$. For our mother the natural $Q$ is the mass of the
colour dipole she forms with the incoming $u$ quark,
$Q^2=(p_b+p_u)^2=2\mb^2-\hat t=1.43\times10^{5}\GeV^2$, that is
$Q=378\GeV$ (the invariants are in the Info Listing of the \CMSSW log,
\cref{sec:ancestral_home:cmssw}), and the shower starts at
$\pt^{\max}=113.3\GeV$ (\cref{sec:mother_gives_birth:scales}). With
$\alphas=0.118$ the probability that she emits nothing above $10\GeV$ is
\begin{equation*}
  \Delta(113.3,10)=\exp\left\{-0.0501\left[\ln^2 37.8-\ln^2 3.34\right]\right\}
  =\exp(-0.588)=0.56 ,
\end{equation*}
and with the running coupling of \cref{sec:mother_gives_birth:cp5} in the
integrand it drops to $0.48$. Scene~2 of the film (\cref{sec:mother_gives_birth:animation})
measured the same quantity from the standalone shower itself and found
about $0.8$ (\cref{fig:mother_gives_birth:sudakov_plot}). The difference is what the leading logarithm leaves out: the
full kernel is smaller than $2C_F/(1-z)$ away from $z=1$, the exact $z$
range of the dipole is narrower than $1-z\ge\pt/Q$, and an emission that
recoils against an incoming parton is further weighted by the PDF ratio
and by a damping factor described in \cref{sec:mother_gives_birth:fsr}.
The estimate is nevertheless worth having: it tells you at a glance that a
$100\GeV$ quark radiates a $10\GeV$ gluon about half the time, a
$1\GeV$ gluon almost always, and it shows why the evolution must be cut off
at some $\pt^{\min}$ where the coupling stops making sense.
\Cref{fig:mother_gives_birth:sudakov_plot} puts the estimate next to the
measured curves: the leading-logarithmic form factor with a fixed and with a
running coupling, both for the mother's dipole, against the survival
probabilities of all four radiators measured from $1500$ proposals each
with the shower script itself, the same measurement as scene~2 of the film.
The two leading-log curves fall below the measured one for the reasons
listed above, and all three agree that the mother emits nothing above
$30\GeV$ in most showers and nothing above $1\GeV$ almost never.

\begin{figure}[htbp]\centering
\includegraphics[page=3,width=\textwidth]{PS/Standalone/theory_parton_showering.pdf}
\caption[The measured Sudakov form factors against the leading-log formula]{The Sudakov form factor, from \file{plot_theory_PS.py}. (a) The probability that each of the four radiators of the hard event has not yet branched when the scale has fallen from $\pt^{\max}$ to \pt, measured from $1500$ proposals per radiator with the veto algorithm of the script; the dashed black curve is their product. For the mother's dipole the leading-logarithmic estimate \cref{eq:mother_gives_birth:double_log_sudakov} is drawn with a fixed coupling (dotted) and with the running coupling in the integrand (dash-dotted). (b) The scale of the first branching, the hardest of the four proposals, coloured by who made it; the $+z$ beam at small $x$ wins most often, as in the seed-1 event.}
\label{fig:mother_gives_birth:sudakov_plot}
\end{figure}

\subsection{From a form factor to a random number: the veto algorithm}
\label{sec:mother_gives_birth:veto}

The form factor gives a probability; the computer needs a scale. Picture
the mother at $113.3\GeV$ in round~1 with her survival curve of
\cref{fig:mother_gives_birth:survival} in front of her. To decide where she
branches next she should draw a uniform random number $R$, go to the height
$R$ on that curve and read off the \pt, exactly as one samples a decay time
from an exponential. The curve, however, is an integral over the full
kernel and the running coupling and has no closed form, so nothing can be
read off it. The way out is the oldest trick of Monte Carlo: replace the
curve by a simpler one that lies \emph{below} it everywhere, that is, one
that overestimates the branching rate; sample from that; then throw away
the surplus with a second random number. \Cref{fig:mother_gives_birth:veto}
shows the four steps with the mother's own numbers; keep it open while you
read them.

\begin{figure}[htbp]\centering
\resizebox{\textwidth}{!}{% ---------------------------------------------------------------------------
% The Sudakov veto algorithm as a flow chart, annotated with the mother's
% first proposal in round 1 of the seed-1 shower (-vv log): three draws,
% 59.67 -> 47.54 -> 5.24 GeV, the four weight factors and R' = 0.062.
% ---------------------------------------------------------------------------
\begin{tikzpicture}[x=1cm,y=1cm,
  every node/.style={font=\footnotesize},
  box/.style={draw,rounded corners=3pt,thick,fill=orange!12,align=center,inner sep=5pt,text width=4.6cm},
  decision/.style={draw,thick,diamond,aspect=2.6,fill=blue!8,align=center,inner sep=1pt,text width=3.6cm,font=\scriptsize},
  numbox/.style={draw,rounded corners=3pt,thin,fill=blue!6,align=left,inner sep=4pt,text width=6.9cm,font=\scriptsize},
  flow/.style={-{Stealth[length=2.5mm,width=2mm]},thick},
  no/.style={flow,red!70!black},
  yes/.style={flow,green!50!black},
  head/.style={font=\small\bfseries},
  lab/.style={font=\scriptsize,align=center}]
\node[head,anchor=west] at (-2.6,1.4) {the veto algorithm, one proposal};
\node[head,anchor=west] at (7.0,1.4) {the mother's proposal in round 1 (\code{-vv} log)};
\node[box] (A) at (0,0) {\textbf{start} at the current scale $t$ (first: $t_{\max}$)};
\node[box,below=6mm of A] (B) {\textbf{draw $R$}, solve $\Delta_{\mathrm{over}}(t,t')=R$ for the next trial scale $t'$ (\cref{eq:mother_gives_birth:veto_solution})};
\node[decision,below=6mm of B] (C) {$t'$ above the cut-off and kinematically allowed?};
\node[box,below=6mm of C] (D) {\textbf{draw $z$} from the overestimate kernel; compute the weight $w$ of \cref{eq:mother_gives_birth:veto_weight}};
\node[decision,below=6mm of D] (E) {draw $R'$: $R'\le w$?};
\node[box,below=6mm of E,fill=green!12] (F) {\textbf{accept}: this end proposes $t'$, $z$ (it still has to win the round)};
\node[box,fill=gray!12] (G) at (-6.4,0 |- C) {\textbf{no emission}: this end is silent in this round};
\draw[flow] (A) -- (B); \draw[flow] (B) -- (C); \draw[yes] (C) -- (D) node[midway,right,lab] {yes}; \draw[flow] (D) -- (E); \draw[yes] (E) -- (F) node[midway,right,lab] {yes};
\draw[no] (C) -- (G) node[pos=0.3,above,lab] {below the cut-off};
\draw[no] (E.east) -- ++(1.6,0) |- (A.east); \node[lab,text=red!70!black,anchor=west] at (4.3,0 |- E) {no: \textbf{continue from $t'$},\\not from $t_{\max}$};
\draw[no] (C.east) -- ++(1.0,0) |- ($(A.east)+(0,-0.25)$); \node[lab,text=red!70!black,anchor=west] at (4.3,0 |- C) {not allowed:\\continue from $t'$};
% the numbers
\node[numbox,anchor=west] (N1) at (7.0,0 |- B) {draw 1: $R=0.51607$ from $113.32\GeV$ $\to$ $t'^{1/2}=59.67\GeV$};
\node[numbox,anchor=west] (N2) at (7.0,0 |- C) {$59.67$: recoiler is the incoming $u$, would need $x'=2.65>1$: not allowed\\draw 2: $R=0.77668$ from $59.67$ $\to$ $47.54\GeV$: $x'=19$, not allowed\\draw 3: $R=0.03959$ from $47.54$ $\to$ $5.24\GeV$: allowed};
\node[numbox,anchor=west] (N3) at (7.0,0 |- D) {$z=0.9819$ from $2C_F/(1-z)$;\quad $w=0.9818\times0.8076\times0.2747\times0.9810=0.2137$\\(kernel ratio, coupling ratio, beam-recoil damping, PDF ratio)};
\node[numbox,anchor=west] (N4) at (7.0,0 |- E) {$R'=0.0623\le0.2137$: accepted};
\node[numbox,anchor=west] (N5) at (7.0,0 |- F) {the mother proposes $5.241\GeV$, $z=0.9819$; the $+z$ beam proposes $53.877\GeV$ and wins the round. Over the event: $3513$ draws, $433$ accepted, $33$ winners.};
\foreach \s/\n in {B/N1,C/N2,D/N3,E/N4,F/N5} {\draw[flow,gray!60] ($(\s.east)+(0.1,0)$) -- (\n.west);}
\end{tikzpicture}
}
\caption[The veto algorithm as a flow chart, with the mother's first proposal]{The veto algorithm as a flow chart. Left: one proposal of one radiator, from the current scale to an accepted trial or to silence; the two red loops are the instruction that makes the algorithm exact, to continue from the rejected scale and not from $t_{\max}$. Right: the mother's proposal in round~1 of the seed-1 shower from the \code{-vv} log, three draws and four weight factors long, which produced the $5.24\GeV$ that lost to the $+z$ beam.}
\label{fig:mother_gives_birth:veto_flow}
\end{figure}

\begin{figure}[htbp]\centering
\includegraphics[width=\textwidth]{PS/Standalone/veto_algorithm.pdf}
\caption[The veto algorithm step by step]{The four steps of the veto
algorithm for the mother's first proposal in round~1, drawn by
\file{plot\_veto\_algorithm.py} from the \code{-vv} log. Step~1: the kernel
$P_{qq}$ and its overestimate $2C_F/(1-z)$; the shaded surplus is removed in
step~4, and the two $z$ values of the mother's emissions are marked.
Step~2: the second-order coupling and its first-order overestimate, with
the ratio that becomes a weight; in the last octave above the cut-off the
two curves cross. Step~3: the overestimated form factor
$\Delta_{\mathrm{over}}$ for three draws in a row, each starting from the
scale the previous one reached: $R=0.516$ gives $59.7\GeV$, rejected
because the incoming $u$ would need $x'>1$; $R=0.777$ from $59.7\GeV$ gives
$47.5\GeV$, rejected for the same reason; $R=0.040$ from $47.5\GeV$ gives
$5.24\GeV$. The blue triangles are the eight draws the $\bar b$ needed in
the same round. Step~4: the four factors of the acceptance weight of the
$5.24\GeV$ trial, the light part of each bar being thrown away, and the
$R'=0.062$ that accepted it.}
\label{fig:mother_gives_birth:veto}
\end{figure}

Here is the part I would like you to do once with pencil and paper, because
it is the engine of every shower and it fits on half a page. We want to
sample the scale of the next branching from \cref{eq:mother_gives_birth:first},
that is, to solve $\Delta(t_{\max},t)=R$ for $t$ with $R$ uniform in
$[0,1]$. For the real $\hat P(t)$ and the real $\alphas(t)$ that cannot be
done in closed form. So replace both by overestimates that can be inverted,
and correct afterwards.

\paragraph{Step 1: overestimate the kernel.} Choose a constant
$\hat P_{\mathrm{over}}\ge\hat P(t)$ for all $t$ in the interval. For
$P_{qq}$ the script uses $2C_F/(1-z)$ integrated over the widest $z$ range
that occurs in the interval, $c\equiv2C_F\ln\frac{1-z_{\min}}{1-z_{\max}}$;
the missing factor $(1+z^2)/2\le1$ is applied later. For the mother's
dipole in round~1 the $z$ range is set by the lowest scale of the
five-flavour segment, $\pt=\mb$, on a dipole of corrected mass squared
$(377.9\GeV)^2-\mb^2=1.428\times10^5\GeV^2$:
$z_{\min}=\tfrac12-\sqrt{\tfrac14-23.04/1.428\times10^5}=1.6\times10^{-4}$,
$z_{\max}=1-z_{\min}$, and $c=\tfrac83\ln(0.99984/0.00016)=23.28$. The
top-left panel of \cref{fig:mother_gives_birth:veto} shows the two
curves; the band between them is what step~4 will throw away.

\paragraph{Step 2: overestimate the coupling.} Use the first-order form
\begin{equation}
  \alphas^{(1)}(t)=\frac{2\pi}{b_0\ln(t/\Lambda^2)},\qquad
  b_0=\frac{33-2n_f}{6},
  \label{eq:mother_gives_birth:alphas_1loop}
\end{equation}
which for the same $\Lambda$ lies above the second-order running (the
correction term is negative) and which has an elementary integral. The
log prints both at the scales that matter: at $34.2\GeV$, where the
$\bar b$ first radiates, $0.1633$ against the true $0.1386$; at $5.2\GeV$,
$0.2608$ against $0.2106$; at $1\GeV$, $0.725$ against $0.530$. The ratio,
$0.85$, $0.81$ and $0.73$, is the price of the overestimate and returns as
a weight in step~4. In the last octave above the cut-off the two curves
cross: at $0.611\GeV$ the first-order form gives $1.485$ and the
second-order $1.563$, a ratio of $1.052$, and the ``overestimate'' is no
longer one. The top-right panel shows the two curves crossing there.

\paragraph{Step 3: invert.} With these two,
\begin{equation}
  -\ln\Delta_{\mathrm{over}}(t_{\max},t)
  =\int_t^{t_{\max}}\frac{\mathrm{d}t'}{t'}\,\frac{c}{b_0\ln(t'/\Lambda^2)}
  =\frac{c}{b_0}\ln\frac{\ln(t_{\max}/\Lambda^2)}{\ln(t/\Lambda^2)},
\end{equation}
because $\mathrm{d}t'/t'=\mathrm{d}\ln t'$ and
$\int\mathrm{d}u/u=\ln u$ with $u=\ln(t'/\Lambda^2)$. Setting this equal to
$-\ln R$ and solving for $t$,
\begin{equation}
  \ln\frac{t}{\Lambda^2}=\ln\frac{t_{\max}}{\Lambda^2}\;R^{\,b_0/c}
  \qquad\Longrightarrow\qquad
  t=\Lambda^2\left(\frac{t_{\max}}{\Lambda^2}\right)^{R^{\,b_0/c}} .
  \label{eq:mother_gives_birth:veto_solution}
\end{equation}
This single line is what the standalone script executes at every trial
(\code{pT2next\_fsr}, \code{pT2next\_isr}); find it on GitHub and you will
recognise $\Lambda^2$, $b_0$ and $c$. Check the limits: $R=1$ gives
$t=t_{\max}$ (no evolution), $R\to0$ gives $t\to\Lambda^2$ (the coupling
blows up), and a large $c$, that is a strongly radiating parton, pushes $t$
towards $t_{\max}$. Now do it with the mother's numbers of round~1,
$t_{\max}=113.32^2\GeV^2$, $\Lambda_5=0.2262\GeV$, $b_0=23/6=3.833$ and
$c=23.28$, so that $b_0/c=0.1646$. The \code{-vv} log drew
$R=0.51607$: $R^{0.1646}=0.8968$, $\ln(t_{\max}/\Lambda_5^2)=12.433$, and
\cref{eq:mother_gives_birth:veto_solution} gives
$\sqrt t=0.2262\,\mathrm{e}^{0.8968\times12.433/2}=59.67\GeV$. One random
number has taken the mother from $113.3$ to $59.7\GeV$; the next two, in
the same log, took her from $59.7$ to $47.5$ and from $47.5$ to
$5.24\GeV$. Why three draws for one proposal is the business of step~4.
The bottom-left panel draws the three curves $\Delta_{\mathrm{over}}$ of
these draws, each one starting where the previous one stopped; a single
curve from $113\GeV$ would have been the wrong distribution.

\paragraph{Step 4: choose $z$ and correct.} Draw $z$ from the overestimate
kernel (for $2C_F/(1-z)$ that is $1-z=(1-z_{\min})[(1-z_{\max})/(1-z_{\min})]^{R'}$),
then accept the trial with probability
\begin{equation}
  w=\frac{P(z)}{P_{\mathrm{over}}(z)}\times\frac{\alphas(t)}{\alphas^{(1)}(t)}
  \times(\text{any further factor}\le1),
  \label{eq:mother_gives_birth:veto_weight}
\end{equation}
and if it is rejected, \emph{continue the evolution from the rejected $t$},
not from $t_{\max}$. That last instruction is the veto algorithm proper, and
it is the step people get wrong. Restarting from $t_{\max}$ would sample
$\Delta_{\mathrm{over}}$ and then thin it; continuing from the rejected scale
samples exactly $\Delta$ with the true kernel, because the survival
probability of the true process factorises through every intermediate scale.
Exercise~4 asks you to prove it for a two-step interval. All further physics,
the second-order coupling, the mass term of the heavy quark, the PDF ratio
of an incoming recoiler, is inserted through the weight of
\cref{eq:mother_gives_birth:veto_weight}, which is why the script can count
how often a weight exceeded one (14 times in our event, all from the
coupling ratio in the last octave above the cut-off, where the second-order
running is no longer below the first-order form).

Here is the mother's first proposal, trial by trial, from the \code{-vv}
log. The first draw, $59.7\GeV$, was rejected before any weight was
computed: the recoiler is the incoming $u$, and absorbing that virtuality
would have needed $x'=2.65$, more than the whole proton. The second,
$47.5\GeV$, failed the same test with $x'=19$. The evolution continued
from $47.5\GeV$, not from $113\GeV$, and the third draw gave $5.24\GeV$
with $z=0.9819$ from the overestimate kernel. Then the weight of
\cref{eq:mother_gives_birth:veto_weight}, factor by factor:
\begin{equation*}
  w=\underbrace{0.9818}_{(1+z^2)/2\text{, dead cone }0.0003}
  \times\underbrace{0.8076}_{0.2106/0.2608}
  \times\underbrace{0.2747}_{w_{\mathrm{damp}}\text{, \cref{eq:mother_gives_birth:dampen}}}
  \times\underbrace{0.9810}_{xf_u(0.3944)/xf_u(0.3901)}
  =0.2137 ,
\end{equation*}
against which $R'=0.0623$ was drawn: accepted (the bottom-right panel of
\cref{fig:mother_gives_birth:veto}). That is the ``$5.241\GeV$,
$z=0.9819$'' the mother proposed in round~1, and it lost to the
$53.9\GeV$ of the $+z$ beam. The same three-level bookkeeping, trials,
accepted proposals, winners, gave the $3513/433/33$ of
\cref{sec:mother_gives_birth:sudakov}.

\section{\PYTHIA's \texorpdfstring{\pt}{pT}-ordered dipole shower}
\label{sec:mother_gives_birth:dipole_shower}

What is the evolution variable $t$? Angle, virtuality and transverse
momentum have all been used; \PYTHIA~8 uses the transverse momentum of the
branching~\cite{Sjostrand:2004ef,Bierlich:2022pfr}, and its radiators are
not single partons but colour dipoles. The standalone script of
\cref{sec:mother_gives_birth:standalone_code} follows that construction
closely enough that its output can be compared with the generator's, and
this section describes the construction with the script's variable names in
brackets, so that you can read the two side by side.

\subsection{Final-state radiation}
\label{sec:mother_gives_birth:fsr}

Whom does the mother radiate against? The natural guess is her twin, the
$\bar b$ she was born with, and it is wrong: the two quarks of our hard
process are not colour-connected to each other at all, and the mother's
partner is a parton in the beam. That one fact decides the recoil of every
one of her emissions, so it has to be settled before the kinematics.
Start from what the two colour columns of \file{event_ME.lhe} record. A
quark carries one of $N_c=3$ colours, an antiquark one of three
anticolours, and a gluon one colour and one anticolour, eight of the nine
combinations, since the colour-singlet one is not a gluon. A generator does
not propagate SU(3) amplitudes through forty branchings. It keeps the
leading term of the expansion in $1/N_c$, in which the colour structure of
any diagram is a set of \emph{lines}, each carrying one colour index
unchanged from where it starts to where it ends: a quark is one line
pointing along its motion, an antiquark one line pointing against it, and
a gluon is a pair of lines with opposite arrows, 't~Hooft's double-line
notation. The Les Houches record numbers the lines from 101 upwards, and
the two colour columns of a parton are the tags of the line and the
anti-line it holds. For our hard process the assignment follows from one
identity. The $s$-channel gluon couples the two quark currents through
\begin{equation}
  T^a_{ij}\,T^a_{kl}
  =\frac12\left(\delta_{il}\delta_{jk}-\frac{1}{N_c}\,\delta_{ij}\delta_{kl}\right),
  \label{eq:mother_gives_birth:fierz}
\end{equation}
where $j$ and $i$ are the colour indices of the incoming $u$ and $\bar u$
and $k$ and $l$ those of the outgoing $b$ and $\bar b$. The first term
identifies the colour of the $u$ with that of the $b$ and the anticolour
of the $\bar u$ with that of the $\bar b$; the second, a colour-singlet
exchange suppressed by $1/N_c$, is dropped. That is
\cref{fig:mother_gives_birth:dipoles}(a): colour~102 runs from the
incoming $u$ through the virtual gluon to the $b$, anticolour~101 from the
incoming $\bar u$ to the $\bar b$, and nothing connects the mother to her
twin. Her colour partner is a beam parton. The price of the approximation
is of order $1/N_c^2$, and \PYTHIA recovers part of it by letting a quark
end radiate with $C_F=4/3$ rather than with the strict large-$N_c$ value
$N_c/2$.

Every final-state parton with a colour line thus has a partner at the
other end of that line, and a gluon has two lines and therefore two
partners. Each such pair is a \emph{dipole}, and each end of a dipole is a
radiator, with the other end as its \emph{recoiler} [\code{dipole\_ends}].
The initial-state shower rethreads the lines too.
\Cref{fig:mother_gives_birth:dipoles}(b) shows branching~1
(\cref{sec:mother_gives_birth:isr}): backward evolution finds that the
$\bar u$ came out of a gluon, and a gluon holds two lines, the 101 that the
$\bar u$ already had, which now begins at the beam gluon \#7, and a new
103 that leaves with the emitted $u$ \#8. The hard process is untouched,
but the dipole 101 has a new end and its mass has grown from $377.9$ to
$420.5\GeV$ (step~2 of the log). The dipole picture is what makes the
shower coherent: a soft gluon at a wide angle sees the total colour charge
of the pair, not of one parton, and interference between the two ends is
built in (\cref{fig:mother_gives_birth:dipoles}(e)). The same lines
outlive the shower: hadronisation will stretch a string along every one of
them (\cref{ch:dressing_children}).

\begin{figure}[htbp]\centering
\resizebox{\textwidth}{!}{% ---------------------------------------------------------------------------
% Colour flow and colour chains of our event. Panels (a) and (b) draw the
% hard process and the first initial-state branching in 't Hooft's double-line
% (colour-flow) notation next to the chains the shower keeps; panels (c) and
% (d) show the chains before and after the first two final-state branchings
% (branchings 4 and 5 of the seed-1 shower); panel (e) is colour coherence.
% Every tag, row number and mass is taken from the -v log of
% standalone_parton_showering.py (example/PS/Standalone/standalone_parton_showering.log,
% sections 2 and 4, steps 1 to 5). Input from chapter_02_mother_gives_birth.tex
% inside a figure environment.
% ---------------------------------------------------------------------------
\begin{tikzpicture}[x=1cm,y=1cm,
  every node/.style={font=\footnotesize},
  parton/.style={draw,circle,minimum size=8mm,inner sep=0pt,fill=white,thick,align=center},
  gluon/.style={parton,double,double distance=1.2pt},
  quark/.style={parton},
  beam/.style={draw,rectangle,rounded corners=1.5pt,minimum height=8mm,minimum width=12mm,inner sep=1pt,fill=gray!15,thick,align=center},
  new/.style={fill=yellow!40},
  line/.style={line width=1.8pt},
  tag/.style={font=\scriptsize,fill=white,inner sep=1pt},
  role/.style={font=\scriptsize\itshape,text=black!65,align=center},
  lab/.style={font=\scriptsize,align=center},
  head/.style={font=\small\bfseries},
  % colour-flow lines: one colour index each, arrow along the index
  % (forwards along a quark, backwards along an antiquark)
  cline/.style={line width=1.4pt,rounded corners=2.5pt},
  marks/.style={postaction={decorate},decoration={markings,
     mark=at position #1 with {\arrow[scale=0.9]{Stealth}}}},
  ghost/.style={gray!45,decorate,decoration={coil,aspect=0,segment length=4pt,amplitude=1.4pt}},
  bigarrow/.style={-{Stealth[length=3mm,width=3mm]},gray!60,line width=2.4pt}]
\colorlet{c101}{blue!70!black}
\colorlet{c102}{red!80!black}
\colorlet{c103}{green!55!black}
\colorlet{c104}{orange!90!black}
\colorlet{c105}{violet}
\colorlet{c106}{teal}
\colorlet{c107}{brown!90!black}
% ---------------- (a) the hard process in colour-flow notation --------------------------
\node[head,anchor=west] at (-0.7,2.9) {(a) the hard process $\bar u u\to g^*\to\bbbar$: two colour lines, two dipoles};
\coordinate (V1) at (2.3,0);
\coordinate (V2) at (5.1,0);
\draw[ghost] (V1) -- (V2);
% colour 102: incoming u -> lower edge of the gluon -> b
\draw[cline,c102,marks=0.05,marks=0.5,marks=0.95]
   (0.9,-1.5) -- ($(V1)+(0,-0.13)$) -- ($(V2)+(0,-0.13)$) -- (6.7,-1.5);
% anticolour 101: bbar -> upper edge of the gluon -> incoming ubar (arrow against the fermion)
\draw[cline,c101,marks=0.05,marks=0.5,marks=0.95]
   (6.7,1.5) -- ($(V2)+(0,0.13)$) -- ($(V1)+(0,0.13)$) -- (0.9,1.5);
\node[tag,text=c102] at (1.8,-1.07) {102};
\node[tag,text=c101] at (1.8,1.07) {101};
\node[tag,text=c102] at (5.72,-1.07) {102};
\node[tag,text=c101] at (5.72,1.07) {101};
\node[lab,anchor=east] at (0.8,-1.5) {$u$ \#4\\from the $-z$ proton};
\node[lab,anchor=east] at (0.8,1.5) {$\bar u$ \#3\\from the $+z$ proton};
\node[lab,anchor=west] at (6.8,-1.5) {$b$ \#6, the mother};
\node[lab,anchor=west] at (6.8,1.5) {$\bar b$ \#5};
\node[lab] at (3.7,0.75) {$g^*$: one line of each kind};
\node[lab,text=black!65] at (3.7,-0.75) {$\sqrt{\shat}=396\GeV$};
% the same as chains
\draw[bigarrow] (8.2,0) -- (9.0,0);
\node[role] at (8.6,0.5) {as the shower\\stores it};
\node[beam] (a3) at (10.4,1.0) {$\bar u$ \#3\\[-1pt]{\tiny beam $+z$}};
\node[quark] (a5) at (13.8,1.0) {$\bar b$\\[-2pt]\tiny\#5};
\node[beam] (a4) at (10.4,-1.0) {$u$ \#4\\[-1pt]{\tiny beam $-z$}};
\node[quark] (a6) at (13.8,-1.0) {$b$\\[-2pt]\tiny\#6};
\draw[line,c101] (a3) -- (a5) node[tag,midway,below] {101} node[role,midway,above] {$m_{\mathrm{dip}}=377.9\GeV$};
\draw[line,c102] (a4) -- (a6) node[tag,midway,below] {102} node[role,midway,above] {$m_{\mathrm{dip}}=377.9\GeV$};
\node[role] at (12.5,-1.85) {2 dipoles: the 2 final-state ends radiate (FSR),\\the 2 beam ends evolve backwards (ISR)};
% ---------------- (b) after branching 1 (ISR) -------------------------------------------
\begin{scope}[yshift=-5.4cm]
\node[head,anchor=west] at (-0.7,2.9) {(b) after branching 1 (ISR at $\ptevol=53.9\GeV$): the $\bar u$ came out of a gluon};
\coordinate (V0) at (2.0,1.4);
\coordinate (V1) at (3.4,0);
\coordinate (V2) at (5.4,0);
\draw[ghost] (V1) -- (V2);
\draw[ghost] (0.4,1.4) -- (V0);
% colour 102 as before
\draw[cline,c102,marks=0.05,marks=0.5,marks=0.95]
   (1.7,-1.5) -- ($(V1)+(0,-0.13)$) -- ($(V2)+(0,-0.13)$) -- (7.1,-1.5);
% anticolour 101: bbar -> gluon -> ubar leg -> lower edge of the beam gluon #7
\draw[cline,c101,marks=0.05,marks=0.37,marks=0.66,marks=0.95]
   (7.1,1.5) -- ($(V2)+(0,0.13)$) -- ($(V1)+(0,0.13)$) -- ($(V0)+(0,-0.13)$) -- (0.4,1.27);
% colour 103: upper edge of the beam gluon #7 -> emitted u #8
\draw[cline,c103,marks=0.3,marks=0.8]
   (0.4,1.53) -- ($(V0)+(0,0.13)$) -- (3.5,2.25);
\node[tag,text=c102] at (2.74,-1.07) {102};
\node[tag,text=c102] at (6.07,-1.07) {102};
\node[tag,text=c101] at (6.07,1.07) {101};
\node[tag,text=c101] at (2.45,0.5) {101};
\node[tag,text=c103] at (2.6,2.1) {103};
\node[lab,anchor=east] at (1.6,-1.5) {$u$ \#4 from the $-z$ proton};
\node[lab,anchor=north east] at (1.3,1.15) {$g$ \#7 from the\\$+z$ proton};
\node[lab,anchor=west] at (3.6,2.25) {$u$ \#8 emitted, $45\GeV$};
\node[lab,anchor=west] at (3.05,0.95) {$\bar u$ \#3, spacelike};
\node[lab,anchor=west] at (7.2,-1.5) {$b$ \#10};
\node[lab,anchor=west] at (7.2,1.5) {$\bar b$ \#9};
\node[lab] at (4.4,0.45) {$g^*$};
\draw[bigarrow] (8.5,0) -- (9.3,0);
\node[beam] (b7) at (10.4,1.0) {$g$ \#7\\[-1pt]{\tiny beam $+z$}};
\node[quark] (b8) at (13.8,2.2) {$u$\\[-2pt]\tiny\#8};
\node[quark] (b9) at (13.8,1.0) {$\bar b$\\[-2pt]\tiny\#9};
\node[beam] (b4) at (10.4,-1.0) {$u$ \#4\\[-1pt]{\tiny beam $-z$}};
\node[quark] (b10) at (13.8,-1.0) {$b$\\[-2pt]\tiny\#10};
\draw[line,c103] (b7) -- (b8) node[tag,pos=0.35,below,sloped] {103} node[role,pos=0.68,above,sloped] {$114.6\GeV$};
\draw[line,c101] (b7) -- (b9) node[tag,midway,above] {101} node[role,midway,below] {$420.5\GeV$\\(was $377.9$)};
\draw[line,c102] (b4) -- (b10) node[tag,midway,below] {102} node[role,midway,above] {partner unchanged};
\node[role] at (12.1,-1.85) {3 final-state partons, 3 dipole ends;\\the beam gluon holds two lines};
\end{scope}
% ---------------- notation box, right of (a) and (b) -------------------------------------
\begin{scope}
\fill[gray!8,rounded corners=4pt] (15.6,2.75) rectangle (21.9,-7.3);
\node[head,anchor=west] at (15.8,2.4) {how to read the colour lines};
\draw[cline,c102,marks=0.6] (15.9,1.75) -- (17.1,1.75);
\node[lab,anchor=west] at (17.25,1.75) {quark: one colour, arrow forwards};
\draw[cline,c101,marks=0.6] (17.1,1.15) -- (15.9,1.15);
\node[lab,anchor=west] at (17.25,1.15) {antiquark: one anticolour, arrow backwards};
\draw[ghost] (15.9,0.55) -- (17.1,0.55);
\draw[cline,c103,marks=0.6] (15.9,0.68) -- (17.1,0.68);
\draw[cline,c101,marks=0.6] (17.1,0.42) -- (15.9,0.42);
\node[lab,anchor=west] at (17.25,0.55) {gluon: one of each, two lines};
\node[font=\scriptsize,align=left,anchor=north west,text width=5.6cm] at (15.8,0.05)
  {A quark carries one of $N_c=3$ colours, an antiquark one of three
   anticolours, a gluon one colour and one anticolour. The generator keeps
   only the leading term in $1/N_c$: every diagram becomes a set of lines,
   each carrying one colour index unchanged from where it starts to where
   it ends, and a gluon is a pair of lines with opposite arrows ('t~Hooft's
   double-line notation). The Les Houches record numbers the lines from
   101 upwards; the two colour columns of a parton are the tags of the
   line and the anti-line it holds. A \emph{dipole} is one line with its
   two ends, and the chains on the left list, parton by parton, which
   lines each one holds. Hadronisation will later stretch a string along
   every one of these lines.};
\end{scope}
% ---------------- (c) before branching 4: after the three ISR branchings ------------------
\begin{scope}[yshift=-11.3cm]
\node[head,anchor=west] at (-0.7,3.3) {(c) after the three ISR branchings: 5 final partons, 7 dipole ends};
\node[beam] (b16) at (0,1.0) {$g$ \#16\\[-1pt]{\tiny beam $+z$}};
\node[quark] (u19) at (3.6,2.4) {$u$\\[-2pt]\tiny\#19};
\node[gluon] (g17) at (3.6,1.0) {$g$\\[-2pt]\tiny\#17};
\node[quark] (bb20) at (7.6,1.0) {$\bar b$\\[-2pt]\tiny\#20};
\node[beam] (u11) at (0,-1.2) {$u$ \#11\\[-1pt]{\tiny beam $-z$}};
\node[gluon] (g18) at (3.6,-1.2) {$g$\\[-2pt]\tiny\#18};
\node[quark] (b21) at (7.6,-1.2) {$b$\\[-2pt]\tiny\#21};
\draw[line,c103] (b16) -- (u19) node[tag,midway,above,sloped] {103};
\draw[line,c105] (b16) -- (g17) node[tag,midway,below] {105};
\draw[line,c101] (g17) -- (bb20) node[tag,midway,below] {101} node[role,midway,above] {$m_{\mathrm{dip}}=843\GeV$};
\draw[line,c104] (u11) -- (g18) node[tag,midway,below] {104};
\draw[line,c102] (g18) -- (b21) node[tag,midway,below] {102} node[role,midway,above] {$m_{\mathrm{dip}}=80.6\GeV$};
\node[role,above=1pt of bb20] {radiator of branching 4\\(anticolour 101)};
\node[role,below=1pt of g17] {recoiler of\\branching 4};
\node[role,below=1pt of b21] {radiator of branching 5\\(colour 102)};
\node[role,below=1pt of g18] {recoiler of\\branching 5};
\end{scope}
% ---------------- (d) after branchings 4 and 5 -------------------------------------------
\begin{scope}[yshift=-17.7cm]
\node[head,anchor=west] at (-0.7,3.3) {(d) after branchings 4 and 5: 7 final partons, 11 dipole ends};
\node[beam] (b16) at (0,1.0) {$g$ \#16\\[-1pt]{\tiny beam $+z$}};
\node[quark] (u19) at (3.0,2.4) {$u$\\[-2pt]\tiny\#19};
\node[gluon] (g24) at (3.0,1.0) {$g$\\[-2pt]\tiny\#24};
\node[gluon,new] (g23) at (5.9,1.0) {$g$\\[-2pt]\tiny\#23};
\node[quark] (bb22) at (8.8,1.0) {$\bar b$\\[-2pt]\tiny\#22};
\node[beam] (u11) at (0,-1.2) {$u$ \#11\\[-1pt]{\tiny beam $-z$}};
\node[gluon] (g27) at (3.0,-1.2) {$g$\\[-2pt]\tiny\#27};
\node[gluon,new] (g26) at (5.9,-1.2) {$g$\\[-2pt]\tiny\#26};
\node[quark] (b25) at (8.8,-1.2) {$b$\\[-2pt]\tiny\#25};
\draw[line,c103] (b16) -- (u19) node[tag,midway,above,sloped] {103};
\draw[line,c105] (b16) -- (g24) node[tag,midway,below] {105};
\draw[line,c101] (g24) -- (g23) node[tag,midway,below] {101} node[role,midway,above] {$534\GeV$};
\draw[line,c106] (g23) -- (bb22) node[tag,midway,below] {106} node[role,midway,above] {$69.8\GeV$};
\draw[line,c104] (u11) -- (g27) node[tag,midway,below] {104};
\draw[line,c102] (g27) -- (g26) node[tag,midway,above] {102} node[role,midway,below] {$4.5\GeV$};
\draw[line,c107] (g26) -- (b25) node[tag,midway,above] {107} node[role,midway,below] {$65.7\GeV$};
\node[role,above=1pt of g23] {emitted gluon (status 51), $67\GeV$,\\inserted into line 101};
\node[role,below=1pt of g24] {recoiler's\\new copy (52)};
\node[role,above=1pt of bb22] {$\bar b$'s new\\copy (51)};
\node[role,above=1pt of g26] {emitted gluon (51), $38\GeV$,\\inserted into line 102};
\node[role,below=1pt of g27] {recoiler's\\new copy (52)};
\node[role,below=1pt of b25] {mother's new\\copy (51)};
\end{scope}
% ---------------- (e) colour coherence ------------------------------------------------------
\begin{scope}[xshift=12.9cm,yshift=-13.9cm]
\node[head,anchor=west] at (-1.2,3.6) {(e) what a soft gluon sees};
\coordinate (O) at (0,0);
\coordinate (B) at (20:4.0);
\coordinate (G) at (-32:4.0);
\fill[red!8] (O) -- (46:3.8) arc[start angle=46,end angle=-6,radius=3.8] -- cycle;
\fill[red!8] (O) -- (-6:3.8) arc[start angle=-6,end angle=-58,radius=3.8] -- cycle;
\draw[thick,c107] (O) -- (B) node[pos=1.1,font=\footnotesize] {$b$ \#25};
\draw[thick,c107,decorate,decoration={coil,aspect=0,segment length=4pt,amplitude=1.6pt}] (O) -- (G)
   node[pos=1.12,font=\footnotesize] {$g$ \#26};
\draw[-{Stealth}] (-32:1.5) arc[start angle=-32,end angle=20,radius=1.5];
\node[font=\scriptsize] at (-6:2.0) {$\theta_{bg}$};
\draw[thick,gray!70,decorate,decoration={coil,aspect=0,segment length=3pt,amplitude=1.3pt}] (O) -- (128:3.0)
   node[pos=1.25,font=\scriptsize,align=center,text=black!70] {soft gluon at $\theta\gg\theta_{bg}$:\\sees the pair's total\\charge, $C_F$};
\node[font=\scriptsize,align=left,anchor=north west,text width=5.4cm] at (-1.2,-2.7)
  {Inside the dipole 107 each end radiates mainly within a cone of about $\theta_{bg}$
   (shaded). A gluon at a wider angle cannot resolve the two ends; it sees the colour
   charge of the pair, which is that of the original $b$, $C_F$, not $C_F+C_A$. That is
   colour coherence, and the dipole construction has it built in.};
\end{scope}
\end{tikzpicture}
}
\caption[Colour flow and dipoles of our event]{The colour lines of our
event, first in colour-flow notation and then as the chains the shower
keeps, with the tags and row numbers of the record. (a)~The hard process
in double-line notation: the $s$-channel gluon carries one colour line and
one anticolour line, so the colour of the incoming $u$ (tag 102) passes
straight through to the $b$ and the anticolour of the $\bar u$ (101) to
the $\bar b$. These two lines are the two dipoles the shower starts from,
both of mass $377.9\GeV$. (b)~After the first initial-state branching: the
$\bar u$ is resolved as coming out of the gluon \#7, whose two lines are
the old 101, now anchored at the beam, and a new 103 that leaves with the
emitted $u$ \#8; the box explains the notation. (c)~After the three
initial-state branchings: the incoming gluon \#16 holds two lines, 103 to
the $u$ quark emitted in branching~1 and 105 to the gluon \#17 of
branching~3, which is tied by line 101 to the $\bar b$; the incoming $u$
\#11 is tied by 104 to the gluon \#18 of branching~2, which holds the
mother on line 102. Every end of every line is a radiator and the other
end its recoiler. (d)~After branchings 4 and 5: each emitted gluon has
inserted itself into the line of its radiator with a new tag (106, 107),
and the recoilers have become new copies with their old tags. (e)~Colour
coherence: a soft gluon at an angle larger than the opening angle of a
dipole cannot resolve its two ends and sees their summed colour charge.}
\label{fig:mother_gives_birth:dipoles}
\end{figure}

\Cref{fig:mother_gives_birth:dipoles}(c) draws the chains as they stand
when the first final-state branching is about to happen, after the three
initial-state branchings have each inserted a parton into them. Read the
log's line for branching~4 against it: ``radiator \#20 $\bar b$ radiates
on its acol line (tag 101) against recoiler \#17 $g$'' is the right end of
the upper chain radiating against its neighbour, and the dipole mass of
$843\GeV$ is the invariant mass of that pair. The mother, at the right end
of the lower chain, radiates against the gluon \#18 that the $-z$ beam
emitted in branching~2, a dipole of only $80.6\GeV$. Panel~(d) is the same
picture two branchings later and is the one to keep in mind for the rest
of the section: the emitted gluons sit \emph{between} their radiator and
recoiler on the chain, every new parton brings two new dipole ends, and
the number of competitors in a round grows by two with every final-state
branching, from seven ends to eleven here and to sixty-seven at the end.

A branching $a\to bc$ of the radiator is parametrised by its virtuality
$m_a$ and the energy fraction $z$ kept by $b$ in the dipole rest frame, and
the evolution variable is
\begin{equation}
  \ptevol^2=z(1-z)\,(m_a^2-m_{a,0}^2),
  \label{eq:mother_gives_birth:pt_evol_fsr}
\end{equation}
where $m_{a,0}$ is the on-shell mass of the radiator ($4.8\GeV$ for the $b$,
zero for a gluon). The mother's first emission, branching~5 of the log,
had $\ptevol=32.070\GeV$ and $z=0.6013$, so her virtuality was
\begin{equation*}
  m_a^2=\mb^2+\frac{\ptevol^2}{z(1-z)}
  =23.04+\frac{1028.49}{0.6013\times0.3987}=4313.0\GeV^2,
  \qquad m_a=65.67\GeV ,
\end{equation*}
the invariant mass of the $b$--gluon pair she became, each daughter on
shell. Check the small-angle limit: for a massless emitter with
energy $E$ and opening angle $\theta$, $m_a^2\simeq z(1-z)E^2\theta^2$, so
$\ptevol\simeq z(1-z)E\theta$, the transverse momentum of either
daughter with respect to the parent. The three kernels are those of
\cref{eq:mother_gives_birth:splitting}, and for the heavy quark the
$q\to qg$ kernel carries the extra mass term of
\cref{sec:mother_gives_birth:dead_cone}. After the branching the dipole is
reassembled: the two daughters and the recoiler share the dipole's momentum
with all three on shell, the recoiler keeping its direction in the dipole
frame and giving up the longitudinal momentum needed for the radiator's
virtuality [\code{fsr\_kinematics}].

Read the construction once with the log's numbers for the $\bar b$'s first
emission, branching~4 ($\ptevol=34.218\GeV$, $z=0.5939$,
$m_a=69.84\GeV$), whose recoiler is the final-state gluon of the third
initial-state branching. In the dipole rest frame, of mass
$843.0\GeV$, the daughter pair gets $E_P=424.41\GeV$, the $\bar b$
$E_1=zE_P=252.08\GeV$ and the gluon $(1-z)E_P=172.33\GeV$, and the
daughter's transverse momentum with respect to the pair's axis,
$\sqrt{E_1^2-\mb^2-p_{1z}^2}=34.15\GeV$, reproduces $\ptevol$ up to the
mass term. The quoted $z$ is a fraction in \emph{this} frame. Back in the
laboratory the accounts read
\begin{equation*}
  \underbrace{2449.447}_{\bar b}+\underbrace{79.117}_{\text{recoiler}}
  =\underbrace{1467.400}_{\bar b'}+\underbrace{982.588}_{g}
  +\underbrace{78.577}_{\text{recoiler}'}\ \GeV ,
\end{equation*}
balanced to $4\times10^{-12}\GeV$ in all four components, and the colour
line is rethreaded: the $\bar b$ carried anticolour~101; after the
branching the daughter $\bar b$ carries anticolour~106, the new gluon
$(106,101)$ and the recoiler keeps $(101,105)$. The gluon has inserted
itself into the existing line, and these new tags decide who radiates
against whom from now on (\cref{fig:mother_gives_birth:dipoles}(d)). One more lesson from the same two branchings:
there are three transverse momenta in the log and they are not the same
quantity (\cref{tab:mother_gives_birth:three_pts}). For the mother's
emission the evolution variable is $32.07\GeV$ but the daughter's
transverse momentum in the construction frame is only $6.25\GeV$, because
her $65.7\GeV$ virtuality uses up most of the $80.6\GeV$ dipole she shares
with a nearby gluon; the emitted gluon nevertheless has $38.2\GeV$ in the
laboratory. Decreasing evolution scales do not mean decreasing laboratory
transverse momenta.

\begin{table}[htbp]\centering\small
\caption[Three transverse momenta in one branching]{Three transverse
momenta that appear in the log of every final-state branching, for the
first emission of each $b$ quark; all in GeV.}
\label{tab:mother_gives_birth:three_pts}
\begin{tabular}{llrr}\toprule
quantity & meaning & $\bar b$, branching 4 & $b$, branching 5 \\ \midrule
$\ptevol$ & evolution variable that orders the shower & $34.218$ & $32.070$ \\
$\pt$ in the dipole frame & daughter's \pt relative to the pair's axis & $34.145$ & $6.246$ \\
gluon's laboratory \pt & relative to the beam, what a detector sees & $67.305$ & $38.234$ \\
\bottomrule
\end{tabular}
\end{table}

Two things happen when the recoiler is an incoming parton, as it is for both
$b$ quarks at the start of our shower. First, the incoming parton cannot
absorb recoil by changing direction; instead its momentum fraction grows,
$x\to x'=x\,(1+\delta)$ with $\delta=(m_a^2-m_{a,0}^2)/(m_{\mathrm{dip}}^2-m_{a,0}^2)$,
and the emission is weighted by the PDF ratio $xf(x')/xf(x)$ at the new
fraction, which for a valence quark at $x=0.39$ is below one. Our event
has one such emission that was actually executed, branching~26, a gluon
emitting at $0.820\GeV$ against the incoming gluon of the $+z$ side:
$\delta=0.00188$, the incoming fraction moves from $0.0171076$ to
$0.0171398$, and the $0.219\GeV$ the daughters gain is exactly what the
beam parton now brings in more. The ``$x$ after ISR'' of the summary line
therefore includes small shifts from final-state emissions too. Second,
\PYTHIA damps such emissions by the factor
\begin{equation}
  w_{\mathrm{damp}}=\frac{\ptrad\,\ptevol}
  {\ptrad\,\ptevol+m_a^2}
  \qquad(\code{TimeShower:dampenBeamRecoil}),
  \label{eq:mother_gives_birth:dampen}
\end{equation}
with $\ptrad$ the transverse momentum of the radiator itself,
which suppresses large virtualities of a parton whose colour partner is a
beam. For the mother's round-1 proposal, $\ptrad=113.2\GeV$,
$\ptevol=5.24\GeV$ and $m_a=39.6\GeV$, it is
$593/(593+1567)=0.275$, the smallest of the four factors in the weight of
\cref{sec:mother_gives_birth:veto}. This one factor turns out to be the single most sensitive ingredient
of the whole shower: switching it off doubles the number of final-state
branchings in \cref{tab:mother_gives_birth:numbers}
(\cref{sec:mother_gives_birth:pythia}).

\Cref{fig:mother_gives_birth:fsr_plot} shows both lessons of this subsection
over the whole event: for the $33$ final-state branchings the evolution
variable and the laboratory \pt of the emitted parton differ by a factor
of three either way as a matter of course, and the damping factor of
\cref{eq:mother_gives_birth:dampen} falls steeply for soft emissions at
large virtuality against a beam recoiler, which is what made the mother's
round-1 proposal so expensive.

\begin{figure}[htbp]\centering
\includegraphics[page=4,width=\textwidth]{PS/Standalone/theory_parton_showering.pdf}
\caption[Evolution variable against laboratory $\pt$, and the beam-recoil damping]{Final-state radiation, from \file{plot_theory_PS.py}. (a) For each of the $33$ final-state branchings of the seed-1 shower, the evolution variable $\ptevol$ against the laboratory \pt of the emitted parton read from the record, coloured by kind; the dotted line is equality, and the first emissions of the two $b$ quarks and the mother's second are labelled. (b) The factor $w_{\mathrm{damp}}$ of \cref{eq:mother_gives_birth:dampen} for a radiator of $\pt=113\GeV$ against a beam recoiler, as a function of $\ptevol$ for three values of $z$; the star is the mother's round-1 proposal with its $0.275$.}
\label{fig:mother_gives_birth:fsr_plot}
\end{figure}

\subsection{Initial-state radiation: backward evolution}
\label{sec:mother_gives_birth:isr}

Which of the two incoming partons of our event radiates first, and how? The
$\bar u$ with $x=2\times10^{-3}$ or the $u$ with $x=0.39$? Most people
guess the valence quark, because it carries more momentum; the answer is
the sea antiquark, in more than half of all showers, and not by emitting a
gluon but by turning out to have been a gluon. Keep the guess and see what
the formula says. The two incoming partons radiate too, and here the shower
runs backwards in time. The hard process fixed the parton $a$ entering it with momentum
fraction $x$ at the scale $\pt^{\max}$; the question is which parton $b$ at
a lower scale and larger fraction $x'=x/z$ branched into $a$ plus an emitted
parton $c$. The probability is the DGLAP kernel weighted by how many partons
$b$ there are to branch relative to how many $a$ there are to be found
\cite{Sjostrand:2004ef},
\begin{equation}
  \mathrm{d}\mathcal P_{a\leftarrow b}
  =\frac{\alphas(\pt^2)}{2\pi}\,
  \frac{\mathrm{d}\pt^2}{\pt^2+p_{\mathrm{T}0}^2}\,
  \mathrm{d}z\,P_{ba}(z)\,
  \frac{x'f_b(x',\pt^2)}{x f_a(x,\pt^2)},
  \qquad
  \ptevol^2=(1-z)\,Q^2 ,
  \label{eq:mother_gives_birth:isr}
\end{equation}
where $Q^2$ is the spacelike virtuality of $a$, and $p_{\mathrm{T}0}=2\GeV$
(\code{SpaceShower:pT0Ref}) regularises the coupling and the measure,
$\alphas(\pt^2)\to\alphas(\pt^2+p_{\mathrm{T}0}^2)$, so that the evolution can be
continued to $\pt^{\min}=0.42\GeV$ without the coupling diverging: at
$\ptevol=1\GeV$ the coupling is evaluated at $\sqrt5=2.24\GeV$ and is
$0.29$ instead of $0.53$, and the measure is a quarter of $\mathrm{d}\pt^2/\pt^2$. The
four channels [\code{isr\_channels}] are $q\leftarrow q$ (a gluon is
emitted), $q\leftarrow g$ (an antiquark is emitted), $g\leftarrow g$ and
$g\leftarrow q$ (a quark is emitted and the gluon continues). After the
branching the new incoming parton $b$ sits on the beam axis with fraction
$x/z$, the emitted parton $c$ takes the transverse momentum, and the whole
downstream system, the other incoming parton untouched, is boosted and
rotated so that everything balances [\code{isr\_kinematics}]. The
\emph{recoil is taken by the entire event}: every final-state parton, the
mother included, shifts. We will see her \pt jump by $20\GeV$ three times
before she has emitted anything herself.

\begin{figure}[htbp]\centering
\resizebox{\textwidth}{!}{% ---------------------------------------------------------------------------
% The three initial-state branchings of our event as one Feynman diagram.
% Row numbers, x values, scales and virtualities from the -v log of
% standalone_parton_showering.py (branchings 1 to 3). Time runs from left to
% right, the backward evolution of the shower from right to left.
% ---------------------------------------------------------------------------
\begin{tikzpicture}[x=1cm,y=1cm,
  every node/.style={font=\footnotesize},
  fermion/.style={thick,postaction={decorate},decoration={markings,mark=at position 0.55 with {\arrow[scale=1.1]{Stealth}}}},
  antifermion/.style={thick,postaction={decorate},decoration={markings,mark=at position 0.45 with {\arrowreversed[scale=1.1]{Stealth}}}},
  gluon/.style={thick,decorate,decoration={coil,aspect=0,segment length=4.5pt,amplitude=2pt}},
  vertex/.style={fill=black,circle,inner sep=1.5pt},
  stepno/.style={draw,circle,inner sep=1pt,font=\scriptsize\bfseries,fill=orange!30,minimum size=4mm},
  lab/.style={font=\scriptsize,align=center},
  xlab/.style={font=\scriptsize,align=center,text=blue!60!black}]
% hard process
\node[circle,draw,thick,fill=gray!20,minimum size=12mm] (H) at (6.8,0) {$\hat\sigma$};
\node[lab,below=1pt of H] {$\bar u u\to\bbbar$\\$\sqrt{\shat}=396\GeV$};
\draw[fermion] (H) -- (9.6,1.0) node[right] {$b$ \#6, $\pt=113\GeV$};
\draw[antifermion] (H) -- (9.6,-1.0) node[right] {$\bar b$ \#5, $\pt=113\GeV$};
% +z side (top): g(#16) -> g(#7) + g(#17);  g(#7) -> ubar(#3) + u(#8)
\coordinate (A) at (1.8,2.6);
\coordinate (B) at (4.4,1.6);
\draw[gluon] (-0.9,3.2) -- (A);
\node[xlab,anchor=south] at (0.1,3.5) {$g$ \#16 from the $+z$ proton\\$x''=1.36\times10^{-2}$};
\draw[gluon] (A) -- (3.2,3.9) node[lab,above right=-2pt] {$g$ \#17 emitted, $42.1\GeV$};
\draw[gluon] (A) -- (B);
\node[xlab,anchor=north] at (2.9,1.7) {$g$ \#7, $x'=2.79\times10^{-3}$\\spacelike, $Q=47.3\GeV$};
\draw[fermion] (B) -- (6.2,3.0) node[lab,above right=-2pt] {$u$ \#8 emitted, $45\GeV$};
\draw[antifermion] (B) -- (H);
\node[xlab,anchor=north east] at (5.7,0.45) {$\bar u$ \#3, $x=2.17\times10^{-3}$\\spacelike, $Q=114.6\GeV$};
\node[vertex] at (A) {}; \node[vertex] at (B) {};
\node[stepno] at ($(A)+(-0.55,-0.35)$) {3};
\node[stepno] at ($(B)+(-0.1,-0.6)$) {1};
% -z side (bottom): u(#11) -> u(#4) + g(#12)
\coordinate (C) at (3.0,-2.4);
\draw[fermion] (-0.9,-3.2) -- (C);
\node[xlab,below] at (0.3,-3.0) {$u$ \#11 from the $-z$ proton\\$x'=0.459$};
\draw[gluon] (C) -- (4.9,-3.7) node[lab,below right=-2pt] {$g$ \#12 emitted, $36.6\GeV$};
\draw[fermion] (C) -- (H);
\node[xlab,above left=-1pt] at (5.4,-1.5) {$u$ \#4, $x=0.390$\\spacelike, $Q=132.4\GeV$};
\node[vertex] at (C) {};
\node[stepno] at ($(C)+(-0.45,0.5)$) {2};
% time and evolution arrows
\draw[-{Stealth},thick] (-0.9,4.6) -- (3.0,4.6) node[midway,above,font=\footnotesize] {time: the proton resolves into partons};
\draw[-{Stealth},thick,orange!80!black] (9.6,4.6) -- (5.2,4.6)
  node[midway,above,font=\footnotesize,text=orange!80!black] {backward evolution: $53.9\to51.4\to42.2\GeV$};
\node[lab,text=black!65,text width=11cm] at (4.35,-4.7) {The emitted partons are labelled by their laboratory \pt;
  the evolution variable $\ptevol^2=(1-z)Q^2$ of the three steps is $53.9$, $51.4$ and $42.2\GeV$.};
\end{tikzpicture}
}
\caption[The initial-state branchings of our event as a diagram]{The three
initial-state branchings of our event drawn as one diagram, with the row
numbers of the record. Time runs from left to right: the $+z$ proton
provides a gluon at $x''=1.36\times10^{-2}$, which radiates a gluon and
continues as a spacelike gluon at $x'=2.79\times10^{-3}$, which splits into
the $\bar u$ at $x=2.17\times10^{-3}$ that enters the hard process and a $u$
quark that goes into the final state; the $-z$ proton provides a $u$ at
$x'=0.459$ that emits a gluon before entering at $x=0.390$. The shower
built this diagram from right to left, at $53.9$, $51.4$ and $42.2\GeV$
(the circled numbers), asking at each step where the parton it already had
came from.}
\label{fig:mother_gives_birth:isr_diagram}
\end{figure}

Branching~1 of the log, the circled~1 of
\cref{fig:mother_gives_birth:isr_diagram}, is \cref{eq:mother_gives_birth:isr}
with every quantity filled in. The $\bar u$ of the hard process, $x=2.17429\times10^{-3}$,
is found to come from a gluon with $z=0.77894$, so the gluon entered with
\begin{equation*}
  x'=\frac{x}{z}=\frac{2.17429\times10^{-3}}{0.77894}=2.79136\times10^{-3},
  \qquad
  Q^2=\frac{\ptevol^2}{1-z}=\frac{2902.735}{0.22106}=13130.8\GeV^2 ,
\end{equation*}
and the record confirms the meaning of $Q^2$: the four-momentum of the
new incoming gluon minus that of the emitted $u$ squares to
$-13130.8\GeV^2$, spacelike. The coupling is evaluated at
$\pt^2+p_{\mathrm{T}0}^2=2906.7\GeV^2$, $\alphas=0.1282$, and the PDF
ratio is $16.247/1.037=15.66$. The emitted $u$ quark has
$(p_x,p_y)=(14.62,-42.81)\GeV$, that is $\pt=45.24\GeV$ in the
laboratory against $\ptevol=53.88\GeV$, and an energy of $175.9\GeV$,
although its parent gluon carries only $18.98\GeV$ along the beam: a
spacelike branching is not a decay, and the energy comes from the whole
scattering system. The $\bbbar$ pair, which has not radiated, is boosted
to balance the $45\GeV$: the $\bar b$ goes from $113.22$ to $115.77\GeV$
and the mother from $113.22$ to $93.74\GeV$, while their invariant mass
stays at $396.10\GeV$ to the last digit. Including the new $u$ the
partonic system now has a mass of $448.8\GeV$, as it must for the larger
$x'$.

The PDF ratio is where the initial state differs most from the final. In
the collinear limit of \cref{eq:mother_gives_birth:collinear_limit} nothing
knew about the proton; here the shower is steered by the parton densities at
every step, and the ratio can be large or small. For the $\bar u$ of our
event at $x=2.17\times10^{-3}$, the chance that it came from a gluon is
governed by $x'f_g(x')/xf_{\bar u}(x)$ at the scale of the first branching,
$\pt=53.9\GeV$: the standalone tables give $1.04$ for the antiquark and
$16.2$ for the gluon at $x'=x/0.779=2.79\times10^{-3}$, a ratio of $15.7$,
multiplied by $P_{qg}(0.779)=0.33$. For the valence $u$ at $x=0.39$ the
$q\leftarrow q$ ratio at its first branching is $0.65$, because the valence
density falls with $x$, multiplied by the large $P_{qq}(0.850)=15.3$. The
small-$x$ side radiates through gluon splitting, the valence side through
gluon emission; that is the general pattern at the LHC, and it is why in
scene~2 of the film (\cref{sec:mother_gives_birth:animation}) the $+z$ beam
won the first branching more than half the time.

\Cref{fig:mother_gives_birth:isr_plot} draws the integrand of
\cref{eq:mother_gives_birth:isr} for both beams as a function of $z$ at the
starting scale, kernel times PDF ratio: on the $+z$ side the
$\bar u\leftarrow g$ channel is ten to a hundred times larger than
$\bar u\leftarrow\bar u$ over the whole range, on the $-z$ side the
$u\leftarrow u$ channel wins and $u\leftarrow g$ is suppressed by the
vanishing of the gluon density at $x'>0.4$. It also shows what the
$p_{\mathrm{T}0}$ regularisation does to the coupling and to the measure in
the last decade above the cut-off.

\begin{figure}[htbp]\centering
\includegraphics[page=5,width=\textwidth]{PS/Standalone/theory_parton_showering.pdf}
\caption[The backward-evolution integrand of the two beams, and the infrared regularisation of the initial state]{Initial-state radiation, from \file{plot_theory_PS.py}. (a) The integrand $P_{ba}(z)\,x'f_b(x')/xf_a(x)$ of \cref{eq:mother_gives_birth:isr} at $Q^2=(\pt^{\max})^2$ for the two channels of each beam of our event, with the $z$ values of branchings~1 and~2 marked: the small-$x$ antiquark evolves backwards into a gluon, the valence quark into a quark. (b) The coupling evaluated at $\pt^2+p_{\mathrm{T}0}^2$ against the final-state $\alphas(\pt^2)$, and the measure factor $\pt^2/(\pt^2+p_{\mathrm{T}0}^2)$, with $p_{\mathrm{T}0}=2\GeV$.}
\label{fig:mother_gives_birth:isr_plot}
\end{figure}

\subsection{Interleaving, the starting scale and the cut-offs}
\label{sec:mother_gives_birth:scales}

In which order do things happen in a shower? Guess before reading: does the
initial state radiate first and the final state afterwards, or do the two
$b$ quarks shower while the protons are still deciding what to send in?
Neither; there is no ``afterwards'', only a scale that runs down, and
whoever proposes the hardest emission at each round goes next, whichever
side of the hard process it belongs to. The scales involved are laid out in
\cref{fig:mother_gives_birth:scale_ladder}.
Both showers evolve in the same variable, so \PYTHIA runs them as one
sequence~\cite{Sjostrand:2004ef}: from the current scale every dipole end
and both incoming partons propose their next branching by the veto
algorithm, the hardest proposal is executed, the event is updated, and all
proposals are made afresh from that scale [\code{Shower.run}]. Nothing is
remembered from the losing proposals; that is legitimate because the veto
algorithm's survival probabilities factorise through every scale. The
ladder of our event reads
\begin{equation*}
  113.32\to53.88\to51.36\to42.16\to34.22\to32.07\to25.62\to\cdots
  \to0.647\to0.644\GeV ,
\end{equation*}
38 rungs, and it is an ordering in resolution, not a clock: the first
three rungs are initial-state branchings reconstructed backwards towards
the beams, the next two are the first emissions of the two $b$ quarks.
Watch what happens to a loser. In round~1 the $-z$ beam proposed
$12.70\GeV$ and lost to $53.88$; in round~2 it proposed afresh, below
$53.88$, and got $51.36\GeV$, which won. The $12.70\GeV$ was never a
scheduled emission. In the
full generator the multi-parton interactions of \cref{ch:uninvited_cousins}
are interleaved in the same way, which is why the \CMSSW log quotes one
``max \pt scale'' for all three.

\begin{figure}[htbp]\centering
\resizebox{\textwidth}{!}{% ---------------------------------------------------------------------------
% The scales of the shower on a logarithmic ladder: the starting scale, the
% quark-mass thresholds of the coupling, pT0 of the initial state, the two
% cut-offs and Lambda_3, with the 38 rungs of our event's ladder (-v log).
% ---------------------------------------------------------------------------
\begin{tikzpicture}[x=1cm,y=1cm,
  every node/.style={font=\footnotesize},
  lab/.style={font=\scriptsize,align=center},
  head/.style={font=\small\bfseries},
  tick/.style={thick},
  ladder/.style={-{Stealth},very thick},
  rung/.style={line width=1.6pt}]
\node[head,anchor=west] at (-0.3,4.1) {the scales of the shower, $\log_{10}(\pt/\mathrm{GeV})$: the interleaved evolution runs from right to left};
% axis: log10 from -0.6 (0.25 GeV) to 2.3 (200 GeV); 1 decade = 5.2 cm; x = (log10 v + 0.6) * 5.2
\draw[ladder] (15.6,0) -- (-0.4,0);
\foreach \l/\t in {0/{$1$},1/{$10$},2/{$100$}} {\draw[tick] ({(\l+0.6)*5.2},-0.12) -- ({(\l+0.6)*5.2},0.12); \node[lab,below] at ({(\l+0.6)*5.2},-0.15) {\t};}
\node[lab,below] at ({(-0.301+0.6)*5.2},-0.15) {$0.5$}; \draw[tick] ({(-0.301+0.6)*5.2},-0.08) -- ({(-0.301+0.6)*5.2},0.08);
\node[lab,below] at ({(0.699+0.6)*5.2},-0.15) {$5$}; \draw[tick] ({(0.699+0.6)*5.2},-0.08) -- ({(0.699+0.6)*5.2},0.08);
\node[lab,below] at ({(1.699+0.6)*5.2},-0.15) {$50$}; \draw[tick] ({(1.699+0.6)*5.2},-0.08) -- ({(1.699+0.6)*5.2},0.08);
% the 38 rungs of our event (pT_evol in GeV), FSR red/blue by family is not known here: orange = ISR, black = FSR
\foreach \v in {53.877,51.361,42.157,5.933,1.245} {\draw[rung,orange!90!black] ({(ln(\v)/ln(10)+0.6)*5.2},0.0) -- ({(ln(\v)/ln(10)+0.6)*5.2},0.75);}
\foreach \v in {34.218,32.070,25.615,22.096,13.377,10.618,5.967,5.802,4.402,2.571,2.485,2.256,1.866,1.771,1.489,1.022,0.985,0.864,0.861,0.823,0.820,0.803,0.767,0.746,0.737,0.735,0.730,0.724,0.714,0.700,0.675,0.647,0.644}
  {\draw[rung,black!70] ({(ln(\v)/ln(10)+0.6)*5.2},0.0) -- ({(ln(\v)/ln(10)+0.6)*5.2},0.5);}
\node[lab,anchor=west] at (11.0,0.95) {orange: the $5$ initial-state branchings};
\node[lab,anchor=west] at (11.0,0.6) {black: the $33$ final-state branchings};
% scales above
\draw[thick,red!70!black] ({(2.0544+0.6)*5.2},0) -- ({(2.0544+0.6)*5.2},2.9); \node[lab,above,text=red!70!black] at ({(2.0544+0.6)*5.2},2.9) {$\pt^{\max}=113.3\GeV$\\$\sqrt{\pthat^2+\mb^2}$, the hard process};
\draw[thick] ({(0.6812+0.6)*5.2},0) -- ({(0.6812+0.6)*5.2},2.0); \node[lab,above] at ({(0.6812+0.6)*5.2},2.0) {$\mb=4.8\GeV$: $n_f$ $5\to4$,\\$g\to\bbbar$ possible above};
\draw[thick] ({(0.1761+0.6)*5.2},0) -- ({(0.1761+0.6)*5.2},2.9); \node[lab,above] at ({(0.1761+0.6)*5.2},2.9) {$m_c=1.5\GeV$: $n_f$ $4\to3$};
\draw[thick,blue!60!black] ({(0.301+0.6)*5.2},0) -- ({(0.301+0.6)*5.2},1.1); \node[lab,above,text=blue!60!black,anchor=south west] at ({(0.301+0.6)*5.2},1.1) {$p_{\mathrm{T}0}=2\GeV$\\regularises ISR};
% below
\draw[thick,gray!70] ({(-0.2140+0.6)*5.2},0) -- ({(-0.2140+0.6)*5.2},-1.0); \node[lab,anchor=north west,text=gray!80!black] at ({(-0.2140+0.6)*5.2},-1.0) {$\pt^{\min,\mathrm{FSR}}=1.6\Lambda_3=0.611\GeV$};
\draw[thick,gray!70] ({(-0.3766+0.6)*5.2},0) -- ({(-0.3766+0.6)*5.2},-1.9); \node[lab,anchor=north west,text=gray!80!black] at ({(-0.3766+0.6)*5.2},-1.7) {$\pt^{\min,\mathrm{ISR}}=1.1\Lambda_3=0.420\GeV$};
\draw[thick,gray!70,dashed] ({(-0.4179+0.6)*5.2},0) -- ({(-0.4179+0.6)*5.2},-2.8); \node[lab,anchor=north west,text=gray!80!black] at ({(-0.4179+0.6)*5.2},-2.4) {$\Lambda_3=0.382\GeV$: the Landau pole of the three-flavour coupling};
\node[lab,anchor=north west,text=black!65] at (6.3,-1.75) {$\alphas$: $0.114$ at $113\GeV$, $0.14$ at $32$, $0.18$ at $10$, $0.21$ at $5$, $0.53$ at $1\GeV$, $1.56$ at the cut-off};
\node[lab,text=black!65,text width=15.6cm,anchor=north west] at (-0.3,-3.1) {The first ten rungs span three octaves ($53.9\to6.0\GeV$), the last eighteen less than one ($1.02\to0.64$): the number of branchings per octave rises as the scale falls and the coupling grows. Emissions harder than $\pt^{\max}$ are forbidden; below the cut-off nobody proposes anything and hadronisation takes over.};
\end{tikzpicture}
}
\caption[The scales of the shower on a logarithmic ladder with the 38 rungs of our event]{The scales of the shower on a logarithmic ladder, the evolution running from right to left: the starting scale \cref{eq:mother_gives_birth:ptmax}, the quark-mass thresholds of the coupling, the $p_{\mathrm{T}0}$ of the initial state, the two cut-offs \cref{eq:mother_gives_birth:ptmin} and the Landau pole $\Lambda_3$ below them, with the $38$ rungs of the seed-1 shower drawn as ticks (orange: initial state). The coupling at the marked scales is that of the \code{-v} log.}
\label{fig:mother_gives_birth:scale_ladder}
\end{figure}

The starting scale is \PYTHIA's factorisation scale of the $2\to2$ process,
\begin{equation}
  \pt^{\max}=\sqrt{\pthat^2+\tfrac12(m_3^2+m_4^2)}=\sqrt{113.22^2+4.8^2}=113.32\GeV ,
  \label{eq:mother_gives_birth:ptmax}
\end{equation}
printed in the log as ``Max pT scale for MPI = 1.133e+02, ISR = 1.133e+02,
FSR = 1.133e+02''. Emissions harder than the hard process itself are not
allowed; this is the crude but standard answer to the question of
\cref{ch:whose_child}. The cut-offs at the bottom are
\begin{equation}
\begin{aligned}
  \pt^{\min,\mathrm{FSR}}&=\max\bigl(0.5\GeV,\,1.6\,\Lambda_3\bigr)=0.611\GeV,\\
  \pt^{\min,\mathrm{ISR}}&=\max\bigl(0.2\GeV,\,1.1\,\Lambda_3\bigr)=0.420\GeV,
\end{aligned}
  \label{eq:mother_gives_birth:ptmin}
\end{equation}
where $\Lambda_3=0.382\GeV$ is the three-flavour $\Lambda$ of the shower's
coupling. The first is the origin of the warning
``\code{SimpleTimeShower::init: pTmin too low, raised to 0.611}'' at the top
of every CP5 log: the tune asks for $0.5\GeV$, and \PYTHIA refuses to go
closer than a factor $1.6$ to the Landau pole. Below the cut-off the
shower stops and hadronisation (\cref{ch:dressing_children}) takes over.
The cut-offs apply to the evolution variable, not to what ends up in the
record: the last branching of our event, number~38, is a $\bar d$ emitting
a gluon at $\ptevol=0.644\GeV$ with $z=0.743$ and a pair mass of
$1.47\GeV$, and the gluon has a laboratory \pt of $0.39\GeV$, below the
cut-off, which is perfectly possible because the two are different
quantities. Nor does the last emission sit at the cut-off: in round~39
nobody proposed anything above it, and that, not reaching $0.611\GeV$, is
what stops the shower. The flavour thresholds of the running coupling,
$m_c=1.5$ and $m_b=4.8\GeV$, also sit inside the evolution range: above
$4.8\GeV$ a gluon may split into $\bbbar$, and in our event one does, at
branching~10, with $\ptevol=5.967\GeV$ and $z=0.350$. Its pair mass,
$\sqrt{35.60/(0.350\times0.650)}=12.51\GeV$, clears the threshold
$2\mb=9.6\GeV$; the two new quarks have $27.2$ and $40.1\GeV$ in the
laboratory and are born $\dR=0.23$ apart. The gluon was the one the
$\bar b$ had radiated at branching~4, so the event now contains two $b$
and two $\bar b$ quarks, and ``a $b$ quark in the final state'' no longer
says where it came from.

\Cref{fig:mother_gives_birth:interleaving_plot} draws the ladder against the
round number together with the thresholds and cut-offs, and counts the
competitors of every round from the \code{-v} log: seven dipole ends and
two beams in round~1, growing by two with every final-state branching to
sixty-nine in round~39, where all of them are silent.

\begin{figure}[htbp]\centering
\includegraphics[page=6,width=\textwidth]{PS/Standalone/theory_parton_showering.pdf}
\caption[The ladder of our event and the number of competitors in every round]{The interleaved evolution, from \file{plot_theory_PS.py}. (a) The $38$ rungs of the seed-1 shower against the round number, coloured by kind, with the starting scale, the quark-mass thresholds, $p_{\mathrm{T}0}$ and the two cut-offs as horizontal lines. (b) The competitors of every round parsed from the \code{-v} log: dipole ends and beams that proposed an emission, and those that evolved below the cut-off and stayed silent.}
\label{fig:mother_gives_birth:interleaving_plot}
\end{figure}

\subsection{The CP5 shower settings and the coupling}
\label{sec:mother_gives_birth:cp5}

How many strong couplings does a CMS event have? One, you would hope; in
fact \PYTHIA keeps four, one each for the matrix element, the two showers
and the multi-parton interactions, and the tune decides whether they
agree (\cref{fig:mother_gives_birth:cp5_sketch}).
The CP5 tune~\cite{CMS:2019csb} touches the showers in one place: it sets
\code{TimeShower:alphaSvalue} and \code{SpaceShower:alphaSvalue} to
$0.118$ with second-order running (\code{alphaSorder = 2}), where \PYTHIA's
defaults are $0.1365$ at first order. These are the lines listed under
``changes only'' in the settings table of the \CMSSW log. The value is the
same as for the matrix element and the multi-parton interactions, and the
running is \PYTHIA's own: second order with flavour thresholds at
$m_c=1.5$, $m_b=4.8$ and $m_t=171\GeV$, matched so that the coupling is
continuous, which fixes
\begin{equation}
  \Lambda_5=0.2262\GeV,\qquad\Lambda_4=0.3282\GeV,\qquad\Lambda_3=0.3820\GeV .
  \label{eq:mother_gives_birth:lambdas}
\end{equation}
The standalone script reproduces these analytically to six digits
[\code{AlphaStrong}]; $\alphas$ is $0.114$ at $\pt^{\max}$, $0.14$ at
$32\GeV$ where the mother first radiates, $0.18$ at $10\GeV$, $0.21$ at
$5\GeV$ and $0.53$ at $1\GeV$. The last number is the reason for the
cut-off, and for the excess of soft emissions in
\cref{tab:mother_gives_birth:numbers}: the last octave above $\pt^{\min}$
is where the coupling is largest and where any difference in the treatment
of the infrared region shows.

\begin{figure}[htbp]\centering
\resizebox{\textwidth}{!}{% ---------------------------------------------------------------------------
% Where the couplings and the PDFs of a CMS CP5 event come from: one
% alpha_s(mZ) = 0.118 at second order for the matrix element, the two showers
% and the MPI, one PDF set for the hard process and the backward evolution.
% ---------------------------------------------------------------------------
\begin{tikzpicture}[x=1cm,y=1cm,
  every node/.style={font=\footnotesize},
  stage/.style={draw,rounded corners=3pt,thick,fill=orange!12,align=center,minimum height=13mm,text width=3.1cm,font=\scriptsize},
  src/.style={draw,rounded corners=3pt,thick,fill=blue!8,align=center,minimum height=10mm,text width=4.6cm,font=\scriptsize},
  flow/.style={-{Stealth[length=2.5mm,width=2mm]},thick,gray!60},
  lab/.style={font=\scriptsize,align=center},
  head/.style={font=\small\bfseries}]
\node[head,anchor=west] at (-0.3,3.3) {one coupling, one PDF set, four places where they enter a CP5 event};
\node[src] (A) at (2.5,2.0) {\textbf{\code{alphaSvalue = 0.118}, \code{alphaSorder = 2}}\\$\Lambda_5=0.226$, $\Lambda_4=0.328$, $\Lambda_3=0.382\GeV$};
\node[src] (P) at (10.5,2.0) {\textbf{\code{PDF:pSet}: NNPDF3.1 NNLO}, member 0\\$xf(x,Q^2)$ for all flavours and the gluon};
\node[stage] (ME) at (0,0) {\textbf{matrix element}\\\code{SigmaProcess:}\\$\alphas(\pthat^2+\mb^2)=0.1143$\\PDFs at $\muF=113\GeV$};
\node[stage] (ISR) at (3.6,0) {\textbf{initial-state shower}\\\code{SpaceShower:}\\$\alphas(\pt^2+p_{\mathrm T0}^2)$\\PDF ratios at $\pt^2$};
\node[stage] (FSR) at (7.2,0) {\textbf{final-state shower}\\\code{TimeShower:}\\$\alphas(\ptevol^2)$, $0.14\to1.56$\\no PDF};
\node[stage,text width=3.7cm] (MPI) at (11.1,0) {\textbf{multi-parton interactions}\\\code{MultipartonInteractions:}\\$\alphas(\pt^2+p_{\mathrm T0}^2)$, PDFs\\(\cref{ch:uninvited_cousins})};
\foreach \s in {ME,ISR,FSR,MPI} {\draw[flow] (A.south) -- (\s.north);}
\foreach \s in {ME,ISR,MPI} {\draw[flow] (P.south) -- (\s.north);}
\node[lab,text=black!65,text width=14.4cm,anchor=north west] at (-1.6,-1.0) {\PYTHIA's defaults would be $0.1365$ at first order for both showers; CP5 ties all four to the same $0.118$ at second order, the ``changes only'' lines of the \CMSSW settings table. The shower $\alphas$ is a tune parameter, not the $\overline{\mathrm{MS}}$ world average; the standalone script reproduces the running and the thresholds analytically (\code{AlphaStrong}) and embeds the PDF grid (\code{TablePDF}).};
\end{tikzpicture}
}
\caption[One coupling and one PDF set, four places where they enter a CP5 event]{Where the coupling and the parton densities of a CP5 event come from: a single $\alphas(m_Z)=0.118$ at second order feeds the matrix element, the two showers and the multi-parton interactions, each evaluating it at its own scale, and a single PDF set serves the hard process, the backward evolution and the multi-parton interactions.}
\label{fig:mother_gives_birth:cp5_sketch}
\end{figure}

The PDFs of the backward evolution are the same NNPDF3.1 set as at the
first stop, now needed for all flavours and the gluon from the cut-off up to
$\pt^{\max}$, and for all $x$ from the smallest incoming fraction up to one.
The set's grid starts at $Q=1.65\GeV$; below that the values are frozen, as
\PYTHIA does. Everything else, the rapidity-ordering option of the initial
state, matrix-element corrections, QED radiation, spin correlations, is
either \PYTHIA default or left out of the standalone version; the list is in
\cref{sec:mother_gives_birth:caveats}.

\Cref{fig:mother_gives_birth:coupling_plot} draws the coupling the script
uses, with its thresholds and the values quoted above, next to the
first-order form that the veto algorithm overestimates it with; their ratio
is the coupling factor of \cref{eq:mother_gives_birth:veto_weight}, between
$0.73$ and $0.87$ over the shower and above one only in the last octave.

\begin{figure}[htbp]\centering
\includegraphics[page=7,width=\textwidth]{PS/Standalone/theory_parton_showering.pdf}
\caption[The CP5 shower coupling with its thresholds, and the first-order overestimate]{The shower coupling, from \file{plot_theory_PS.py}. (a) \PYTHIA's second-order $\alphas(Q^2)$ with $\alphas(m_Z)=0.118$ and the flavour thresholds at $m_c$ and $m_b$, as reproduced by \code{AlphaStrong}, with the values of section~[1] of the \code{-v} log marked, and the first-order form \cref{eq:mother_gives_birth:alphas_1loop} with the same $\Lambda_{n_f}$. (b) Their ratio, the coupling factor of the acceptance weight: the price of the overestimate, and the place where it fails.}
\label{fig:mother_gives_birth:coupling_plot}
\end{figure}

\subsection{From formula to event: the shower step by step}
\label{sec:mother_gives_birth:step_by_step}

We now have everything needed to turn the hard-process record into a
showered event, and as at the previous stop I want you to see the machine,
not just the output. The steps below are those of the standalone script
you ran in \cref{sec:mother_gives_birth:standalone_code}.

\begin{figure}[htbp]\centering
\resizebox{\textwidth}{!}{% ---------------------------------------------------------------------------
% The six steps of the standalone shower as a flow chart, with the numbers of
% the seed-1 event from the -v log (round 1 proposals, winner, record size).
% ---------------------------------------------------------------------------
\begin{tikzpicture}[x=1cm,y=1cm,
  every node/.style={font=\footnotesize},
  stepbox/.style={draw,rounded corners=3pt,thick,fill=orange!12,align=left,inner sep=5pt,minimum height=11mm,text width=5.4cm},
  numbox/.style={draw,rounded corners=3pt,thin,fill=blue!6,align=left,inner sep=4pt,text width=8.4cm,font=\scriptsize},
  stepno/.style={draw,circle,inner sep=1pt,font=\scriptsize\bfseries,fill=orange!40,minimum size=5mm},
  flow/.style={-{Stealth[length=2.5mm,width=2mm]},thick},
  loop/.style={-{Stealth[length=2.5mm,width=2mm]},thick,red!70!black},
  fn/.style={font=\scriptsize\ttfamily,text=black!65},
  head/.style={font=\small\bfseries}]
\node[head,anchor=west] at (-0.6,1.5) {the recipe};
\node[head,anchor=west] at (6.5,1.5) {our event (seed 1), from the \code{-v} log};
\node[stepbox] (S1) at (2.4,0) {\textbf{1\; read the hard process, find the dipoles}\\ incoming and outgoing partons with colour tags; every end of a colour line is a radiator};
\node[stepbox,below=8mm of S1] (S2) {\textbf{2\; set the scales}\\ $\pt^{\max}$ from the hard process, $\pt^{\min}$ from $\Lambda_3$, $\Lambda_{3,4,5}$ from $\alphas(m_Z)$};
\node[stepbox,below=8mm of S2] (S3) {\textbf{3\; let everybody propose}\\ veto algorithm for each dipole end (FSR) and each beam (ISR, with the PDF ratio)};
\node[stepbox,below=8mm of S3] (S4) {\textbf{4\; the hardest wins}\\ execute it: new rows, status codes, recoil, colour tags handed on};
\node[stepbox,below=8mm of S4] (S5) {\textbf{5\; repeat from the winning scale}\\ until nobody proposes anything above the cut-off};
\node[stepbox,below=8mm of S5] (S6) {\textbf{6\; write the record, summarise}\\ final state = rows with positive status; follow each $b$ to its last copy; family within $\dR<0.4$};
\foreach \a/\b in {S1/S2,S2/S3,S3/S4,S4/S5,S5/S6} {\draw[flow] (\a) -- (\b);}
\draw[loop] (S5.west) -- ++(-0.6,0) |- (S3.west) node[pos=0.25,left,font=\scriptsize,text=red!70!black,align=center] {38\\rounds};
\foreach \n/\s in {1/S1,2/S2,3/S3,4/S4,5/S5,6/S6} {\node[stepno] at ($(\s.west)+(-0.35,0)$) {\n};}
\node[fn,anchor=south east] at (S1.north east) {read\_event, dipole\_ends};
\node[fn,anchor=south east] at (S3.north east) {pT2next\_fsr, pT2next\_isr};
\node[fn,anchor=south east] at (S4.north east) {branch\_fsr, branch\_isr};
\node[fn,anchor=south east] at (S5.north east) {Shower.run};
\node[fn,anchor=south east] at (S6.north east) {family, summarise};
\node[numbox,anchor=west] (N1) at (6.5,0 |- S1) {$\bar u$ \#3 ($x_A=2.17\times10^{-3}$) and $u$ \#4 ($x_B=0.390$) in; $\bar b$ \#5 and $b$ \#6 out, $\pt=113.22\GeV$ each\\ dipoles: $\bar b$--$\bar u$ on line 101, $b$--$u$ on line 102, both $m_{\mathrm{dip}}=377.9\GeV$; four ends and two beams can radiate};
\node[numbox,anchor=west] (N2) at (6.5,0 |- S2) {$\pt^{\max}=\sqrt{113.22^2+4.8^2}=113.32\GeV$;\quad $\pt^{\min}$: $0.611$ (FSR), $0.420\GeV$ (ISR);\\ $\Lambda_5=0.2262$, $\Lambda_4=0.3282$, $\Lambda_3=0.3820\GeV$; $7.5$ octaves to fill};
\node[numbox,anchor=west] (N3) at (6.5,0 |- S3) {round 1: $\bar b$ proposes $2.159\GeV$ ($z=0.953$), $b$ $5.241$ ($0.982$), $+z$ beam $53.877$ ($\bar u\leftarrow g$, $0.779$), $-z$ beam $12.700$ ($u\leftarrow u$, $0.797$)};
\node[numbox,anchor=west] (N4) at (6.5,0 |- S4) {ISR $\bar u\leftarrow g$ at $53.877\GeV$: rows 7 ($g$, $x'=2.79\times10^{-3}$, status $-41$) and 8 ($u$, $45.2\GeV$, status 43); the $\bbbar$ pair recoils, $113.22\to115.77$ and $93.74\GeV$};
\node[numbox,anchor=west] (N5) at (6.5,0 |- S5) {$113.32\to53.88\to51.36\to42.16\to34.22\to32.07\to\cdots\to0.647\to0.644\GeV$; in round 39 all $69$ competitors are silent};
\node[numbox,anchor=west] (N6) at (6.5,0 |- S6) {158 rows, 40 final ($27$ gluons, $13$ quarks); $x_A=0.0171$, $x_B=0.470$;\\ $b$: $113.22\to78.12\GeV$, family $100.45\GeV$, $m=13.52$, $n=3$;\quad $\bar b$: $73.18$, family $92.13\GeV$};
\foreach \s/\n in {S1/N1,S2/N2,S3/N3,S4/N4,S5/N5,S6/N6} {\draw[flow,gray!60] (\s.east) -- (\n.west);}
\end{tikzpicture}
}
\caption[The six steps of the shower, with the numbers of our event]{The six steps of the standalone shower as a flow chart, with the function names of the script. Right: the same steps for the seed-1 shower of our event, every number from the \code{-v} log, from the two dipoles of the hard process to the $158$ rows of the record.}
\label{fig:mother_gives_birth:shower_flow}
\end{figure}

\paragraph{Step 1: read the hard process and find the dipoles.} From
\file{event_ME.lhe} the script takes the two incoming and two outgoing
partons with their colour tags [\code{read\_event}]. The two incoming
partons are as unequal as partons can be: the $\bar u$ carries
$14.785\GeV$, $x_A=2.17\times10^{-3}$ of its proton, the $u$ carries
$2652.9\GeV$, $x_B=0.390$, so that their laboratory energy is
$2667.7\GeV$ but $\sqrt{\shat}=\sqrt{x_Ax_Bs}=396.1\GeV$, and both $b$
quarks fly towards $-z$ with $\pt=113.221\GeV$ each. Colour~102 links the $b$
to the incoming $u$, anticolour~101 the $\bar b$ to the incoming $\bar u$:
two dipoles, four ends, and two beams that can radiate
[\code{dipole\_ends}]. Both dipoles have a mass of $377.9\GeV$ and could
host emissions up to $188.9\GeV$; the starting scale is
$\min(188.9,113.3)=113.3\GeV$, because phase space does not override the
hard process.

\paragraph{Step 2: set the scale.} $\pt^{\max}=113.32\GeV$ from
\cref{eq:mother_gives_birth:ptmax}; $\pt^{\min}$ from
\cref{eq:mother_gives_birth:ptmin}; $\Lambda_{3,4,5}$ from
\cref{eq:mother_gives_birth:lambdas}.

\paragraph{Step 3: let everybody propose.} For each dipole end, the veto
algorithm of \cref{sec:mother_gives_birth:veto} with the FSR kernels and
\cref{eq:mother_gives_birth:pt_evol_fsr} [\code{pT2next\_fsr}]; for each
beam, the same with the ISR kernels and the PDF ratio of
\cref{eq:mother_gives_birth:isr} [\code{pT2next\_isr}]. Each returns the
scale of its next branching, or nothing if it evolved below the cut-off.
Round~1 of our event:
\begin{center}\small
\begin{tabular}{llr}\toprule
competitor & proposed branching & $\ptevol$ [GeV] \\ \midrule
$\bar b$ against the incoming $\bar u$ & $\bar b\to\bar b g$, $z=0.953$ & $2.159$ \\
$b$ against the incoming $u$ & $b\to bg$, $z=0.982$ & $5.241$ \\
$+z$ beam, $\bar u$ at $x=2.17\times10^{-3}$ & $\bar u\leftarrow g$, $z=0.779$ & $\mathbf{53.877}$ \\
$-z$ beam, $u$ at $x=0.390$ & $u\leftarrow u$, $z=0.797$ & $12.700$ \\
\bottomrule
\end{tabular}
\end{center}

\paragraph{Step 4: the hardest wins.} Execute that branching: new partons
are appended to the record with \PYTHIA's status codes, the parents are
marked as replaced (negative status) and the recoil is distributed
[\code{branch\_fsr}, \code{branch\_isr}]. Colour tags are handed on so that
the new dipoles can be found.

\paragraph{Step 5: repeat from the winning scale.} Back to step~3 with
$\pt^{\max}$ replaced by the scale just reached, until nobody proposes
anything above the cut-off [\code{Shower.run}].
\Cref{fig:mother_gives_birth:competition_plot} is steps~3 to~5 for the
whole event: every proposal of every round from the \code{-v} log, the
winner of each round joined by the black line, and the scale the next
round starts from. Read round~1 on the left, the four numbers of the table
above; read how the losers of a round reappear below the winner's scale in
the next; and watch the proposals thin out towards the cut-off until, in
round~39, there are none.

\begin{figure}[htbp]\centering
\includegraphics[page=8,width=\textwidth]{PS/Standalone/theory_parton_showering.pdf}
\caption[Every proposal of every round of our event, and the winner]{Everybody proposes, the hardest wins: every proposal of every round of the seed-1 shower, parsed from the \code{-v} log by \file{plot_theory_PS.py}. Green dots are proposals of dipole ends, orange squares those of the two beams, the open circles joined by the black line the winners, and the red dashed line the scale each round starts from, which is the previous winner. Losers are thrown away and propose afresh below the winner's scale.}
\label{fig:mother_gives_birth:competition_plot}
\end{figure}

\paragraph{Step 6: write the record and summarise.} The final state is every
row with positive status. The script follows each hard $b$ quark down the
record to its final copy, collects the final partons within $\dR<0.4$ of it
and prints the family's \pt, energy, mass and multiplicity
[\code{family}, \code{summarise}]. The record of our event has $158$
rows, of which $40$ have positive status: $2+5+33$, since every executed
branching adds one parton. They are $27$ gluons and $13$ quarks, and you
can count the gluons from the branching list itself,
$4$ (emitted by ISR) $+\,9$ ($q\to qg$) $+\,19$ ($g\to gg$, one net each)
$-\,5$ ($g\to q\bar q$) $=27$, and the quarks as
$2+1+2\times5=13$. The two incoming partons now carry
$x_A=0.01714$ and $x_B=0.46999$, that is $116.55$ and $3195.91\GeV$, which
is the final-state energy sum of $3312.46\GeV$; the rest of the two
protons is not in the record.

For our event the sequence is 38 branchings long, and its first five are
worth reading in order, because they set everything that follows:
\begin{enumerate}[nosep]
  \item ISR on the $+z$ side at $\pt=53.9\GeV$: the $\bar u$ is resolved as
    coming from a gluon, $g\to\bar u+u$, and a $u$ quark of $54\GeV$ enters
    the final state; the incoming fraction rises from $2.17\times10^{-3}$ to
    $2.79\times10^{-3}$.
  \item ISR on the $-z$ side at $51.4\GeV$: the valence $u$ came from a $u$
    that emitted a gluon; $x$ rises from $0.390$ to $0.459$.
  \item ISR on the $+z$ side again at $42.2\GeV$: the gluon came from a
    gluon, $x\to1.36\times10^{-2}$.
  \item FSR at $34.2\GeV$: the $\bar b$ emits its first gluon, $z=0.594$.
  \item FSR at $32.1\GeV$: the mother emits her first gluon, $z=0.601$, at a
    virtuality of $65.7\GeV$.
\end{enumerate}
Three initial-state branchings, each harder than anything the mother does,
came first. That is not an accident of the seed: scene~2 of the film
(\cref{sec:mother_gives_birth:animation}) showed that the $+z$ beam, at small
$x$ where the gluon density is huge, wins the first branching in $53\%$ of
the trials. Then follow 33 more, softer and softer: gluons splitting into
gluons, a gluon splitting into $s\bar s$ at $6\GeV$, another into
$\bbbar$ at $6\GeV$, and from $2.6\GeV$ down a long tail of emissions
within a factor of four of the cut-off. The full list is printed by the
script, and \cref{sec:mother_gives_birth:our_jet} reads the record.

\section{The heavy mother: the dead cone}
\label{sec:mother_gives_birth:dead_cone}

Our mother is not massless, and this is where it shows. Go back to
\cref{eq:mother_gives_birth:propagator} and keep $\beta_b<1$: the
denominator $1-\beta_b\cos\theta$ no longer vanishes at $\theta=0$, but
levels off at $1-\beta_b\simeq\mb^2/2E_b^2$. Collinear radiation is
suppressed inside a cone of half-angle
\begin{equation}
  \theta_0=\frac{\mb}{E_b},
  \label{eq:mother_gives_birth:theta0}
\end{equation}
the \emph{dead cone}. For the mother at her birth, $E_b=252\GeV$ in the
laboratory and $\theta_0=0.019$; in the $(\eta,\phi)$ plane, where angles
near the transverse direction scale with $\pt$ rather than $E$, the relevant
number is $\mb/\pt=0.042$, still an order of magnitude below the family
cone. \Cref{fig:mother_gives_birth:dead_cone} draws the cone and the
angular distribution behind it: for a soft gluon the massless
$\mathrm{d}\theta^2/\theta^2$ of \cref{eq:mother_gives_birth:collinear_limit}
becomes
\begin{equation}
  \frac{\mathrm{d}\theta^2}{\theta^2}\;\longrightarrow\;
  \frac{\theta^2\,\mathrm{d}\theta^2}{(\theta^2+\theta_0^2)^2},
  \label{eq:mother_gives_birth:dead_cone_angle}
\end{equation}
which is the massless form for $\theta\gg\theta_0$, a quarter of it at
$\theta=\theta_0$, and $\theta^2\,\mathrm{d}\theta^2/\theta_0^4$, vanishing,
inside the cone. Check the first limit yourself: drop $\theta_0$ in the
denominator and the $\theta^2$ cancel. The radiation is not forbidden
inside the cone, it is starved of its collinear enhancement, which is why
the name is a little too dramatic.

\begin{figure}[htbp]\centering
\resizebox{\textwidth}{!}{% ---------------------------------------------------------------------------
% The dead cone of the heavy mother. Left: geometry (the cone is drawn far
% wider than theta_0 = m_b/E_b = 0.019). Right: the soft-gluon emission
% density per unit ln(theta) relative to a massless quark,
% [theta^2/(theta^2+theta_0^2)]^2 (Dokshitzer, Khoze, Troyan).
% ---------------------------------------------------------------------------
\begin{tikzpicture}[x=1cm,y=1cm,
  every node/.style={font=\footnotesize},
  gluon/.style={thick,decorate,decoration={coil,aspect=0,segment length=4pt,amplitude=1.8pt}}]
% ---------------- left: the cone ------------------------------------------------------------
\begin{scope}
\coordinate (O) at (0,0);
\fill[red!14] (O) -- (14:6.3) arc[start angle=14,end angle=-14,radius=6.3] -- cycle;
\draw[red!70!black,dashed] (O) -- (14:6.3);
\draw[red!70!black,dashed] (O) -- (-14:6.3);
\draw[very thick,-{Stealth}] (O) -- (6.6,0) node[right,align=left] {$b$\\$E_b=252\GeV$};
\node[red!70!black,align=center] at (4.6,-0.55) {dead cone $\theta<\theta_0=\mb/E_b$\\(not to scale: $\theta_0=0.019$ here)};
\draw[gluon] (O) -- (30:4.6) node[right] {$g$};
\draw[gluon] (O) -- (55:3.6) node[above right] {$g$};
\draw[gluon] (O) -- (-40:3.8) node[right] {$g$};
\draw[gluon,gray!60] (O) -- (6:3.6) node[right,text=gray!80!black,font=\scriptsize] {suppressed by $\bigl(\theta^2/(\theta^2+\theta_0^2)\bigr)^2$};
\draw[-{Stealth}] (3.1,0) arc[start angle=0,end angle=30,radius=3.1] node[midway,right] {$\theta$};
\draw[-{Stealth},red!70!black] (1.4,0) arc[start angle=0,end angle=14,radius=1.4] node[pos=1.0,above right=-2pt,font=\scriptsize] {$\theta_0$};
\node[circle,fill=black,inner sep=1.5pt] at (O) {};
\node[below left,align=right] at (O) {the mother leaves\\the hard process};
\end{scope}
% ---------------- right: the emission density --------------------------------------------
\begin{scope}[xshift=9.4cm,yshift=-1.8cm]
\draw[gray!50] (0,3) -- (5,3);
\draw[gray!50] (0,0.75) -- (2.5,0.75);
\draw[-{Stealth}] (0,0) -- (5.6,0) node[below left=2pt and -2pt] {$\theta/\theta_0$};
\draw[-{Stealth}] (0,0) -- (0,3.7) node[above,align=center,font=\scriptsize] {emission density per $\mathrm{d}\ln\theta$,\\relative to a massless quark};
\foreach \x/\l in {0/0.1,2.5/1,5/10} {\draw (\x,0) -- (\x,-0.1) node[below] {\l};}
\foreach \y/\l in {0.75/$\tfrac14$,3/1} {\draw (0,\y) -- (-0.1,\y) node[left] {\l};}
\draw[thick,dashed,gray!80!black] (0,3) -- (5,3) node[above left,font=\scriptsize,text=gray!80!black] {massless quark: $\mathrm{d}\theta^2/\theta^2$, flat};
\draw[very thick,red!80!black] plot[domain=-1:1,samples=80,variable=\L]
   ({2.5*(\L+1)},{3*(pow(10,2*\L)/(1+pow(10,2*\L)))^2});
\draw[dotted,red!70!black] (2.5,0) -- (2.5,0.75);
\node[red!80!black,font=\scriptsize,align=left,anchor=west] at (2.7,0.45) {heavy quark: $\dfrac{\theta^2\,\mathrm{d}\theta^2}{(\theta^2+\theta_0^2)^2}$};
\node[font=\scriptsize,align=left,anchor=north west,text=black!70] at (0.1,2.75) {at $\theta=\theta_0$ one quarter\\of the massless rate,\\at $\theta=\theta_0/3$ one percent};
\end{scope}
\end{tikzpicture}
}
\caption[The dead cone]{Left: the mother leaves the hard process and
radiates gluons at angles $\theta$ to her direction; inside the cone
$\theta<\theta_0=\mb/E_b$ the radiation is suppressed (the cone is drawn far
wider than the $0.019$ of our mother). Right: the soft-gluon emission
density per unit $\ln\theta$ of \cref{eq:mother_gives_birth:dead_cone_angle}
relative to a massless quark, for which it is flat.}
\label{fig:mother_gives_birth:dead_cone}
\end{figure}

In the shower the effect enters through the quasi-collinear splitting
function
\begin{equation}
  P_{qq}(z,\pt^2)=C_F\left[\frac{1+z^2}{1-z}
  -\frac{2\mb^2\,z(1-z)}{\ptevol^2}\right],
  \label{eq:mother_gives_birth:dead_cone_kernel}
\end{equation}
which the script applies as a correction weight in step~4 of the veto
algorithm and which is zero or negative, that is forbidden, when
$\ptevol^2\le2\mb^2z(1-z)^2/(1+z^2)$, roughly
$\ptevol\lesssim\mb(1-z)$. Translate that back through
$\ptevol\simeq(1-z)E\theta$ and you recover
$\theta\lesssim\mb/E$.

How much does it matter for our event? The mother's two own emissions in the
record had $(\ptevol,z)=(32.1\GeV,0.60)$ and $(4.4\GeV,0.86)$; the mass
term reduces the kernel by $0.3\%$ and $2\%$. The softest $b$-quark emission
of the whole event, branching~17 of the log, in which the $\bar b$ born from
the gluon splitting of branching~10 radiates at $(1.9\GeV,0.80)$, is
reduced by $25\%$. The hard emissions do not notice the mass at all; only
the softest ones, near the cut-off, are visibly suppressed. Averaged over many
showers the dead cone removes a few percent of the soft, collinear radiation
of a $100\GeV$ $b$ quark, and this is one of the reasons $b$ jets are
slightly narrower and carry a slightly larger fraction of their energy in
the leading hadron than light-quark jets. The effect was observed directly
by ALICE in 2022 for charm quarks by reclustering jets and comparing the
angular distribution of the splittings to that of light jets
\cite{ALICE:2021aqk}; for the top quark, at $m_t/E_t\sim0.4$ for a
$400\GeV$ top, it is not a correction but the dominant feature, and it
returns in \cref{ch:three_families}.
\Cref{fig:mother_gives_birth:dead_cone_plot} shows the mass correction of
\cref{eq:mother_gives_birth:dead_cone_kernel} as a function of $\ptevol$
and $z$ with the mother's two emissions on it, and the angular density of
\cref{eq:mother_gives_birth:dead_cone_angle} with the actual opening angles
of those emissions: both are tens of dead-cone angles away from her
direction.

\begin{figure}[htbp]\centering
\includegraphics[page=9,width=\textwidth]{PS/Standalone/theory_parton_showering.pdf}
\caption[The mass correction of the $b\to bg$ kernel and the dead-cone angular density, with the mother's emissions]{The dead cone in numbers, from \file{plot_theory_PS.py}. (a) The ratio of the quasi-collinear kernel \cref{eq:mother_gives_birth:dead_cone_kernel} to the massless one as a function of $\ptevol$ for four values of $z$, with the mother's two emissions of the seed-1 shower marked and their reductions quoted. (b) The soft-gluon density per unit $\ln\theta$ of \cref{eq:mother_gives_birth:dead_cone_angle} relative to a massless quark, against $\theta/\theta_0$, with the opening angles of the same two emissions computed from the record.}
\label{fig:mother_gives_birth:dead_cone_plot}
\end{figure}

There is a second, less obvious consequence of the mass. The virtuality
window available to a heavy quark is smaller: in
\cref{eq:mother_gives_birth:pt_evol_fsr} the branching needs
$m_a^2>\mb^2$, and the kinematic limit on the dipole,
$(m_{\mathrm{dip}}-m_{\mathrm{rec}})^2-\mb^2$, is reduced. Both are handled
exactly by the script; both are small at $113\GeV$ and would not be at
$10\GeV$.

\section{Children who run away: out-of-cone radiation and the first energy loss}
\label{sec:mother_gives_birth:runaways}

The ledger asks how much of the mother's $113.22\GeV$ is still ``hers''
after this stage, and the record answers in three layers.

\paragraph{The quark herself.} Follow the $b$ down the record, copy by copy:
$113.2\GeV$ at birth; $93.7$, $91.7$ and $112.5\GeV$ after the recoils of
the three hard initial-state branchings, before she has radiated anything;
$99.9\GeV$ after her own $32\GeV$ emission; $80.9\GeV$ after the $4.4\GeV$
one; $78.1\GeV$ at the end. Two lessons. The initial-state recoil moved her
by $20\GeV$ in both directions before her own shower began, which is why
``the \pt after ISR'' is not a monotonic quantity and why the mother's
family and the initial-state partons can only be separated by their
ancestry, never by their kinematics. And her second emission, with an
evolution variable of only $4.4\GeV$, cost her $19\GeV$ of laboratory \pt:
$\ptevol$ is a transverse momentum in the dipole frame with
respect to the dipole axis, not the \pt of the gluon in the detector (the
gluon in question has $20\GeV$ in the laboratory). Do not read a shower
history as a list of jet constituents.

\begin{figure}[htbp]\centering
\resizebox{\textwidth}{!}{% ---------------------------------------------------------------------------
% (a) The mother's pT copy by copy through the shower, with the cause of each
% change (-v log: ISR recoils of branchings 1-3 and 11, her own emissions 5
% and 13, the recoils of branchings 20 and 27); (b) the family cone of dR < 0.4 around her
% final copy with the children inside and the runaway outside (record rows).
% ---------------------------------------------------------------------------
\begin{tikzpicture}[x=1cm,y=1cm,
  every node/.style={font=\footnotesize},
  copy/.style={draw,circle,thick,fill=red!15,minimum size=9mm,inner sep=0pt,font=\scriptsize,align=center},
  cause/.style={font=\scriptsize,align=center,text=black!65},
  flow/.style={-{Stealth[length=2.5mm,width=2mm]},thick},
  isr/.style={flow,orange!90!black},
  own/.style={flow,red!80!black},
  lab/.style={font=\scriptsize,align=center},
  head/.style={font=\small\bfseries}]
% ---------------- (a) the chain of copies -------------------------------------------------------
\node[head,anchor=west] at (-0.5,2.3) {(a) the mother's \pt, copy by copy (record rows), and what moved it};
\node[copy] (c6) at (0,0) {\#6\\113.2};
\node[copy] (c10) at (2.4,0) {\#10\\93.7};
\node[copy] (c15) at (4.8,0) {\#15\\91.7};
\node[copy] (c21) at (7.2,0) {\#21\\112.5};
\node[copy] (c25) at (9.6,0) {\#25\\99.9};
\node[copy] (c46) at (12.0,0) {\#46\\99.5};
\node[copy] (c60) at (14.4,0) {\#60\\80.9};
\node[copy] (c124) at (16.8,0) {\#124\\78.1};
\draw[isr] (c6) -- (c10) node[midway,above=4pt,cause] {ISR 1\\$53.9\GeV$\\recoil};
\draw[isr] (c10) -- (c15) node[midway,below=4pt,cause] {ISR 2\\$51.4\GeV$\\recoil};
\draw[isr] (c15) -- (c21) node[midway,above=4pt,cause] {ISR 3\\$42.2\GeV$\\recoil};
\draw[own] (c21) -- (c25) node[midway,below=4pt,cause,text=red!80!black] {her 1st emission\\$\ptevol=32.1$, $z=0.60$\\gluon $38.2\GeV$ (lab)};
\draw[isr] (c25) -- (c46) node[midway,above=4pt,cause] {ISR 11\\$5.9\GeV$\\recoil};
\draw[own] (c46) -- (c60) node[midway,below=4pt,cause,text=red!80!black] {her 2nd emission\\$\ptevol=4.4$, $z=0.86$\\gluon $20.7\GeV$ (lab)};
\draw[flow,black!60] (c60) -- (c124) node[midway,above=4pt,cause] {ISR 20 and\\FSR 27 recoils\\($1.2$, $0.8\GeV$)};
\node[lab,text=black!65,text width=8.4cm,anchor=north west] at (8.8,-3.3) {Three initial-state branchings move her by $\pm20\GeV$ before she has radiated anything: ``\pt after ISR'' is not monotonic and the recoil is taken by the whole event. She radiates twice. Her second emission, with an evolution variable of only $4.4\GeV$, costs $19\GeV$ of laboratory \pt, because $\ptevol$ is a transverse momentum in the dipole frame, not the gluon's \pt in the detector.};
% ---------------- (b) the family cone ---------------------------------------------------------------
\begin{scope}[xshift=4.6cm,yshift=-6.0cm]
\node[head,anchor=west] at (-5.1,2.7) {(b) the family: $\dR<0.4$ around her last copy};
\draw[thick,red!80!black,fill=red!6] (0,0) circle (1.3);
\draw[thick,red!80!black,dashed] (0,0) circle (2.6);
\node[lab,text=red!80!black] at (0,1.55) {$\dR=0.4$};
\node[lab,text=red!80!black] at (0,-2.85) {$\dR=0.8$ (large-radius jets)};
\node[circle,fill=red!80!black,inner sep=0pt,minimum size=4.4mm] at (0,0) {};
\node[lab,anchor=west] at (0.25,0.05) {$b$ \#124, $78.1\GeV$};
\node[circle,fill=black!70,inner sep=0pt,minimum size=3.0mm] at (-0.55,0.55) {};
\node[lab,anchor=east] at (-0.7,0.6) {$g$, $20.7\GeV$};
\node[circle,fill=black!70,inner sep=0pt,minimum size=1.8mm] at (0.35,-0.75) {};
\node[lab,anchor=north] at (0.35,-0.9) {$g$, $1.8\GeV$};
\node[diamond,fill=green!50!black,inner sep=0pt,minimum size=3.2mm] at (-3.3,0.3) {};
\node[diamond,fill=green!50!black,inner sep=0pt,minimum size=3.0mm] at (-3.6,-0.2) {};
\node[lab,anchor=north,text=green!50!black] at (-3.4,-0.5) {$u\bar u$ from her\\first gluon,\\$17.8+14.9\GeV$\\at $\dR\approx1.0$};
\node[lab,text=black!65,text width=8.4cm,anchor=north west] at (4.2,0.3) {Inside the cone: $78.1+20.7+1.8\to\pt=100.45\GeV$, $m=13.5\GeV$, $n=3$. Outside: the $38\GeV$ gluon of her first emission, which left at $\dR=1.0$ and split into $u\bar u$. The energy loss of this stage, $113.22-100.45=12.77\GeV$, is radiation that left the cone.};
\end{scope}
\end{tikzpicture}
}
\caption[The mother's $\pt$ copy by copy, and the family cone around her last copy]{(a) The mother's \pt through the record, copy by copy, with what moved it: the recoils of the initial-state branchings (orange), her own two emissions (red) and the final-state recoils, from the \code{-v} log. (b) The family cone of $\dR<0.4$ around her last copy with its three members, the $0.8$ cone of the large-radius jets of later chapters, and the runaway: the $u\bar u$ pair from her first gluon at $\dR\approx1$.}
\label{fig:mother_gives_birth:recoil_chain}
\end{figure}

\paragraph{The family.} The final partons within $\dR<0.4$ of the $b$ are
the $b$ itself at $78.1\GeV$ and two gluons at $20.7$ and $1.8\GeV$;
together they carry $\pt=100.45\GeV$ with an invariant mass of
$13.5\GeV$. The energy loss of this stage, $113.22-100.45=12.77\GeV$ or
$11\%$, is radiation that left the cone: the mother's first gluon, emitted
at $32\GeV$ evolution scale and $38\GeV$ in the laboratory, left at
$\dR=1.0$ from her and split into a $u\bar u$ pair; its descendants are at
$\eta\approx-2.4$, outside the $0.4$ cone and outside even the $0.8$ cone
of the large-radius jets of \cref{ch:three_families}. The family has acquired a mass of $13.5\GeV$ from $4.8$: the
jet mass, which will be measured at the inn, is born here as the invariant
mass of a quark and its collinear children.

\paragraph{The system.} At Born level the $b$ and the $\bar b$ were exactly
back to back with equal \pt. After the shower their vector sum is
$13.3\GeV$ at $\Delta\phi=2.98$: the pair has been kicked. The 24
initial-state partons carry a net $64.4\GeV$ of transverse momentum, the 14
final-state emissions $52.1\GeV$, and the sum over all 40 partons is zero to
the millielectronvolt [\code{momentum\_sum}]. The initial state radiated
more than the final state did, and in the opposite direction: a
$\bbbar$-pair recoiling against a $60\GeV$ initial-state jet is the
normal appearance of a $2\to2$ event at the LHC, and the ``third jet'' of
the later chapters is very often not a child of either quark. The two
incoming partons have changed identity as well: the $\bar u$ at
$x=2.2\times10^{-3}$ is now a gluon at $x=1.7\times10^{-2}$ and the $u$ at
$0.39$ is a $u$ at $0.47$. The proton remnants that must balance those
fractions are not in the record yet; they arrive in
\cref{ch:uninvited_cousins}.

\paragraph{The twin.} The $\bar b$ lost more: $113.2\to73.2\GeV$, family
$92.1\GeV$. At $\eta=-3.4$ she is deep in the forward calorimeter, and the
gluon she radiated at $34\GeV$ split into a second $\bbbar$ pair,
so that the final state contains two $b$ and two $\bar b$ quarks. The
shower does this in a few percent of events, and at the registry office
(\cref{ch:registry_office}) it is one of the reasons a ``$b$ jet'' needs a
definition rather than a label: which of the two $b$ quarks at
$\eta\approx-3.3$ is the jet's flavour? Keep the question.

\Cref{fig:mother_gives_birth:runaways_plot} is the same ledger as a function of
the cone radius: the family grows from the bare quark at $78\GeV$ to
$100.45\GeV$ at $R=0.4$, where the $20.7\GeV$ gluon is inside, and only
beyond $R\approx1$, when the $u\bar u$ pair of the first gluon is captured,
does it climb back towards the Born value, together with partons that are
not her descendants at all. The radius a jet algorithm chooses is the
radius at which this curve is read.

\begin{figure}[htbp]\centering
\includegraphics[page=10,width=\textwidth]{PS/Standalone/theory_parton_showering.pdf}
\caption[The mother's $\pt$ chain, her descendants, and the family $\pt$ against the cone radius]{Children who run away, from \file{plot_theory_PS.py} and the record \file{event_PS.lhe}. (a) The \pt of the mother's current copy against the branching number: the three initial-state recoils, her own emissions (dotted lines) and the final-state recoils. (b) Every final parton in the plane of $\dR$ from her last copy and \pt, her descendants in red and everything else in grey, with the $0.4$ and $0.8$ cones. (c) The vector-summed \pt inside a cone of radius $R$ around her, for all partons and for her descendants alone; the ledger reads this curve at $R=0.4$.}
\label{fig:mother_gives_birth:runaways_plot}
\end{figure}

\section{Validation from official packages}
\label{sec:mother_gives_birth:validation}

As at the first stop, the last word goes to the official packages. Two are
available here: \PYTHIA itself, driven from a script so that it showers
exactly our hard event and nothing else, and the full CMS generator job of
\cref{ch:ancestral_home}, whose log and event record we can now read with
the shower's status codes in hand. A shower is a random process, so the
comparison with \PYTHIA is made on averages over two thousand showers, and
the comparison with the single \CMSSW event is made on structure, not on
numbers.

\subsection{The same shower with \PYTHIA: the validation script}
\label{sec:mother_gives_birth:pythia}

\file{validate_shower_with_pythia.py} in the same directory writes the hard
event of \cref{ch:ancestral_home} two thousand times into a Les Houches
file, with the scale \code{SCALUP} set to $\pt^{\max}=113.32\GeV$, and
showers the copies with \PYTHIA~8.315 in the environment of
\file{example/ME/Standalone/} (\code{pythia8mc} with the LHAPDF grid file
read through \PYTHIA's own \code{LHAGrid1} reader, since the Python wheel
has no LHAPDF plugin). Multi-parton interactions, hadronisation and QED are
switched off, the CP5 shower couplings are set, and
\code{SpaceShower:rapidityOrder} is switched off to match the standalone
implementation (\code{--rapidity-order} keeps it). The script computes the
same per-event quantities as \code{--repeat}, writes \PYTHIA's first
showered event in the standalone record format to
\file{validate_shower_with_pythia_event1.txt}, and prints the two sets of
averages side by side:
\begin{minted}[linenos=false,frame=single,framesep=1.5mm,fontsize=\small]{bash}
cd example/PS/Standalone
PY=../../ME/Standalone/.venv/bin/python
$PY validate_shower_with_pythia.py -n 2000 --standalone-stats s.stats.json
$PY validate_shower_with_pythia.py -n 2000 --no-isr
$PY validate_shower_with_pythia.py -n 2000 --rapidity-order
\end{minted}
\PYTHIA showers the two thousand copies in about a minute. The first
command produces the first two columns of
\cref{tab:mother_gives_birth:numbers}; the README in the directory has the
full tables. Read them as a physicist, not as a referee. Everything that
carries momentum agrees: the number of partons above $20\GeV$, the \pt and
mass of each $b$ family, the number of initial-state branchings and the
momentum fractions after them, all within one to three percent, which is
the level at which the standalone tables and interpolations differ from
LHAPDF. The disagreement is in the soft tail: the standalone shower makes
$7\%$ more final-state branchings, all between $1$ and $5\GeV$, and $14\%$
more when ISR is off on both sides and only the two $b$ quarks radiate.
Raising the FSR cut-off to $2\GeV$ on both sides leaves a $7\%$ excess in
the last octave, so this is a normalisation of the soft-emission rate near
the cut-off, where the coupling is $0.3$ to $0.5$ and no two implementations
of the infrared region agree exactly. Two further numbers from the README
teach more than the agreement does. With the beam-recoil damping of
\cref{eq:mother_gives_birth:dampen} switched off the standalone shower
makes twice as many final-state branchings; and with \PYTHIA's default
rapidity ordering of the initial state switched back on, \PYTHIA itself
makes $3.4$ instead of $5.1$ initial-state branchings and $16$ instead of
$20$ final-state ones, with the $b$ families unchanged to a percent. The
softest activity of a shower is a matter of options; the hard structure is
not.

\subsection{How CMS does it in practice: the shower inside the GEN step}
\label{sec:mother_gives_birth:cmssw}

In CMS nobody runs the shower separately. The \code{GEN} step of
\cref{sec:ancestral_home:cmssw} hands the hard process to \PYTHIA's
\code{PartonLevel}, which runs the interleaved ISR, FSR and multi-parton
interactions, adds the beam remnants and passes the result to
hadronisation, all inside the one \code{cmsRun} job; the output is the
\code{genParticles} collection, in which every stage is still visible
through the status codes. The book's own analysis package,
\code{JitJetSW}, adds an analyser that draws each stage separately. To
reproduce the figures of this chapter:
\begin{minted}[linenos=false,frame=single,framesep=1.5mm,fontsize=\small]{bash}
cmsrel CMSSW_15_0_5
cd CMSSW_15_0_5/src
cmsenv
git clone git@github.com:ravindkv/JitJetSW.git
cd JitJetSW && scram b
cd Journey/test
cmsRun step1_GEN_cfg.py 2>&1 | tee step1_GEN_PS_cfg.log
cmsRun genPartAnalyzer_cfg.py
\end{minted}
The first \code{cmsRun} is the configuration of \cref{ch:ancestral_home}
with the \code{customiseJitJet} line active, which adds a
\code{ParticleListDrawer} that prints the whole generator record of the
first event; the second reads the output file and writes one figure per
stage, \file{genPartAnalyzer_stage1_hardProcess_run1_event1.pdf} through
stage~5. The log and the first two figures are in
\file{example/PS/CMSSW/}.

\subsubsection{Reading the log}

Three things in \file{step1_GEN_PS_cfg.log} close the loop with this
chapter. At initialisation, the warning
``\code{SimpleTimeShower::init: pTmin too low , raised to 0.611}'' is
\cref{eq:mother_gives_birth:ptmin}, and the settings table lists the four
CP5 shower lines of \cref{sec:mother_gives_birth:cp5}. In the Info Listing
of the first event, after the invariants we used at the previous stop, come
\begin{minted}[linenos=false,frame=single,framesep=1.5mm,fontsize=\small]{text}
 Max pT scale for MPI =  1.133e+02, ISR =  1.133e+02, FSR =  1.133e+02.
 Number of MPI =    32, ISR =    28, FSRproc =    64, FSRreson =     0.
\end{minted}
The first line is \cref{eq:mother_gives_birth:ptmax}, the same number for
all three interleaved processes. The second is not directly comparable with
\cref{tab:mother_gives_birth:numbers}: with multi-parton interactions on,
the 32 additional scatterings radiate too, and the 28 initial-state and 64
final-state branchings count all of them. The single warning
``\code{SimpleSpaceShower::pT2nextQCD: weight above unity}'' is the same
event as the script's counter of weights above one: the PDF-ratio
overestimate of the veto algorithm was exceeded once, and \PYTHIA, like the
script, accepts the branching and tells you.

After the hard-process listing comes the drawer's table, 1846 rows for our
event. It is \PYTHIA's record seen through \CMSSW's \code{GenParticle}
interface: the status column is \PYTHIA's, always positive, and the last
column carries the \emph{status flags} that CMS analyses use instead of
walking the record, \code{isHardProcess}, \code{isFirstCopy},
\code{isLastCopyBeforeFSR}, \code{isLastCopy}. Follow the mother: row~5 is
the hard-process $b$ at $113.22\GeV$ (\code{HardProcess FirstCopy});
rows~9 and~14 are her copies after ISR recoils, status~44, at $112.6$ and
$122.8\GeV$; rows~49 and~182, status~52, are her copies after she recoiled
against two emissions of her dipole partner, at $106.0$ and $91.59\GeV$;
in this shower she never branched herself above the cut-off (no copy has
status~51), and her whole loss is recoil; row~515 is status~62, after the
beam remnants of \cref{ch:uninvited_cousins} have taken their share,
$91.74\GeV$; and row~824, status~71 and \code{LastCopy}, is the $b$ that is
handed to hadronisation at the same momentum. The \code{LastCopyBeforeFSR}
flag is the CMS way of asking ``what was the \pt of the $b$ before it
showered''; it is not set on any copy of a $b$ that never radiated, and
the copy that plays that role here is row~14 at $122.8\GeV$, $9\GeV$
\emph{above} the Born value, entirely from ISR recoil.

\subsubsection{The figure}

\Cref{fig:mother_gives_birth:cmssw} is the analyser's stage-2 figure,
the parton level of the \CMSSW event: 57 partons, 26 from the initial state
and 29 from the final state, the $b$ at $91.59\GeV$ and $\eta=-1.57$, the
$\bar b$ at $118.57\GeV$ and $\eta=-3.73$, their vector sum $27.3\GeV$. Set
it next to \cref{fig:mother_gives_birth:journey}, the standalone shower of
the same hard event, drawn by \file{plot_journey_PS.py} in the same layout
on purpose. The two are different showers of the same event: in one the
mother kept $92\GeV$, in the other $78$; in one the initial state radiated
$27\GeV$ net against the pair, in the other $13$. Neither is more right.
Both sit inside the $94\pm29\GeV$ of two thousand showers, and the
comparison to make between the two codes is
\cref{tab:mother_gives_birth:numbers}, not this pair of pictures. What the
pictures share is the structure: a busy backward hemisphere where the
$\bbbar$ pair and its children are, initial-state partons spread over the
whole rapidity range, and a $b$ quark that has moved a few tenths in
$\eta$ and $\phi$ from where the hard process put it.

\begin{figure}[htbp]\centering
\includegraphics[width=\textwidth]{PS/CMSSW/genPartAnalyzer_stage2_partonShower_run1_event1.pdf}
\caption[The parton level of our event in \CMSSW]{The parton level of our
event in \CMSSW, drawn by the \code{genPartAnalyzer} of \code{JitJetSW} from
the \code{genParticles} of the GEN step (stage~2 of five). Red and blue are
the last copies of the $b$ and $\bar b$ before hadronisation, magenta their
vector sum, orange triangles the initial-state partons and green triangles
the final-state emissions; marker size grows with $\log\pt$ and partons
above $5\GeV$ are labelled. The same event showered by the standalone
script is \cref{fig:mother_gives_birth:journey}.}
\label{fig:mother_gives_birth:cmssw}
\end{figure}

\section{Our jet at this stage: multiplicity, \texorpdfstring{\pt}{pT} and mass of the parton-level family}
\label{sec:mother_gives_birth:our_jet}

Here is the family at the end of the shower. The full record,
\file{example/PS/event_PS.lhe}, has 158 rows; the first 28 are the hard
process, the three hard initial-state branchings and the first emission of
each $b$ quark:

\filelistinglines[\widelisting]{text}{example/PS/event_PS.lhe}{1}{31}{the hard process and the first five branchings of the shower; momenta and masses in GeV}

Read rows~7 and~8 with \cref{sec:mother_gives_birth:step_by_step} in hand:
the $\bar u$ of row~3 has acquired mother~7, a gluon of $x=2.79\times10^{-3}$
(status $-41$, $p_z=+18.98\GeV$), and the $u$ quark of row~8 (status $-43$,
later replaced by its own daughters) is the emitted parton, with the
transverse momentum $\sqrt{14.6^2+42.8^2}=45.2\GeV$ that the whole
downstream system now balances; rows~9 and~10 are the two $b$ quarks after
the recoil, status $-44$. Rows~11 to~21 repeat the pattern twice. Row~21 is
the mother just before her first emission, at $112.5\GeV$; rows~25 and~26
are her daughters, a $b$ of $99.9\GeV$ and a gluon of $38.2\GeV$, status
$-51$ because both branch again, and row~27 is the recoiling gluon of the
dipole. Far down the record her last copy is
\filelistinglines[\widelisting]{text}{example/PS/event_PS.lhe}{126}{131}{the mother's final copy (row 124, status 52) among her neighbours}
at $78.12\GeV$, and the momentum sum at the bottom of the file,
\filelistinglines[\widelisting]{text}{example/PS/event_PS.lhe}{162}{162}{the momentum sum over the final state}
is that of the two current incoming partons: $p_z=-3079.4\GeV$ and
$E=3312.5\GeV$, that is $x_A=0.0171$ and $x_B=0.470$ of $6800\GeV$
each, and a mass of $1220.6\GeV$ where the hard process had $396$. The
initial-state shower has tripled the mass of the system the two protons
handed over; the extra energy is in the 24 initial-state partons.

\begin{table}[htbp]\centering\small
\caption[Kinematics of the two $b$ quarks and their families after the shower]{Kinematics
of the two $b$ quarks before and after the shower, and of their families
(final partons within $\dR<0.4$), from \file{event_PS.lhe} as printed by
\file{standalone_parton_showering.py} and \file{plot_journey_PS.py}. The
$b$ quark is the mother we follow.}
\label{tab:mother_gives_birth:kinematics}
\footnotesize
\begin{tabular}{llrrrrrr}\toprule
 & & $\pt$ [GeV] & $E$ [GeV] & $m$ [GeV] & $\eta$ & $\phi$ & $n$ \\ \midrule
$b$ & hard process (row 6) & $113.22$ & $252.06$ & $4.80$ & $-1.439$ & $-2.151$ & $1$ \\
    & after the shower (row 124) & $78.12$ & $165.08$ & $4.80$ & $-1.379$ & $-2.728$ & $1$ \\
    & family, $\dR<0.4$ & $100.45$ & $224.59$ & $13.52$ & $-1.441$ & $-2.706$ & $3$ \\ \midrule
$\bar b$ & hard process (row 5) & $113.22$ & $2415.66$ & $4.80$ & $-3.753$ & $0.991$ & $1$ \\
    & after the shower (row 95) & $73.18$ & $1106.94$ & $4.80$ & $-3.409$ & $0.577$ & $1$ \\
    & family, $\dR<0.4$ & $92.13$ & $1414.97$ & $7.02$ & $-3.424$ & $0.589$ & $3$ \\ \midrule
$b+\bar b$ & hard process & $0$ & $2667.7$ & $396.1$ & & & \\
    & after the shower & $13.31$ & $1272.0$ & $236.1$ & $y=-2.37$ & $2.45$ & \\ \midrule
ISR partons & after the shower & $64.39$ & & & & $-1.04$ & $24$ \\
FSR partons & after the shower & $52.06$ & & & & $2.02$ & $14$ \\
all final partons & & $0.00$ & $3312.5$ & $1220.6$ & & & $40$ \\
\bottomrule
\end{tabular}
\end{table}

\Cref{tab:mother_gives_birth:kinematics} and \cref{fig:mother_gives_birth:journey}
show what the record means. The mother has kept her direction to a tenth
in $\eta$ and half a radian in $\phi$, lost $35\GeV$ of \pt as a bare
quark and $13\GeV$ as a family, and her family has a mass of $13.5\GeV$
and three members, of which one, the $1.8\GeV$ gluon, would not survive
the detector thresholds of the later chapters. Her twin at $\eta=-3.4$ is
in the forward calorimeter with two extra $b$ quarks for company. Between
the two, in the backward hemisphere at $\eta\approx-2.4$ and
$\phi\approx-2.5$, sits a third cluster of four partons within $\dR<0.25$
of each other: the $u\bar u$ pair from the mother's first gluon, $17.8$ and
$14.9\GeV$, and two initial-state gluons of $10.6$ and $6.7\GeV$. Half of
it is the mother's family by ancestry and half is not, and nothing but the
mother columns of the record, which no detector will read, tells the two
halves apart. The $u$ quark of the first initial-state branching, now at
$39.5\GeV$, is a fourth cluster on its own at $\phi=-1.3$.

\begin{figure}[htbp]\centering
\includegraphics[width=\textwidth]{PS/journey_PS.pdf}
\caption[The showered event in the $(\eta,\phi)$ plane]{The parton level of
our event after the standalone shower, drawn by
\file{example/PS/plot_journey_PS.py} from \file{example/PS/event_PS.lhe} in
the layout of \cref{fig:mother_gives_birth:cmssw}. Red and blue are the
final copies of the $b$ and $\bar b$, magenta their vector sum; orange
triangles descend from an initial-state emission and green triangles from
one of the two hard quarks; marker size grows with $\log\pt$ and partons
above $5\GeV$ are labelled. The header lists the net transverse momentum of
each group; the incoming partons and their momentum fractions before and
after the initial-state shower are noted at the two edges. The second
$\bbbar$ pair from a gluon splitting sits next to the $\bar b$ in the
forward calorimeter.}
\label{fig:mother_gives_birth:journey}
\end{figure}

The ledger gets its second and third rows. The family row is the one the
later stages will be compared with; the bare quark is listed because the
\CMSSW status flags will keep asking for her.

\begin{ledger}{The Mother Gives Birth}
  \ledgerrow{$b$ quark in the hard-process record}{113.22}{4.80}{1}{--}
  \ledgerrow{$b$ family after the shower ($\dR<0.4$)}{100.45}{13.52}{3}{$-12.77$}
  \ledgerrow{\quad of which the $b$ quark itself}{78.12}{4.80}{1}{--}
\end{ledger}

\section{What this shower is, and what it is not}
\label{sec:mother_gives_birth:caveats}

The agreement in \cref{tab:mother_gives_birth:numbers} was obtained without
tuning anything to \PYTHIA's output; it holds because the standalone shower
reproduces the generator's construction. Each choice in that construction
is a caveat to carry to the later stops:
\begin{itemize}
  \item It is a \emph{leading-logarithmic} shower with leading-order
    splitting functions. The first emission of each $b$ quark is not
    corrected to the $q\bar q\to\bbbar g$ matrix element (\PYTHIA does that
    for some processes; the standalone script does not), and emissions
    harder than $\pt^{\max}$ are simply forbidden. \Cref{ch:whose_child} is
    about doing better.
  \item Multi-parton interactions and beam remnants are absent
    (\cref{ch:uninvited_cousins}); the record's incoming partons carry $x$
    fractions with nothing to balance them. Hadronisation is absent
    (\cref{ch:dressing_children}); the ``partons'' below a few GeV are not
    objects any detector will see individually.
  \item \PYTHIA's default rapidity ordering of the initial-state shower is
    not implemented, and the \CMSSW run has it on. It changes the number of
    soft branchings by a third and the $b$ families by a percent.
  \item QED radiation, azimuthal spin correlations and the recoil options
    \PYTHIA offers beyond the default dipole scheme are left out.
  \item The shower's coupling is \PYTHIA's second-order $\alphas$ with
    $\alphas(m_Z)=0.118$; changing the tune changes it, and a shower
    $\alphas$ is not the $\overline{\mathrm{MS}}$ coupling of the world
    average (the CMW rescaling of $\Lambda$ that some showers apply is not
    switched on in CP5).
  \item The family cone of $\dR<0.4$ is a bookkeeping choice; jets have
    not been defined yet.
  \item One shower is one sample. Every number in
    \cref{sec:mother_gives_birth:our_jet} would be different with another
    seed, and the ledger rows are those of \emph{this} event, not of an
    average $b$ quark.
\end{itemize}

\section{Further reading}
\label{sec:mother_gives_birth:further_reading}
The \pt-ordered dipole shower and its interleaving with multi-parton
interactions were introduced in \cite{Sjostrand:2004ef}; the \PYTHIA~8.3
guide~\cite{Bierlich:2022pfr}, chapters on parton showers and on the
generator's settings, is the reference for every option named here and
documents the newer dipole-recoil and \code{Vincia}~\cite{Fischer:2016vfv}
and \code{Dire}~\cite{Hoche:2015sya} showers that the standalone script
does not attempt; the angular-ordered alternative is Herwig's
\cite{Bellm:2019zci}. The splitting functions and the Sudakov form factor
are reviewed in the QCD section of the Particle Data
Group~\cite{ParticleDataGroup:2024cfk} and derived at textbook length in
\cite{Ellis:1996mzs}. The CP5 tune is documented in \cite{CMS:2019csb}, the
automated shower variations of \PYTHIA in \cite{Mrenna:2016sih}. The direct
observation of the dead cone is \cite{ALICE:2021aqk}, its exposure with
jet-substructure techniques \cite{Maltoni:2016ays}. Matching and merging,
the answer to what happens above $\pt^{\max}$, are the subject of the next
chapter and of \cite{Mangano:2006rw,Frederix:2012ps,Lonnblad:2001iq}.

\begin{pitfalls}
\begin{itemize}[nosep]
  \item \textbf{One shower is not a prediction.} $78$, $92$ and $94\GeV$
    for the mother's \pt are one distribution with a width of $29\GeV$.
    Compare averages between codes, never single events, and never ``fix''
    a seed to match another program.
  \item \textbf{$\ptevol$ is not the gluon's \pt.} The evolution
    variable is a transverse momentum in the dipole frame. The mother's
    $4.4\GeV$ emission produced a $20\GeV$ gluon in the laboratory and cost
    her $19\GeV$. A shower history is not a list of jet constituents.
  \item \textbf{ISR moves everything.} Every initial-state branching
    recoils against the whole event, so the $b$ quark's \pt changes before
    she has radiated anything. ``\pt after ISR'' and ``\pt before FSR'' are
    the same copy of the quark (\code{LastCopyBeforeFSR}) and neither is
    the Born value.
  \item \textbf{The starting scale is a choice.} $\pt^{\max}=113.3\GeV$ is
    \PYTHIA's factorisation scale for this process; \code{pTmaxMatch} and
    \code{pTmaxFudge} change it, and an LHE file from \MG carries its own
    \code{SCALUP}. A different starting scale is a different shower.
  \item \textbf{The cut-off is not $0.5\GeV$.} CP5 asks for $0.5$;
    \PYTHIA raises it to $1.6\Lambda_3=0.611\GeV$ and prints a warning that
    every CP5 log contains and nobody reads. The standalone script applies
    the same rule.
  \item \textbf{Counting branchings is meaningless across settings.} The
    number of soft emissions depends on the cut-off, on the recoil damping,
    on rapidity ordering and on which subsystems are included (the \CMSSW
    log counts MPI systems too). Count partons above a scale, or families.
  \item \textbf{Weights above unity.} Both \PYTHIA and the script warn when
    an overestimate of the veto algorithm is exceeded. One or a few per
    event near the cut-off is harmless; many per event means the
    overestimate is wrong and the shower is biased.
  \item \textbf{Two $b$ quarks are not two $b$ jets.} Gluon splitting into
    $\bbbar$ adds heavy quarks to the final state; our event has four. The
    flavour of a jet is a definition (\cref{ch:registry_office}), not a
    count.
  \item \textbf{The dead cone is small at $100\GeV$.} $\theta_0=\mb/E=0.02$
    for our mother. Do not expect a $b$ jet at this \pt to look
    qualitatively different from a light jet at parton level; the
    differences that matter for tagging come from the $B$ hadron's
    lifetime and decay, in \cref{ch:dressing_children}.
  \item \textbf{The record's $x$ values are not the proton's.} After ISR
    the two incoming partons of \file{event_PS.lhe} carry fractions
    $0.017$ and $0.47$ and nothing balances them; the remnants come with the
    next chapter. The momentum sum at the bottom of the file is that of the
    partonic system only.
  \item \textbf{Which Python?} The shower and plotting scripts run with a
    plain \code{python3} that has NumPy, SciPy and matplotlib; the
    validation script needs the virtual environment of
    \file{example/ME/Standalone/} with \code{pythia8mc} and LHAPDF. The
    scripts print what they loaded on their first line.
\end{itemize}
\end{pitfalls}

\begin{exercises}
\item \textbf{Run it.} Run \file{standalone_parton_showering.py} with no
  options and confirm 5 ISR and 33 FSR branchings and a $b$ family of
  $100.45\GeV$. Run it with \code{--seed 2}, \code{3} and \code{4} and
  tabulate the family \pt; then \code{--repeat 200 --quiet --history r.json}
  and compute the mean and spread from \file{r.stats.json}. How many
  showers are needed before the mean is known to $1\GeV$?
\item \textbf{The gluon kernel.} Show that the two forms of $P_{gg}$ in
  \cref{eq:mother_gives_birth:splitting} are equal. Then integrate
  $P_{qg}(z)$ over $0<z<1$ and $P_{qq}(z)$ over $0<z<1-\epsilon$; the first
  is finite and the second is $2C_F\ln(1/\epsilon)+\text{const}$. Find
  the constant. This is why the script overestimates $P_{qq}$ by
  $2C_F/(1-z)$ and $P_{qg}$ by a flat $T_R$.
\item \textbf{The double logarithm.} Reproduce
  \cref{eq:mother_gives_birth:double_log_sudakov} and evaluate
  $\Delta(113.3,\pt)$ for $\pt=50$, $20$, $10$, $5$ and $2\GeV$ with
  $\alphas=0.118$; then repeat with the running coupling of
  \code{AlphaStrong} inside the integral (a ten-line numerical integral).
  Compare with the red curve of scene~2 of the film, and list the reasons
  for the difference given in \cref{sec:mother_gives_birth:sudakov}. Which
  of them can you test with \code{--no-dampen}?
\item \textbf{The veto algorithm is exact.} Split the interval
  $[t,t_{\max}]$ into two halves and consider a true rate $r(t)\le R$
  with an overestimate $R$. Show that the probability of no accepted
  branching in the whole interval, when rejected trials continue from the
  rejected scale, is $\exp[-\int r\,\mathrm{d}\ln t]$ and not
  $\exp[-\int R\,\mathrm{d}\ln t]$ thinned by the average of $r/R$.
  (Sum over the number of rejected trials.)
\item \textbf{Invert the form factor by hand.} Derive
  \cref{eq:mother_gives_birth:veto_solution}. Take $t_{\max}=(\pt^{\max})^2$,
  $n_f=5$, $\Lambda_5=0.2262\GeV$ and $c=2C_F\ln(1/\epsilon)$ with
  $\epsilon=\pt^{\min}/Q$ for $Q=378\GeV$, and compute the median scale of
  the mother's first trial emission. Compare with the $32\GeV$ at which she
  actually radiated in the seed-1 event and with the red histogram of
  scene~2.
\item \textbf{The PDF ratio.} With \code{TablePDF.xf} from the script,
  evaluate $x'f_g(x')/xf_{\bar u}(x)$ at $x=2.17\times10^{-3}$, $z=0.779$
  and $Q=53.9\GeV$, and $x'f_u(x')/xf_u(x)$ at $x=0.390$, $z=0.850$ and
  $Q=51.4\GeV$, confirming $15.7$ and $0.65$. Multiply each by the
  corresponding kernel of \cref{eq:mother_gives_birth:splitting} and explain
  why the $+z$ beam nevertheless wins the first branching in only $53\%$ and
  not $95\%$ of the trials. (Look at the $z$ ranges.)
\item \textbf{Follow the recoil.} From the rows of \file{event_PS.lhe}
  compute the \pt of the mother's copies 6, 10, 15, 21, 25, 46, 60, 89 and
  124 and confirm the chain of \cref{sec:mother_gives_birth:runaways}. For
  the branching from row~21 to rows~25 and~26, boost into the rest frame of
  the dipole (row~21 plus row~18) and verify that the energy fraction of the
  $b$ is $z=0.601$ and that
  \cref{eq:mother_gives_birth:pt_evol_fsr} gives $32.07\GeV$ with
  $m_a=65.67\GeV$.
\item \textbf{The dead cone by numbers.} For the mother's three emissions
  evaluate the two terms of \cref{eq:mother_gives_birth:dead_cone_kernel}
  and confirm the suppressions of $0.3\%$, $2\%$ and $25\%$. Then find, in
  the record, the emission of the $\bar b$ closest to the dead-cone limit.
  Modify the kernel weight in \code{pT2next\_fsr} to drop the mass term,
  rerun with \code{--repeat 500} and \code{--no-isr}, and measure the change
  in the number of FSR branchings and in the family mass.
\item \textbf{Cone size.} Rerun the shower with \code{--radius 0.2},
  \code{0.4}, \code{0.8} and \code{1.0} for seeds 1 to 50 and plot the mean
  family \pt against the radius. At which radius does the family recover the
  Born \pt on average, and why does it overshoot at large radius? (Where are
  the initial-state partons?)
\item \textbf{Off switches.} Run with \code{--no-isr}, with
  \code{--no-fsr}, and with \code{--no-dampen}, each for 200 seeds, and
  tabulate the number of branchings and the family \pt and mass of the $b$.
  Which switch changes the family, and which only the rest of the event?
\item \textbf{Read the \CMSSW record.} In \file{step1_GEN_PS_cfg.log} find
  every copy of the $b$ quark (id~5) and list its status and \pt. Identify
  the rows that correspond to the standalone statuses $44$, $52$, $62$ and
  $71$, and explain what the flags \code{FirstCopy},
  \code{LastCopyBeforeFSR} and \code{LastCopy} mean for it. Which copy would
  you call ``the $b$ quark'' for a $b$-jet energy scale study, and why?
\item \textbf{\PYTHIA's first event.} Compare
  \file{validate_shower_with_pythia_event1.txt} with \file{event_PS.lhe}:
  same format, different shower. Run \file{plot_journey_PS.py} on it and
  count the partons per family. Then run the validation with
  \code{--rapidity-order} and read from its output how the initial-state
  multiplicity changes; the $b$ families should not.
\item \textbf{Film another mother.} Run the animation with
  \code{--seed 5 --scene shower} and, before watching, predict from the
  Sudakov curves of scene~2 whether the first branching is more likely
  initial or final state, and roughly at what \pt. Then run
  \code{--scene sudakov --no-dampen} and describe what happened to the red
  and blue curves.
\item \textbf{Two fates, continued.} The $\bar b$ family at $\eta=-3.4$
  now contains three $b$-flavoured quarks. Which of them will the forward
  calorimeter see as one object, and can any detector tell that the
  cluster is heavy-flavoured? Keep the answer for \cref{ch:registry_office}.
\item \textbf{Read the competition.} In the \code{-v} log, list for rounds
  1 to 5 every proposal, its scale and its fate, and reproduce panel~(b) of
  \cref{fig:mother_gives_birth:interleaving_plot} for those rounds by hand.
  Then explain, with \cref{eq:mother_gives_birth:first}, why the losers of
  a round may be thrown away and re-drawn below the winner's scale without
  biasing the result.
\item \textbf{The ledger against the cone.} Reproduce panel~(c) of
  \cref{fig:mother_gives_birth:runaways_plot} from \file{event_PS.lhe} with
  the reader of \file{plot_journey_PS.py}: the vector-summed \pt of all
  final partons within $R$ of the mother's last copy, for $R$ from $0$ to
  $3$. At which $R$ does the curve first exceed the Born value, and which
  partons are responsible? Repeat with \code{--seed 2} and \code{3} and
  compare the $R=0.4$ and $R=0.8$ values with \cref{tab:mother_gives_birth:numbers}.
\end{exercises}

\subsection*{Open problems in the context of Phase-2 LHC}
\label{sec:mother_gives_birth:open_problems}

The shower is where the simulation stops being a calculation and becomes a
model, and the High-Luminosity LHC will measure jets precisely enough to
notice~\cite{Azzi:2019yne}. Two developments frame everything below. The
first is the question of what a shower is actually accurate for: the
programme started in \cite{Dasgupta:2018nvj} showed that the standard
dipole showers, \PYTHIA's included, fail to be even leading-logarithmic
for some observables, and the PanScales showers of
\cite{Dasgupta:2020fwr,Hamilton:2020rcu,vanBeekveld:2022zhl,vanBeekveld:2023chs}
and the angular-ordered analysis of \cite{Bewick:2019rbu} are the response.
The second is the cost of running any of this for $3000\,\mathrm{fb}^{-1}$
\cite{HSFPhysicsEventGeneratorWG:2020gxw}. Each of the following can be
started with the scripts of this chapter.
\begin{enumerate}
  \item \textbf{Shower uncertainties as weights.} \PYTHIA can vary the
    shower's $\alphas$ and the non-singular terms of its kernels on the fly
    through its \code{UncertaintyBands} settings and return one weight per
    variation~\cite{Mrenna:2016sih}, Herwig does the
    same~\cite{Bellm:2016rhh}, and CMS stores such weights in its samples.
    Reproduce the idea in the standalone script: scale $\alphas$ by a factor
    $\sqrt2$ up and down in the veto weight, keep the same random numbers,
    and measure how much the family \pt of
    \cref{tab:mother_gives_birth:numbers} moves. Is the weight approach
    equivalent to rerunning with \code{--alphas}, and where does it break
    down, in particular when the weight of \cref{eq:mother_gives_birth:veto_weight}
    would have to exceed one?
  \item \textbf{The last octave.} The $7\%$ excess of soft branchings in
    \cref{sec:mother_gives_birth:pythia} is invisible at parton level and
    may not be after hadronisation. CMS has already used jet-substructure
    observables in $\ttbar$ events to constrain the shower
    coupling~\cite{CMS:2018ypj}, and a Phase-2 measurement of constituent
    multiplicity at $\pt\sim100\GeV$ will have a statistical error far
    below $7\%$. Which hadron-level observables are sensitive to the shower
    cut-off region and which are protected by hadronisation? Groomed
    observables such as the soft-drop mass isolate the collinear
    dynamics~\cite{Anderle:2020mxj}; the answer decides what a Run-4 tune
    can constrain.
  \item \textbf{Rapidity ordering and the third jet.} \PYTHIA's rapidity
    ordering of ISR changes the soft initial-state activity by a third
    (\cref{sec:mother_gives_birth:pythia}). Run the validation script with
    and without \code{--rapidity-order} and compare the \pt spectrum and
    rapidity of the hardest initial-state parton, the ``third jet'' of a
    dijet event. This is a region where the logarithmic accuracy of dipole
    showers is known to be deficient~\cite{Dasgupta:2018nvj,vanBeekveld:2022zhl};
    which existing CMS measurement constrains the option, and would a
    $14\TeV$ measurement of the third-jet rate at high rapidity do better?
  \item \textbf{The dead cone at the HL-LHC.} ALICE saw the dead cone in
    charm jets at low \pt~\cite{ALICE:2021aqk}, following the
    reclustering strategy of \cite{Maltoni:2016ays}. Using \code{--repeat}
    with the mass term switched on and off in \code{pT2next\_fsr}, predict
    the ratio of the angular distributions of the first splitting for $b$
    jets at $\pt=100\GeV$ and $500\GeV$. With $3000\,\mathrm{fb}^{-1}$ of
    $b$ jets, could CMS measure the $b$-quark dead cone, and what would the
    result tell us about the $b$ mass in the shower, $4.8\GeV$ here and a
    tune parameter elsewhere?
  \item \textbf{Which \pt is the jet's \pt?} The ledger row of this chapter
    is the family within $\dR<0.4$; CMS status flags offer
    \code{LastCopyBeforeFSR} ($122.8\GeV$ in the \CMSSW event), the last
    copy ($91.7\GeV$) and the hard-process copy ($113.2\GeV$) as ``the $b$
    quark''. Jet energy corrections in \cref{ch:fair_price} are derived
    against generator jets, not quarks, and the same ambiguity is at the
    heart of the interpretation of the top-quark
    mass~\cite{Nason:2017cxd,Hoang:2020iah}. Quantify, with \code{--repeat}
    and different \code{--radius}, how much of the difference between the
    three quark definitions and the family is out-of-cone radiation and how
    much is ISR recoil, as a function of radius
    (\cref{fig:mother_gives_birth:runaways_plot}c); then ask which
    definition a Phase-2 top-mass measurement should use.
  \item \textbf{A shower for $14\TeV$.} The CP5 tune~\cite{CMS:2019csb}
    was fitted at $13\TeV$ with the shower $\alphas$ tied to the
    matrix-element and MPI values. Change the beam energy in
    \file{event_ME.lhe} and the starting scale, and see whether the
    initial-state shower of a $\pthat=113\GeV$ event changes at all between
    $13.6$ and $14\TeV$. If it does not, which observables of the Run-4 tune
    will actually constrain the shower rather than the underlying event?
  \item \textbf{Beyond leading logarithms.} The shower of this chapter, like
    \PYTHIA's, is built from the collinear limit and recovers soft
    wide-angle radiation only through the dipole picture. Compare the
    standalone shower's distribution of the first emission against the
    mother with the NLL expectation for a global observable such as the
    jet mass~\cite{Dasgupta:2020fwr}, using the LL double logarithm of
    \cref{eq:mother_gives_birth:double_log_sudakov} as the first step. Where
    in $(\ln1/\theta,\ln E_g)$ of \cref{fig:mother_gives_birth:propagator_sketch}
    does the dipole shower deviate, and will Run-4 jet-substructure
    measurements at the per-cent level be able to tell a PanScales shower
    from the one CMS uses today?
\end{enumerate}
Later stops will add their own lists; the ones about the detector will be
longer.
